% ============================================================
%  PROTEA — Doctoral Thesis (PhD)
%  Main document entry point
% ============================================================
\documentclass[12pt, a4paper, twoside, openright]{report}

% ── Encoding & language ─────────────────────────────────────
\usepackage[utf8]{inputenc}
\usepackage[T1]{fontenc}
\usepackage[english]{babel}

% ── Typography ──────────────────────────────────────────────
\usepackage{lmodern}
\usepackage{microtype}
\usepackage{setspace}
\onehalfspacing
\setlength{\emergencystretch}{3em}

% ── Page layout ─────────────────────────────────────────────
\usepackage[
  a4paper,
  top=3cm, bottom=3cm,
  inner=3.5cm, outer=2.5cm,
  headheight=28pt
]{geometry}

% ── Headers & footers ───────────────────────────────────────
\usepackage{fancyhdr}
\pagestyle{fancy}
\fancyhf{}
\fancyhead[LE,RO]{\thepage}
\fancyhead[RE]{\textit{\leftmark}}
\fancyhead[LO]{\textit{\rightmark}}
\renewcommand{\headrulewidth}{0.4pt}

% ── Colours ─────────────────────────────────────────────────
\usepackage[dvipsnames]{xcolor}
\definecolor{proteagreen}{RGB}{38, 120, 78}
\definecolor{linkblue}{RGB}{0, 70, 140}

% ── Math ────────────────────────────────────────────────────
\usepackage{amsmath, amssymb, amsthm}

% ── Figures & tables ────────────────────────────────────────
\usepackage{graphicx}
\graphicspath{{figures/}}
\usepackage{booktabs}
\usepackage{longtable}
\usepackage{array}
\usepackage{multirow}
\usepackage{caption}
\usepackage{subcaption}
\usepackage{float}

% ── TikZ for native vector diagrams ─────────────────────────
\usepackage{tikz}
\usetikzlibrary{positioning, arrows.meta, fit, shapes.geometric, calc, backgrounds, shadows.blur}
\usepackage{pgfplots}
\pgfplotsset{compat=1.18}
\usepgfplotslibrary{groupplots, statistics}
\pgfplotsset{
  protea plot/.style={
    width=6.2cm, height=4.6cm,
    tick label style={font=\scriptsize},
    label style={font=\scriptsize},
    title style={font=\small\bfseries},
    legend style={font=\scriptsize, draw=black!30, fill=white, fill opacity=0.85, text opacity=1, inner sep=2pt, row sep=-1pt},
    major grid style={dashed, draw=black!15},
    minor grid style={dotted, draw=black!10},
    axis line style={draw=black!50},
    every axis plot/.append style={line width=1pt, mark size=1.6pt},
    cycle list={
      {proteagreen, mark=*, mark options={fill=proteagreen}},
      {linkblue, mark=square*, mark options={fill=linkblue}},
      {BurntOrange, mark=triangle*, mark options={fill=BurntOrange}},
      {RedViolet, mark=diamond*, mark options={fill=RedViolet}},
      {black!60, mark=pentagon*, mark options={fill=black!60}},
    },
  },
}
\tikzset{
  proteablock/.style={
    draw=proteagreen!70!black, very thick, fill=proteagreen!10,
    rounded corners=4pt, align=center, font=\small,
    minimum height=10mm, minimum width=28mm, blur shadow={shadow blur steps=4}
  },
  proteastore/.style={
    draw=linkblue!70!black, very thick, fill=linkblue!10,
    cylinder, shape border rotate=90, aspect=0.18, align=center,
    font=\small, minimum height=14mm, minimum width=24mm
  },
  proteaqueue/.style={
    draw=BurntOrange!70!black, very thick, fill=BurntOrange!10,
    rounded corners=2pt, align=center, font=\footnotesize\ttfamily,
    minimum height=7mm, minimum width=42mm
  },
  proteastate/.style={
    draw=black!70, thick, fill=black!5,
    rounded corners=3pt, align=center, font=\footnotesize\ttfamily,
    minimum height=7mm, minimum width=18mm
  },
  proteaarrow/.style={-{Stealth[length=2.4mm]}, thick, draw=black!70},
  proteadashed/.style={proteaarrow, dashed},
}

% ── Code listings ───────────────────────────────────────────
\usepackage{listings}
\usepackage{minted}          % requires --shell-escape
\lstset{
  basicstyle=\ttfamily\small,
  breaklines=true,
  frame=single,
  numbers=left,
  numberstyle=\tiny\color{gray},
  keywordstyle=\color{proteagreen}\bfseries,
  commentstyle=\color{gray}\itshape,
  stringstyle=\color{BrickRed},
}

% ── Cross-references & hyperlinks ───────────────────────────
\usepackage[
  colorlinks=true,
  linkcolor=linkblue,
  citecolor=proteagreen,
  urlcolor=linkblue,
  pdftitle={PROTEA: A Framework for Scalable Protein Functional Annotation},
  pdfauthor={Francisco Miguel P\'{e}rez Canales},
  pdfkeywords={protein function prediction, gene ontology, embeddings, bioinformatics},
]{hyperref}
\usepackage{cleveref}

% ── Bibliography ────────────────────────────────────────────
\usepackage[
  backend=biber,
  style=numeric,
  sorting=nyt,
  maxbibnames=6,
]{biblatex}
\addbibresource{bibliography/references.bib}

% ── Glossary / acronyms ─────────────────────────────────────
\usepackage[acronym, toc]{glossaries}
\makeglossaries
% ── Acronyms used throughout the thesis ─────────────────────

% ── Biological databases and resources ──────────────────────
\newacronym{go}{GO}{Gene Ontology}
\newacronym{goa}{GOA}{Gene Ontology Annotation}
\newacronym{uniprot}{UniProt}{Universal Protein Resource}
\newacronym{swissprot}{Swiss-Prot}{Swiss Protein database (reviewed section of UniProt)}
\newacronym{trembl}{TrEMBL}{Translated EMBL Nucleotide Sequence Data Library}
\newacronym{pdb}{PDB}{Protein Data Bank}
\newacronym{ncbi}{NCBI}{National Center for Biotechnology Information}
\newacronym{obo}{OBO}{Open Biological and Biomedical Ontologies}
\newacronym{gaf}{GAF}{GO Annotation File}
\newacronym{bfd}{BFD}{Big Fantastic Database}

% ── GO aspects and CAFA evaluation notation ─────────────────
\newacronym{mf}{MF}{Molecular Function}
\newacronym{bp}{BP}{Biological Process}
\newacronym{cc}{CC}{Cellular Component}
\newacronym{mfo}{MFO}{Molecular Function Ontology}
\newacronym{bpo}{BPO}{Biological Process Ontology}
\newacronym{cco}{CCO}{Cellular Component Ontology}

% ── GO evidence codes ────────────────────────────────────────
\newacronym{iea}{IEA}{Inferred from Electronic Annotation}
\newacronym{exp}{EXP}{Experimental evidence (GO evidence code family covering
  EXP, IDA, IMP, IGI, IEP, and related codes)}
\newacronym{ic}{IC}{Information Content}

% ── Sequence analysis tools ──────────────────────────────────
\newacronym{blast}{BLAST}{Basic Local Alignment Search Tool}
\newacronym{hmm}{HMM}{Hidden Markov Model}
\newacronym{nw}{NW}{Needleman-Wunsch}
\newacronym{sw}{SW}{Smith-Waterman}
\newacronym{lca}{LCA}{Lowest Common Ancestor}
\newacronym{fasta}{FASTA}{Fast Alignment (sequence format)}
\newacronym{diamond}{DIAMOND}{Double Index Alignment of Next-generation Sequencing Data}

% ── Machine learning and embeddings ─────────────────────────
\newacronym{llm}{LLM}{Large Language Model}
\newacronym{plm}{pLM}{Protein Language Model}
\newacronym{esm}{ESM}{Evolutionary Scale Modelling}
\newacronym{knn}{kNN}{k-Nearest Neighbours}
\newacronym{ann}{ANN}{Approximate Nearest Neighbour}
\newacronym{pca}{PCA}{Principal Component Analysis}
\newacronym{mil}{MIL}{Multiple-Instance Learning}
\newacronym{faiss}{FAISS}{Facebook AI Similarity Search}
\newacronym{gpu}{GPU}{Graphics Processing Unit}

% ── Software engineering and infrastructure ──────────────────
\newacronym{api}{API}{Application Programming Interface}
\newacronym{rest}{REST}{Representational State Transfer}
\newacronym{orm}{ORM}{Object-Relational Mapping}
\newacronym{amqp}{AMQP}{Advanced Message Queuing Protocol}
\newacronym{sql}{SQL}{Structured Query Language}
\newacronym{uuid}{UUID}{Universally Unique Identifier}
\newacronym{uri}{URI}{Uniform Resource Identifier}
\newacronym{abc}{ABC}{Abstract Base Class}
\newacronym{adr}{ADR}{Architecture Decision Record}
\newacronym{hpc}{HPC}{High-Performance Computing}
\newacronym{sdk}{SDK}{Software Development Kit}

% ── Ontology structure ───────────────────────────────────────
\newacronym{dag}{DAG}{Directed Acyclic Graph}

% ── Benchmarks and evaluation ────────────────────────────────
\newacronym{cafa}{CAFA}{Critical Assessment of Functional Annotation}
\newacronym{lafa}{LAFA}{Live Assessment of Functional Annotation}

% ── Observability and operations ─────────────────────────────
\newacronym{otel}{OTel}{OpenTelemetry}
\newacronym{slurm}{SLURM}{Simple Linux Utility for Resource Management}


% ── Front/main/backmatter for report class ──────────────────
\newcommand{\frontmatter}{\pagenumbering{roman}}
\newcommand{\mainmatter}{\cleardoublepage\pagenumbering{arabic}}
\newcommand{\backmatter}{\cleardoublepage}

% ── Utilities ───────────────────────────────────────────────
\usepackage{todo}            % \todo{...} inline notes
\usepackage{csquotes}
\usepackage{enumitem}

% ── Theorem-like environments ───────────────────────────────
\theoremstyle{definition}
\newtheorem{definition}{Definition}[chapter]
\newtheorem{example}{Example}[chapter]

% ============================================================
\begin{document}

% ── Front matter ────────────────────────────────────────────
\frontmatter
\begin{titlepage}
  \centering
  \vspace*{1.2cm}

  % To add the official Universidad de Sevilla emblem, drop the logo file into
  % figures/ (e.g. figures/us-logo.pdf) and uncomment the line below.
  % \includegraphics[height=2.4cm]{us-logo}\\[0.8cm]

  {\Large\scshape Universidad de Sevilla\par}
  \vspace{0.4cm}
  {\large Escuela T\'{e}cnica Superior de Ingenier\'{i}a Inform\'{a}tica\par}
  \vspace{0.15cm}
  {\large Programa de Doctorado en Ingenier\'{i}a Inform\'{a}tica\par}

  \vspace{2.6cm}

  {\color{proteagreen}\rule{\linewidth}{0.4pt}}\\[0.6cm]
  {\huge\bfseries PROTEA\par}
  \vspace{0.35cm}
  {\Large\bfseries PROtein funcTional Embedding-based Annotation\par}
  \vspace{0.6cm}
  {\color{proteagreen}\rule{\linewidth}{0.4pt}}\\[1.3cm]

  {\Large\itshape Doctoral Thesis\par}
  \vspace{0.35cm}
  {\large\itshape submitted in partial fulfilment\\
  of the requirements for the degree of\\
  Doctor of Philosophy\par}

  \vspace{2.0cm}

  \begin{minipage}[t]{0.46\textwidth}
    \begin{flushleft}
      \large
      {\scshape\color{proteagreen} Author}\\[0.15cm]
      Francisco Miguel P\'{e}rez Canales
    \end{flushleft}
  \end{minipage}
  \hfill
  \begin{minipage}[t]{0.46\textwidth}
    \begin{flushright}
      \large
      {\scshape\color{proteagreen} Co-supervisors}\\[0.15cm]
      David Orellana-Mart\'{i}n\\
      Ana M.\ Rojas
    \end{flushright}
  \end{minipage}

  \vfill

  {\large Academic Year 2025/2026\par}
\end{titlepage}

\cleardoublepage
\chapter*{Abstract}
\addcontentsline{toc}{chapter}{Abstract}

Predicting the biological function of proteins at scale remains one of the central
challenges in computational biology. Experimental characterisation cannot keep pace
with the rate at which new sequences are deposited in public databases, leaving the
vast majority of known proteins with no experimentally validated functional annotation.
The gap is widening: large-scale protein language models (PLMs) and structural
predictors generate dense representations of every UniProt sequence in days, yet the
infrastructure required to turn those representations into auditable, reproducible Gene
Ontology predictions is fragmented across notebooks, prototype scripts and one-off
benchmarks.

This thesis presents \textbf{PROTEA} (\textit{PROtein funcTional annotation framEwork
and plAtform}), an open-source distributed platform that automates the full pipeline
from raw protein sequence to structured Gene Ontology predictions, and provides an
extensible foundation on top of which third parties can plug in new PLMs, annotation
sources, and ranking methods. PROTEA combines protein language model embeddings (eight
backends evaluated, including ESM-2 in three sizes, ESM-C, ProstT5, ProtT5, Ankh and
ESM-3 component encoder), approximate $k$-nearest-neighbour search over 527,000 reviewed
UniProt sequences, and evidence transfer from versioned annotation sets to produce
ranked, interpretable predictions. A versioned relational data model links every
prediction to the exact ontology release, annotation set, embedding configuration and
code commit that produced it, enabling exact replay of any past run.

To improve prediction quality beyond embedding cosine similarity, PROTEA incorporates a
plugin-based feature engineering layer that computes Needleman--Wunsch and
Smith--Waterman alignment statistics, NCBI-taxonomy-based distance metrics, ancestry
embeddings of the GO directed acyclic graph, and per-PLM PCA projections, yielding a
schema of approximately 52 features per candidate annotation. A LightGBM re-ranker is
trained on these features using a temporal holdout protocol spanning twelve consecutive
GOA releases, with Information Accretion weighting, per-category models (No-Knowledge,
Limited-Knowledge, Prior-Knowledge), and selective application that falls back to the
embedding-only baseline when re-ranking would degrade performance.

Beyond the prediction method, the thesis describes the engineering pillars that make
the platform deployable across heterogeneous environments: a seven-repository plugin
architecture discovered via Python entry points, externalised configuration via
\texttt{pydantic-settings}, OpenTelemetry-based distributed tracing, multi-target
container packaging covering development, cloud, HPC (Apptainer) and airgapped
environments, and an operational narrative model that makes every Job and
\texttt{ExperimentRun} self-describing.

Evaluation follows the CAFA temporal protocol on held-out GOA splits and reports the
performance of the canonical pipeline across the eight PLM backends and three category
slices, contextualises the impact of each feature family through ablation studies, and
documents the methodological pivots that shaped the canonical configuration. Three
pre-built containers are submitted to the LAFA benchmark to demonstrate adoption-ready
deployment.

\bigskip
\noindent\textbf{Keywords:} protein function prediction, gene ontology, protein
language models, sequence embeddings, nearest-neighbour search, learning-to-rank,
LightGBM, distributed systems, reproducibility, container deployment, bioinformatics
infrastructure.

\cleardoublepage
\chapter*{Acknowledgements}
\addcontentsline{toc}{chapter}{Acknowledgements}

I am grateful to my co-supervisors, David Orellana-Mart\'{i}n and Ana M.\ Rojas,
for their guidance throughout this doctoral research. Their complementary perspectives
on computational methods and biological application have shaped the direction of this
work, and their willingness to engage with technical decisions at every stage has been
invaluable.

The second-place team result in CAFA~6 reflects a collective effort. I contributed the
technical infrastructure and prediction pipeline described in this thesis, and I thank
the broader CAFA~6 team for the collaboration that made that placement possible.
PROTEA, as presented here, is the productisation of that technical contribution and
stands as an independent open-source platform.

% TODO: funding line

My family and closest friends have provided steady support throughout this process,
and I am grateful for it.

\cleardoublepage

\tableofcontents
\cleardoublepage
\listoffigures
\listoftables
\cleardoublepage

\printglossary[type=\acronymtype, title=Acronyms]
\cleardoublepage

% ── Main matter ─────────────────────────────────────────────
\mainmatter
\chapter{Introduction}
\label{chap:introduction}

\section{Motivation}
\label{sec:motivation}

The pace at which biological sequence data accumulates has long outstripped our
capacity to characterise proteins experimentally. As of early 2026, \gls{uniprot}
contains more than 250 million sequences, yet fewer than 0.02\,\% carry
experimentally validated functional annotations~\cite{uniprot2023}. The gap is not
merely quantitative: the proteins that remain unannotated are disproportionately
found in under-studied organisms, precisely where novel biology, and potential
therapeutic targets, are most likely to reside.

Computational function prediction attempts to bridge this gap by transferring
knowledge from characterised proteins to unknown ones. Classical approaches based on
sequence similarity (BLAST, profile \glspl{hmm}) remain widely used, but their
accuracy degrades for distantly related sequences. Over the last five years, however,
\glspl{plm} trained on hundreds of millions of sequences have produced dense vector
representations that capture structural and functional signal well beyond what
classical alignment can detect, and a generation of structure-conditioned models
(ESM-3, ProstT5) has begun to integrate sequence and structure into a single
representation. Each new PLM unlocks a new operating point in the accuracy/throughput
trade-off, but also introduces a new dependency, a new tokeniser, a new GPU memory
footprint, and a new failure mode. Comparing PLMs at scale, on a fair benchmark, with
reproducible numbers, is itself a research engineering problem.

PROTEA was designed to address both faces of the problem at once: as a method, it
operationalises embedding-based annotation transfer with selective learning-to-rank
re-finement; as a platform, it provides the extensible foundation that lets new
backends, annotation sources, and ranking methods be plugged in without ad-hoc
glue. The platform is deliberately built around the assumption that the canonical
pipeline of today will be obsolete tomorrow. What must remain stable is the contract
between layers, the trace of every run, and the pathway from a FASTA file to a
ranked, interpretable Gene Ontology prediction.

\section{Research Questions}
\label{sec:research-questions}

This thesis addresses the following questions:

\begin{enumerate}[label=\textbf{RQ\arabic*.}]
  \item Can embedding-based \gls{knn} annotation transfer match or exceed
        alignment-based methods for \gls{go} term prediction, particularly at low
        sequence identity, and how does the answer change as new PLMs replace
        older ones?
  \item Do pairwise alignment statistics, taxonomic distance, ontology-graph
        ancestry features, and per-PLM PCA projections provide complementary signal
        that improves the ranking of candidate annotations once the embedding
        baseline is already strong?
  \item How can a reproducible, scalable annotation pipeline be designed so that
        researchers can re-run analyses as ontologies and databases are updated, and
        so that the pipeline itself can be deployed in development, cloud, on-premise
        HPC, and airgapped environments without code branching?
  \item How can the journey of a research platform, with its dead ends and pivots,
        be captured operationally so that the canonical pipeline is justified by an
        auditable trace rather than by post-hoc rationalisation?
\end{enumerate}

\section{Contributions}
\label{sec:contributions}

The main contributions of this work are:

\begin{itemize}
  \item \textbf{PROTEA}: an open-source distributed platform implementing the full
        annotation pipeline, from FASTA upload to structured \gls{go} predictions
        with optional feature engineering and selective learning-to-rank re-ranking.
  \item A \textbf{plugin architecture} based on Python \texttt{entry\_points},
        organised across seven code repositories, that lets new annotation sources,
        runners and PLM backends be added without modifying the core platform.
  \item A \textbf{versioned data model} for ontology snapshots, annotation sets,
        embedding configurations, datasets, re-ranker models and experiment runs,
        with provenance complete enough to exactly replay any past prediction.
  \item A \textbf{benchmark of eight PLM backends} (ESM-2 in three sizes, ESM-C,
        ProstT5, ProtT5, Ankh, ESM-3 component encoder) on a common evaluation
        surface, with per-aspect, per-category and per-backend numbers.
  \item A \textbf{feature engineering layer} that augments embedding similarity with
        Needleman--Wunsch / Smith--Waterman alignment scores,
        NCBI-taxonomy-based distance metrics, GO-graph ancestry embeddings, and
        per-PLM PCA projections, organised as a registry of approximately 52
        features.
  \item A \textbf{selective learning-to-rank re-ranker} trained per category and per
        aspect, that falls back to the embedding-only baseline whenever re-ranking
        would degrade performance for the current cell.
  \item A \textbf{multi-target deployment story} (development \texttt{docker
        compose}, cloud Helm/Compose, HPC Apptainer, airgapped bundle) covering the
        same canonical pipeline without code branching.
  \item An \textbf{operational narrative model} (\texttt{Job.findings},
        \texttt{ExperimentRun}, comments) that captures the why of every run and
        provides the audit trace from which the evaluation chapter is distilled.
  \item An \textbf{empirical evaluation} comparing PROTEA predictions against
        baseline methods on held-out \gls{goa} splits following the CAFA temporal
        protocol, complemented by three submissions to the LAFA public benchmark.
\end{itemize}

\section{Thesis Structure}
\label{sec:structure}

The remainder of this document is organised as follows.
\Cref{chap:background} introduces the biological concepts underpinning the work,
including the modern landscape of protein language models and their structural
extensions.
\Cref{chap:related} surveys existing computational approaches to protein function
prediction up to 2025, with attention to the most recent embedding-based and
graph-based methods.
\Cref{chap:design} presents the architecture of PROTEA, its plugin-based
extensibility model, and the operational narrative layer.
\Cref{chap:implementation} describes key implementation decisions across the seven
repositories, including the configuration system, the artifact-store boundary, the
queue topology, and the deployment containers.
\Cref{chap:evaluation} reports the experimental results on the eight-PLM benchmark,
the impact of feature families through ablation, and the journey that shaped the
canonical configuration.
\Cref{chap:conclusion} summarises the findings, discusses limitations, and outlines
future directions including the LAFA submissions and the post-defense agenda.

\chapter{Biological Background}
\label{chap:background}

This chapter provides the biological foundation required to understand the
problem that PROTEA addresses. Readers with a background in molecular biology
may skip directly to \cref{chap:related}.

% ────────────────────────────────────────────────────────────
\section{Proteins: Structure and Function}
\label{sec:proteins}

Proteins are macromolecules composed of chains of amino acids linked by peptide
bonds. The linear sequence of amino acids — encoded by the genome — is referred
to as the \textit{primary structure}. This sequence folds, driven by
thermodynamic forces, into a specific three-dimensional conformation that
determines the protein's biological activity.

Protein functions are extraordinarily diverse: enzymes catalyse chemical
reactions, structural proteins provide mechanical support, signalling proteins
relay information between cells, and transport proteins shuttle molecules across
membranes. Understanding \textit{what} a protein does — its molecular function,
the biological process it participates in, and the cellular compartment where it
operates — is fundamental to biology and medicine.

\subsection{Amino Acids and the Sequence Space}
\label{subsec:amino-acids}

Twenty standard amino acids, differing in their side-chain chemistry, constitute
the alphabet of protein sequences. A protein of $n$ residues can in principle
take $20^n$ distinct sequences, an astronomically large space. Yet the sequences
that fold into stable, functional structures represent a tiny, structured subset
of this space, one in which similar sequences tend to adopt similar folds and
perform similar functions.

\subsection{Sequence Identity and Homology}
\label{subsec:homology}

Two proteins are said to be \textit{homologous} if they share a common
evolutionary ancestor. Homology is typically inferred from sequence similarity:
proteins with high pairwise sequence identity are almost always homologous and
usually perform the same function. As sequence identity drops below roughly
30\,\%, however, the relationship between similarity and function becomes
ambiguous — proteins may be structurally similar yet functionally diverged, or
functionally convergent with unrelated folds.

% ────────────────────────────────────────────────────────────
\section{Protein Databases}
\label{sec:databases}

\subsection{UniProt}
\label{subsec:uniprot}

\gls{uniprot} is the central repository for protein sequence and functional
information~\cite{uniprot2023}. It is divided into two sections:

\begin{description}
  \item[\gls{swissprot}] A manually curated, high-quality database. Each entry
    has been reviewed by expert annotators and contains experimental or
    computationally inferred functional annotations, with explicit evidence codes.
  \item[\gls{trembl}] An automatically annotated, computationally analysed
    database derived from genome translations. It is several orders of magnitude
    larger than Swiss-Prot but of lower average annotation quality.
\end{description}

Each protein in UniProt is identified by a stable accession number. Canonical
isoforms are represented by a bare accession (e.g.\ \texttt{P12345}); alternative
splicing variants are denoted \texttt{P12345-2}, \texttt{P12345-3}, and so on.

\subsection{Protein Data Bank}
\label{subsec:pdb}

The \gls{pdb} archives three-dimensional structures of biological macromolecules
determined experimentally by X-ray crystallography, cryo-electron microscopy, or
NMR~\cite{berman2000pdb}. Although PROTEA operates primarily on sequences, PDB
structures are an important source of ground-truth functional information and are
used indirectly through UniProt cross-references.

% ────────────────────────────────────────────────────────────
\section{The Gene Ontology}
\label{sec:gene-ontology}

\subsection{Overview}
\label{subsec:go-overview}

The \gls{go} is a standardised, species-independent controlled vocabulary for
describing gene product attributes~\cite{ashburner2000go,goconsortium2023}.
It is structured as three independent \glspl{dag}, one for each of the three
top-level \textit{aspects}:

\begin{description}
  \item[\gls{mf} (GO:0003674)] The biochemical activity of the gene product,
    independent of where or when it occurs (e.g.\ \textit{ATP binding},
    \textit{serine-type endopeptidase activity}).
  \item[\gls{bp} (GO:0008150)] The larger biological programme to which the
    activity contributes (e.g.\ \textit{apoptotic process},
    \textit{DNA repair}).
  \item[\gls{cc} (GO:0005575)] The place in the cell where the gene product
    is active (e.g.\ \textit{nucleus}, \textit{plasma membrane}).
\end{description}

Terms within each aspect are connected by \textit{is-a} and \textit{part-of}
relationships, forming a hierarchy from broad root terms to highly specific
leaves. A protein annotated with a term is implicitly annotated with all of its
ancestors up to the root — the \textit{true path rule}.

\subsection{Evidence Codes}
\label{subsec:evidence-codes}

Every GO annotation carries an evidence code that indicates the type of support
for the association. Codes range from experimental (e.g.\ \texttt{EXP},
\texttt{IDA}, \texttt{IMP}) to computationally inferred (\texttt{ISS}, \texttt{IEA})
and author-stated (\texttt{TAS}, \texttt{NAS}). Experimental evidence codes are
the gold standard; \texttt{IEA} (Inferred from Electronic Annotation) is the most
common but least reliable.

\subsection{Ontology Versioning}
\label{subsec:go-versioning}

The GO is actively maintained and released regularly in \gls{obo} format.
New terms are added, obsolete terms are deprecated, and relationships are
refined with each release. Reproducible research therefore requires that
analyses record which ontology version was used. PROTEA stores each imported
release as an \texttt{OntologySnapshot} and associates all annotation sets
and predictions with a specific snapshot.

% ────────────────────────────────────────────────────────────
\section{Sequence Alignment}
\label{sec:alignment}

Pairwise sequence alignment is the classical approach to measuring sequence
similarity. Two algorithms are central to this thesis:

\begin{description}
  \item[Needleman--Wunsch (\gls{nw})] A global alignment algorithm that aligns
    two sequences end-to-end, maximising a cumulative substitution score minus
    gap penalties~\cite{needleman1970general}. It is suitable for comparing
    proteins of similar length.
  \item[Smith--Waterman (\gls{sw})] A local alignment algorithm that finds the
    highest-scoring contiguous aligned region, ignoring the tails of each
    sequence~\cite{smith1981identification}. It is more appropriate when one
    sequence is a domain or fragment of the other.
\end{description}

Both algorithms use a \textit{substitution matrix} — typically BLOSUM62 for
protein sequences~\cite{henikoff1992amino} — to score residue substitutions
based on their observed frequency in alignments of related proteins.

% ────────────────────────────────────────────────────────────
\section{Protein Language Models and Embeddings}
\label{sec:plm}

\Glspl{plm} are neural networks trained on large corpora of protein sequences
using self-supervised objectives borrowed from natural language processing.
Models such as ESM-2~\cite{lin2023esm2} and ProtTrans~\cite{elnaggar2022prottrans}
learn to predict masked residues, in doing so developing internal representations
that capture evolutionary, structural, and functional information.

A protein sequence can be passed through such a model to obtain a dense vector
(embedding) that encodes its properties in a continuous latent space. Proteins
with similar functions tend to cluster together in this space even when their
sequences have diverged below the homology detection threshold of alignment-based
tools. This property makes embeddings attractive for annotation transfer.

% ────────────────────────────────────────────────────────────
\section{Taxonomic Classification}
\label{sec:taxonomy}

Biological organisms are classified into a hierarchy — the NCBI taxonomy — that
reflects their evolutionary relationships~\cite{schoch2020ncbi}. Each node is
assigned a numeric \textit{taxon ID}. The distance between two organisms can be
defined operationally as the number of edges between them through their
\gls{lca} in the taxonomy tree.

Taxonomic distance is a useful co-variate for function prediction: a protein from
a closely related organism is more likely to perform the same function than one
from a distant lineage, even when embedding similarity is comparable.

\chapter{Related Work}
\label{chap:related}

Automated protein function prediction has been an active research area for more
than two decades. Methods can be grouped into four broad families: sequence
alignment transfer, profile and domain matching, supervised machine learning, and
embedding-based approaches. This chapter surveys each family, identifies their
limitations, and contextualises PROTEA within the landscape.

% ────────────────────────────────────────────────────────────
\section{Alignment-Based Annotation Transfer}
\label{sec:related-alignment}

The foundational insight of alignment-based methods is that sequence similarity
implies functional similarity. If a query protein shares high sequence identity
with a characterised protein, one can transfer the known annotations to the query.

\textbf{BLAST}~\cite{altschul1990blast} remains the most widely deployed tool for
this purpose. It uses a heuristic seed-and-extend strategy to find local
alignments in sublinear time. \texttt{BLASTP} (protein--protein) identifies
homologues in Swiss-Prot and transfers their GO terms to the query. This simple
approach is accurate when identity exceeds roughly 50\,\% and degrades
progressively as identity falls. Below the ``twilight zone'' of $\approx$30\,\%
identity, alignment scores lose statistical power and functional inference becomes
unreliable~\cite{rost1999twilight}.

\textbf{DIAMOND}~\cite{buchfels2021diamond} offers BLAST-comparable sensitivity
at orders-of-magnitude higher throughput, making it practical for genome-scale
annotation. Its speed advantage comes from reduced-alphabet seeds and better
SIMD utilisation; the fundamental twilight-zone limitation remains.

\textbf{Annotation propagation rules} add another layer of nuance: not all GO
terms are equally transferable. Terms annotated with \texttt{IEA} codes can be
propagated freely, while experimental annotations require explicit curation
policies. The CAFA community has established benchmark protocols that penalise
propagation errors~\cite{zhou2019cafa3}.

% ────────────────────────────────────────────────────────────
\section{Profile and Domain Methods}
\label{sec:related-hmm}

To extend sensitivity below the twilight zone, profile methods compare a query
against position-specific models built from multiple sequence alignments of
protein families.

\textbf{HMMER}~\cite{eddy2011hmmer} implements profile \glspl{hmm} for this
purpose. A family is represented as a probabilistic model over residue
frequencies at each position, capturing conserved and variable sites more
faithfully than a single representative sequence. HMMER achieves greater
sensitivity than BLAST at the cost of requiring pre-built profiles.

\textbf{InterPro}~\cite{blum2021interpro} integrates over a dozen family and
domain databases — Pfam, PRINTS, ProSite, TIGRFAM, and others — into a unified
annotation resource. A protein is assigned InterPro entries whenever it matches
any member database; GO terms are then inferred via curated mappings. InterPro
annotation is one of the most reliable computational methods available, but it
is limited to protein families and domains that have been explicitly curated and
modelled.

\textbf{Cascade homology search}~\cite{park1998sequence} iteratively extends
sensitivity by feeding BLAST hits into further BLAST searches (PSI-BLAST) or
HMM construction, at the cost of introducing false positives in remote homology
regimes.

% ────────────────────────────────────────────────────────────
\section{Supervised Machine Learning Approaches}
\label{sec:related-ml}

Supervised methods learn classifiers — one per GO term of interest — from
labelled protein--function pairs. Early approaches used \glspl{hmm} or support
vector machines on curated feature sets (amino acid composition, physicochemical
properties, secondary structure predictions). More recent work uses deep neural
networks applied directly to sequences.

\textbf{DeepGO}~\cite{kulmanov2018deepgo} and its successor
\textbf{DeepGOPlus}~\cite{kulmanov2020deepgoplus} train convolutional networks on
sequence $k$-mers and combine the output with sequence similarity scores. They
demonstrate that learned sequence features and alignment evidence are
complementary: the two signals fuse into a more reliable predictor than either
alone.

\textbf{NetGO}~\cite{you2019netgo} augments sequence-based models with
protein–protein interaction network features, achieving state-of-the-art results
in the CAFA3 challenge. Its dependency on interaction data limits applicability
for poorly studied organisms.

A fundamental challenge for supervised methods is the extreme imbalance and
hierarchy of the GO label space. The number of proteins annotated with any
specific GO term decreases rapidly as one descends the ontology DAG; for many
leaf terms, fewer than ten training examples exist. Hierarchical loss functions
and ontology-aware evaluation metrics (Fmax, AUPR) have been developed to address
this~\cite{clark2013information}.

% ────────────────────────────────────────────────────────────
\section{Protein Language Model Embeddings}
\label{sec:related-plm}

The advent of large-scale self-supervised \glspl{plm} has opened a qualitatively
different avenue for function prediction. Models trained on hundreds of millions
of sequences develop internal representations — embeddings — that capture
evolutionary relationships, secondary and tertiary structure propensities, and
functional sites, without any explicit structural or functional supervision.

\textbf{ESM-1b} and \textbf{ESM-2}~\cite{lin2023esm2} (Meta AI) are transformer
models trained on UniRef50/90. ESM-2 scales from 8M to 15B parameters; the
larger variants generate embeddings that outperform BLAST-based transfer at
low-identity regimes. ESM-2 embeddings serve as the primary representation in
PROTEA.

\textbf{ProtTrans}~\cite{elnaggar2022prottrans} provides a suite of encoder
models (ProtBERT, ProtXLNet, ProtT5) trained on BFD and UniRef100 using
distributed HPC resources. ProtT5-XL achieves the best per-residue and per-protein
accuracy among the ProtTrans variants.

\textbf{Annotation transfer via nearest neighbours} is the inference strategy
most naturally paired with embeddings. Given a set of reference proteins with
known annotations and their embeddings, the $k$ nearest reference neighbours of
a query (by cosine or Euclidean distance) are retrieved and their annotations
aggregated to form a prediction. This approach was explored in
\textcite{littmann2021embeddings} and forms the core inference mechanism of
PROTEA.

\textbf{Scalability} is the central engineering challenge for embedding-based
transfer. Swiss-Prot alone contains over 570,000 entries; computing exact cosine
distances against all of them for a large query set requires $O(NQ)$ dot products,
which becomes infeasible at scale. Approximate nearest-neighbour libraries such
as FAISS~\cite{johnson2021faiss} reduce this to $O(Q \sqrt{N})$ with minimal
accuracy loss for IVFFlat indices.

% ────────────────────────────────────────────────────────────
\section{CAFA Benchmark and Evaluation Standards}
\label{sec:related-cafa}

The Critical Assessment of Functional Annotation (CAFA)~\cite{zhou2019cafa3}
challenge provides the community standard for evaluating function prediction
methods. Participants submit predictions for a set of target proteins; ground
truth is established retrospectively from experimental annotations deposited
after the submission deadline.

The primary metric is $F_{\max}$: the maximum $F_1$ score across all prediction
confidence thresholds, computed separately for each GO aspect. AUPR (area under
the precision--recall curve) and Smin (semantic distance) are complementary
measures. PROTEA's evaluation in \cref{chap:evaluation} follows CAFA conventions.

% ────────────────────────────────────────────────────────────
\section{Annotation Platforms and Pipelines}
\label{sec:related-platforms}

Several platforms have been developed to operationalise function prediction at
genome scale.

\textbf{Argot2}~\cite{falda2012argot2} combines BLAST and HMMER hits with a
graph propagation scheme over the GO DAG. It produces confidence scores by
integrating multiple lines of evidence.

\textbf{PANNZER2}~\cite{toronen2018pannzer2} is a web server that combines
BLAST-based nearest-neighbour transfer with a SANS scoring model. It provides
an interactive interface and supports batch submission.

\textbf{eggNOG-mapper}~\cite{cantalapiedra2021eggnogmapper} maps sequences
against clusters of orthologous groups (COGs) and transfers functional
annotations from the eggNOG database. It is widely used in comparative genomics
pipelines.

None of these platforms provides a reproducible, versioned annotation pipeline
backed by a relational data model that records which ontology version, reference
annotation set, and embedding configuration produced each prediction — a gap
that PROTEA addresses directly.

% ────────────────────────────────────────────────────────────
\section{Summary and Positioning of PROTEA}
\label{sec:related-summary}

\Cref{tab:related-comparison} summarises the methods discussed above along four
axes relevant to this thesis.

\begin{table}[htbp]
  \centering
  \caption{Comparison of protein function prediction approaches.}
  \label{tab:related-comparison}
  \resizebox{\textwidth}{!}{%
  \begin{tabular}{lcccc}
    \toprule
    \textbf{Method} & \textbf{Low-identity} & \textbf{Scalable} & \textbf{Versioned} & \textbf{Open source} \\
    \midrule
    BLAST transfer      & \texttimes & \checkmark & \texttimes & \checkmark \\
    HMMER/InterPro      & \checkmark & \checkmark & \texttimes & \checkmark \\
    DeepGO/DeepGOPlus   & \checkmark & \checkmark & \texttimes & \checkmark \\
    PANNZER2            & \texttimes & \checkmark & \texttimes & \checkmark \\
    ESM-2 kNN transfer  & \checkmark & $\sim$     & \texttimes & \checkmark \\
    \textbf{PROTEA}     & \checkmark & \checkmark & \checkmark & \checkmark \\
    \bottomrule
  \end{tabular}
  }
\end{table}

PROTEA occupies a niche that none of the surveyed systems fills: an end-to-end
distributed pipeline that combines \gls{plm} embeddings with optional alignment
and taxonomy features, records all versioning information in a relational schema,
and exposes the full pipeline through a REST API and web interface.

\chapter{System Design}
\label{chap:design}

This chapter describes the architecture of PROTEA at the level of design
decisions. Implementation details — specific libraries, migration tooling,
frontend components — are deferred to \cref{chap:implementation}.

% ────────────────────────────────────────────────────────────
\section{Requirements and Design Goals}
\label{sec:design-goals}

The design of PROTEA is governed by five requirements derived from the
limitations identified in \cref{chap:related}:

\begin{description}
  \item[R1 — Reproducibility.] A prediction produced today must be exactly
    reproducible in the future. This requires recording the ontology version,
    reference annotation set, and embedding model configuration used for every
    prediction run.
  \item[R2 — Scalability.] The system must handle reference sets of hundreds of
    thousands of proteins and query sets of thousands without holding all data
    in memory simultaneously.
  \item[R3 — Separation of concerns.] Domain logic (what to compute), execution
    flow (how jobs are dispatched and tracked), and infrastructure (database,
    message queue) must be independently replaceable.
  \item[R4 — Observability.] Every job must produce a structured audit trail so
    that failures can be diagnosed without replaying the computation.
  \item[R5 — Accessibility.] Researchers without machine-learning infrastructure
    expertise must be able to submit sequences and retrieve predictions through a
    web interface or a REST API.
\end{description}

% ────────────────────────────────────────────────────────────
\section{High-Level Architecture}
\label{sec:architecture}

PROTEA decomposes into four horizontal layers, illustrated in
\cref{fig:architecture}:

\begin{description}
  \item[Presentation layer.] A Next.js single-page application exposes the
    user-facing workflow: uploading FASTA files, triggering pipeline jobs,
    monitoring progress, and downloading results.
  \item[API layer.] A FastAPI application provides a REST interface. All state
    mutations flow through this layer; it writes job requests to PostgreSQL and
    publishes messages to RabbitMQ.
  \item[Worker layer.] A set of Python worker processes consume messages from
    RabbitMQ queues. Workers are stateless with respect to domain logic — they
    resolve and execute registered operations and update job state in the database.
  \item[Data layer.] PostgreSQL (extended with the pgvector extension for vector
    storage) is the single source of truth. RabbitMQ is used exclusively for
    task dispatch; it carries no persistent state.
\end{description}

\begin{figure}[htbp]
  \centering
  \resizebox{\textwidth}{!}{% ============================================================
%  Figure: PROTEA high-level architecture (four layers)
% ============================================================
\begin{tikzpicture}[
  node distance=6mm and 8mm,
  every node/.style={font=\small},
]

  % ── Presentation layer ──────────────────────────────────────
  \node[proteablock, minimum width=46mm] (web) {Next.js Web UI\\\footnotesize FASTA upload, job monitor, TSV export};

  % ── API layer ───────────────────────────────────────────────
  \node[proteablock, minimum width=46mm, below=of web] (api) {FastAPI REST API\\\footnotesize 17 routers · session injection};

  % ── Worker layer ────────────────────────────────────────────
  \node[proteablock, minimum width=22mm, below=14mm of api, xshift=-30mm] (qc) {QueueConsumer\\\footnotesize tracked jobs};
  \node[proteablock, minimum width=22mm, right=4mm of qc]                  (oc) {OperationConsumer\\\footnotesize batch sub-tasks};
  \node[proteablock, minimum width=20mm, right=4mm of oc]                  (gpu) {GPU Inference\\\footnotesize ESM-2 / ESM-C / ProstT5 / ProtT5 / Ankh};

  % ── Data layer ──────────────────────────────────────────────
  \node[proteastore, below=14mm of qc, xshift=10mm]    (pg)  {PostgreSQL 16\\\footnotesize + pgvector};
  \node[proteastore, right=14mm of pg]                 (mq)  {RabbitMQ\\\footnotesize 10 queues};

  % ── Layer brackets ──────────────────────────────────────────
  \begin{scope}[on background layer]
    \node[draw=proteagreen!40, very thick, dashed, rounded corners=4pt,
          fit=(web), inner sep=4mm, label={[font=\footnotesize\itshape, text=proteagreen!70!black]left:Presentation}] {};
    \node[draw=proteagreen!40, very thick, dashed, rounded corners=4pt,
          fit=(api), inner sep=4mm, label={[font=\footnotesize\itshape, text=proteagreen!70!black]left:API}] {};
    \node[draw=proteagreen!40, very thick, dashed, rounded corners=4pt,
          fit=(qc)(oc)(gpu), inner sep=3mm, label={[font=\footnotesize\itshape, text=proteagreen!70!black]left:Workers}] {};
    \node[draw=linkblue!40, very thick, dashed, rounded corners=4pt,
          fit=(pg)(mq), inner sep=3mm, label={[font=\footnotesize\itshape, text=linkblue!70!black]left:Data}] {};
  \end{scope}

  % ── Edges ───────────────────────────────────────────────────
  \draw[proteaarrow] (web) -- node[right, font=\footnotesize]{HTTPS / JSON} (api);
  \draw[proteaarrow] (api.south) -- ++(0,-3mm) -| (qc.north);
  \draw[proteaarrow] (api.south) -- ++(0,-3mm) -| (oc.north);
  \draw[proteaarrow] (api) -- ++(35mm,0) |- (mq.east);

  \draw[proteaarrow] (qc) -- (pg);
  \draw[proteaarrow] (oc) -- (pg);
  \draw[proteaarrow] (oc.south east) -- ++(0,-3mm) -| (gpu.south);
  \draw[proteaarrow] (qc.east) -- (oc.west);

  \draw[proteadashed] (mq.north) to[bend left=18] node[right, font=\footnotesize]{dispatch} (oc.east);
  \draw[proteadashed] (mq.north) to[bend left=10] (qc.east);

\end{tikzpicture}
}
  \caption[High-level architecture of PROTEA]{%
    High-level architecture of PROTEA. The four horizontal layers
    (presentation, API, workers, data) communicate exclusively through well-defined
    interfaces: HTTPS/JSON between the browser and the API, RabbitMQ for task
    dispatch, and PostgreSQL as the single source of truth. GPU inference is
    isolated in dedicated workers so that compute capacity can be scaled
    independently of API throughput.%
  }
  \label{fig:architecture}
\end{figure}

% ────────────────────────────────────────────────────────────
\section{The Operation Abstraction}
\label{sec:operation-abstraction}

The central design principle of PROTEA is the \textit{Operation protocol}: every
unit of domain logic is a class that implements exactly two members:

\begin{itemize}
  \item \texttt{name: str} — a unique string identifier used to route messages.
  \item \texttt{execute(session, payload, *, emit) -> OperationResult} — the
    computation itself.
\end{itemize}

Operations receive a SQLAlchemy session and a validated Pydantic payload; they
report progress through an \texttt{emit} callback that writes \texttt{JobEvent}
rows to the database. Operations do not manage their own sessions, do not
publish messages, and do not know about the queue infrastructure. This strict
interface makes them independently testable and trivially substitutable.

The \texttt{OperationRegistry} maps string names to instances at startup. Workers
resolve the correct operation by name when they receive a message, enabling
new operations to be added without modifying any worker code.

% ────────────────────────────────────────────────────────────
\section{Job Lifecycle and Worker Patterns}
\label{sec:job-lifecycle}

PROTEA distinguishes two worker patterns that correspond to different
observability requirements:

\subsection{QueueConsumer (tracked jobs)}
\label{subsec:queue-consumer}

Long-running, user-visible jobs — inserting proteins, loading ontologies,
computing embeddings, predicting GO terms — are submitted through the API as
\texttt{Job} rows with a JSONB payload. The workflow is:

\begin{enumerate}
  \item The API creates a \texttt{Job} row in state \texttt{QUEUED} and publishes
    the job UUID to the appropriate RabbitMQ queue.
  \item A \texttt{QueueConsumer} worker receives the UUID, opens a
    \textit{claim session}, transitions the job to \texttt{RUNNING}, and flushes
    the \texttt{started\_at} timestamp. The claim session is committed and closed
    before execution begins.
  \item A separate \textit{execute session} resolves the operation and calls
    \texttt{execute()}. On success the job transitions to \texttt{SUCCEEDED}; on
    exception to \texttt{FAILED}, with the error code and message stored on the
    row.
\end{enumerate}

The two-session design ensures that the claim is durable even if the execute
session rolls back. \texttt{JobEvent} rows, written via \texttt{emit()}, form an
append-only audit log that records progress milestones, error context, and
structured diagnostic fields.

\subsection{OperationConsumer (fire-and-forget sub-tasks)}
\label{subsec:operation-consumer}

GPU-intensive batch tasks — running the embedding model on a group of sequences,
or computing kNN predictions for a batch — do not create child \texttt{Job} rows.
Instead, the coordinator operation publishes \textit{operation messages}
(containing the payload directly, not a UUID) to dedicated queues.
\texttt{OperationConsumer} workers execute these messages and report progress
via an atomic increment on the parent job's counter, avoiding write contention
from hundreds of concurrent batches.

\Cref{fig:job-lifecycle} illustrates the two patterns and the queue topology.

\subsection{Parent-Child Job Hierarchy}
\label{subsec:parent-child}

Coordinator operations (\texttt{compute\_embeddings}, \texttt{predict\_go\_terms})
split work across many parallel workers using a parent-child pattern. The
coordinator returns \texttt{OperationResult(deferred=True)}, which signals
\texttt{BaseWorker} to \emph{not} transition the parent job to
\texttt{SUCCEEDED} immediately. Instead, the parent remains in state
\texttt{RUNNING} while batch workers process their messages independently.

The \texttt{Job} model carries two progress fields for this purpose:
\texttt{progress\_total} (set by the coordinator to the number of batches
dispatched) and \texttt{progress\_current} (incremented atomically by each
write worker). When the last write worker detects
$\texttt{progress\_current} = \texttt{progress\_total}$, it closes the parent
job as \texttt{SUCCEEDED}. This atomic counter design avoids write contention
from hundreds of concurrent batches updating the same row.

\subsection{RetryLaterError}
\label{subsec:retry-later}

When a shared resource is unavailable — for instance, the GPU is already
occupied by a running embedding job — an operation raises
\texttt{RetryLaterError} instead of failing permanently. \texttt{BaseWorker}
catches this exception and:

\begin{enumerate}
  \item Resets the job status back to \texttt{QUEUED}.
  \item Writes a \texttt{job.retry\_later} event recording the reason and
    the requested delay.
  \item Re-publishes the job UUID to its queue after the specified delay
    (typically 60 seconds).
\end{enumerate}

The embedding coordinator queue (\texttt{protea.embeddings}) is serialised
precisely because of this mechanism: at most one embedding coordinator runs at
a time, with subsequent requests queued and retried automatically.

\subsection{Cancellation}
\label{subsec:cancellation}

\texttt{POST /jobs/\{id\}/cancel} transitions a \texttt{QUEUED} or
\texttt{RUNNING} job to \texttt{CANCELLED}. If the job is already in a
terminal state (\texttt{SUCCEEDED}, \texttt{FAILED}, or \texttt{CANCELLED})
the request is a no-op. Any queued child jobs are also cancelled atomically
in the same transaction.

Cancellation of a \texttt{RUNNING} job is a \textit{soft cancel}: the database
row is marked \texttt{CANCELLED} immediately, but the worker process is not
interrupted. The worker will still complete its operation and attempt a final
status write; however, because \texttt{CANCELLED} is already committed, the
frontend reflects the cancelled state regardless of how the worker exits.

\begin{figure}[htbp]
  \centering
  \resizebox{\textwidth}{!}{% ============================================================
%  Figure: Job lifecycle and queue topology
%  Source: figures/fig_job_lifecycle.tex @ 57fa4d4
%  Static TikZ diagram. Verify fidelity against:
%    protea/api/app.py, protea/workers/base_worker.py,
%    scripts/manage.sh (10 queues: ping, jobs, training, embeddings,
%    embeddings.batch, embeddings.write, predictions, predictions.batch,
%    predictions.write, evaluations)
% ============================================================
\begin{tikzpicture}[
  node distance=6mm and 8mm,
  every node/.style={font=\small},
]

  % ── State machine (left column) ─────────────────────────────
  \node[proteastate]                       (queued)    {QUEUED};
  \node[proteastate, below=8mm of queued]  (running)   {RUNNING};
  \node[proteastate, below=8mm of running] (succeeded) {SUCCEEDED};
  \node[proteastate, right=10mm of succeeded] (failed)  {FAILED};
  \node[proteastate, left=10mm of succeeded]  (cancel) {CANCELLED};

  \draw[proteaarrow] (queued)  -- node[right, font=\footnotesize]{claim session} (running);
  \draw[proteaarrow] (running) -- node[right, font=\footnotesize]{execute session} (succeeded);
  \draw[proteaarrow] (running.east) -- ++(8mm,0) |- (failed.north);
  \draw[proteaarrow] (running.west) -- ++(-8mm,0) |- (cancel.north);

  % RetryLater self-loop
  \draw[proteadashed] (running.west) to[out=160, in=200, looseness=2.4]
    node[left, font=\footnotesize, align=right]{RetryLaterError\\(re-queue +60s)} (queued.west);

  \node[draw=black!20, dashed, rounded corners=3pt,
        fit=(queued)(running)(succeeded)(failed)(cancel),
        inner sep=4mm,
        label={[font=\footnotesize\itshape]above:Job state machine}] (sm) {};

  % ── Queue topology (right column) ───────────────────────────
  \node[proteablock, minimum width=30mm, right=24mm of queued.north east, anchor=north west] (coord) {Coordinator op\\\footnotesize \texttt{compute\_embeddings}\\\texttt{predict\_go\_terms}};

  \node[proteaqueue, below=8mm of coord]                      (qjobs)    {protea.jobs};
  \node[proteaqueue, below=2mm of qjobs]                      (qemb)     {protea.embeddings};
  \node[proteaqueue, below=2mm of qemb]                       (qebatch)  {protea.embeddings.batch};
  \node[proteaqueue, below=2mm of qebatch]                    (qewrite)  {protea.embeddings.write};
  \node[proteaqueue, below=2mm of qewrite]                    (qpbatch)  {protea.predictions.batch};
  \node[proteaqueue, below=2mm of qpbatch]                    (qpwrite)  {protea.predictions.write};
  \node[proteaqueue, below=2mm of qpwrite]                    (qping)    {protea.ping};
  \node[proteaqueue, below=2mm of qping]                      (qpred)    {protea.predictions};
  \node[proteaqueue, below=2mm of qpred]                      (qtrain)   {protea.training};
  \node[proteaqueue, below=2mm of qtrain]                     (qeval)    {protea.evaluations};

  \node[draw=BurntOrange!40, dashed, rounded corners=3pt,
        fit=(coord)(qjobs)(qeval),
        inner sep=4mm,
        label={[font=\footnotesize\itshape, text=BurntOrange!70!black]above:RabbitMQ topology (10 queues)}] (qg) {};

  % Coordinator → fan-out
  \draw[proteaarrow] (coord.south) -- ++(0,-1mm) -| (qebatch.north);
  \draw[proteaarrow] (coord.south) -- ++(0,-1mm) -| (qpbatch.north);

  % Batch → write chain
  \draw[proteaarrow] (qebatch.east) to[bend left=18] node[right, font=\footnotesize]{vectors} (qewrite.east);
  \draw[proteaarrow] (qpbatch.east) to[bend left=18] node[right, font=\footnotesize]{rows}    (qpwrite.east);

  % Atomic counter on parent job
  \node[draw=proteagreen!50, very thick, fill=proteagreen!8, rounded corners=3pt,
        right=12mm of qewrite, font=\footnotesize, align=center]
        (counter) {atomic\\progress\_current\\++};
  \draw[proteaarrow] (qewrite.east) -- (counter.west);
  \draw[proteaarrow] (qpwrite.east) to[bend right=18] (counter.south);
  \draw[proteadashed] (counter.north) to[bend left=12]
        node[right, font=\footnotesize]{closes parent} (running.east);

\end{tikzpicture}
}
  \caption[Job lifecycle and queue routing in PROTEA]{%
    Job lifecycle and queue routing in PROTEA. The state machine on the left
    is enforced by \texttt{BaseWorker} through two independent SQLAlchemy
    sessions: a short \emph{claim session} that commits the
    \texttt{QUEUED}~$\to$~\texttt{RUNNING} transition, and a longer
    \emph{execute session} that runs the operation and commits the terminal
    state. \texttt{RetryLaterError} resets the job to \texttt{QUEUED} and
    re-publishes it after a delay (typically 60\,s) so that GPU contention
    never produces hard failures. The right column shows the seven RabbitMQ
    queues; coordinator operations fan out work into the
    \texttt{*.batch}~$\to$~\texttt{*.write} chains, where every write worker
    atomically increments \texttt{progress\_current} on the parent job. The
    last write closes the parent as \texttt{SUCCEEDED}, avoiding write
    contention from hundreds of concurrent batches.%
  }
  \label{fig:job-lifecycle}
\end{figure}

% ────────────────────────────────────────────────────────────
\section{Queue Topology}
\label{sec:queue-topology}

PROTEA routes messages across seven RabbitMQ queues, each with a distinct
semantic purpose:

\begin{description}
  \item[\texttt{protea.ping}] Smoke-test queue. Used to verify end-to-end
    connectivity.
  \item[\texttt{protea.jobs}] General-purpose job queue. Handles protein
    insertion, metadata fetch, ontology loading, annotation loading, and the
    coordinator phase of embeddings and predictions.
  \item[\texttt{protea.embeddings}] Serialised embedding coordinator queue. At
    most one coordinator runs at a time; if the GPU is busy the coordinator
    raises \texttt{RetryLaterError} and is re-queued with a 60-second delay.
  \item[\texttt{protea.embeddings.batch}] GPU inference queue. Each message
    carries a batch of sequences; \texttt{OperationConsumer} workers run the
    forward pass of the protein language model and publish results to the
    write queue.
  \item[\texttt{protea.embeddings.write}] Bulk insert queue. Workers receive
    computed embedding vectors and write them to the \texttt{SequenceEmbedding}
    table via pgvector.
  \item[\texttt{protea.predictions.batch}] kNN prediction queue. Each message
    carries a batch of query embeddings; workers search the reference index and
    transfer GO annotations.
  \item[\texttt{protea.predictions.write}] Bulk prediction insert queue. Workers
    write \texttt{GOPrediction} rows for a batch of query--reference pairs.
\end{description}

This topology separates GPU-bound inference, I/O-bound database writes, and
CPU-bound coordination into independent queues that can be scaled independently
(e.g.\ \texttt{manage.sh scale protea.predictions.batch 4}).

% ────────────────────────────────────────────────────────────
\section{Data Model}
\label{sec:data-model}

The relational schema is organised into five logical groups:

\subsection{Sequence and Protein}
\label{subsec:dm-sequence}

\texttt{Sequence} stores deduplicated amino acid strings, keyed by MD5 hash.
\texttt{Protein} stores one row per UniProt accession; multiple accessions
(canonical and isoforms) may reference the same \texttt{Sequence}. This
deduplication prevents redundant embedding computation.

\subsection{Ontology and Annotations}
\label{subsec:dm-ontology}

\texttt{OntologySnapshot} records one complete GO release (OBO format), versioned
by the \texttt{obo\_version} string found in the OBO file header. \texttt{GOTerm}
and \texttt{GOTermRelationship} store the DAG within a snapshot.

\texttt{AnnotationSet} groups a batch of \texttt{ProteinGOAnnotation} rows by
source (\texttt{goa} or \texttt{quickgo}) and links them to an
\texttt{OntologySnapshot}. This design enables side-by-side comparison of
annotation sets from different sources or dates.

\subsection{Embeddings}
\label{subsec:dm-embeddings}

\texttt{EmbeddingConfig} is an immutable recipe: it records the model identifier,
chunking strategy, pooling method, and any other parameters that affect the
embedding geometry. Its UUID primary key is stable; changing any parameter
creates a new configuration.

\texttt{SequenceEmbedding} stores one pgvector \texttt{VECTOR} per
(\texttt{sequence}, \texttt{config}, \texttt{chunk}) triple. Chunking support
allows sequences that exceed the model's context window to be split into
overlapping segments.

\subsection{Predictions}
\label{subsec:dm-predictions}

\texttt{PredictionSet} is the result container, linking a query set, an embedding
configuration, an annotation set, and an ontology snapshot.
\texttt{GOPrediction} stores one row per (query protein, GO term) pair, including
the cosine distance from the nearest reference neighbour and the optional
feature-engineering columns described in \cref{sec:feature-engineering}.

\subsection{Reranker Models}
\label{subsec:dm-reranker}

\texttt{RerankerModel} records a trained LightGBM booster promoted from the
research sandbox described in \cref{sec:reranker-promotion}. The row carries
the artefact URI resolved by PROTEA's storage abstraction
(\texttt{file:///\ldots} for local filesystem or \texttt{s3://\ldots} for
MinIO), the \texttt{feature\_schema\_sha} fingerprint used to validate
feature-set alignment at inference time, the originating embedding
configuration and ontology snapshot, and the full experiment specification
YAML used to reproduce the run. The inline \texttt{model\_data} column is
retained only as a nullable compatibility slot for rows produced by the
pre-abstraction pipeline; new registrations never populate it.

\subsection{Query Sets}
\label{subsec:dm-querysets}

\texttt{QuerySet} represents a user-uploaded FASTA dataset submitted for custom
prediction queries. Each \texttt{QuerySetEntry} row preserves the original
accession header from the FASTA file and links to the deduplicated
\texttt{Sequence} row, reusing existing sequences when the amino-acid string is
already present in the database. Query sets are created via
\texttt{POST /query-sets} and can be referenced in subsequent embedding
computation and prediction jobs.

\subsection{Jobs}
\label{subsec:dm-jobs}

\texttt{Job} implements a state machine with states \texttt{QUEUED},
\texttt{RUNNING}, \texttt{SUCCEEDED}, \texttt{FAILED}, and \texttt{CANCELLED}.
\texttt{JobEvent} is an append-only log of timestamped, structured events
emitted during execution. Both tables use JSONB columns for extensible payloads
and metadata.

% ────────────────────────────────────────────────────────────
\section{Embedding Pipeline}
\label{sec:embedding-pipeline}

The embedding pipeline runs in three phases coordinated through the queue
topology described above:

\begin{enumerate}
  \item \textbf{Coordinator} (\texttt{compute\_embeddings}): queries the database
    for all sequences in the target set that do not yet have an embedding for the
    requested \texttt{EmbeddingConfig}; partitions them into fixed-size batches;
    publishes one operation message per batch to \texttt{protea.embeddings.batch}.
  \item \textbf{Batch inference} (\texttt{compute\_embeddings\_batch}): loads the
    protein language model on the configured device; tokenises the sequence
    batch; runs a forward pass; mean-pools the per-residue representations into
    a single vector per chunk; publishes the computed vectors to
    \texttt{protea.embeddings.write}.
  \item \textbf{Store} (\texttt{store\_embeddings}): receives the computed vectors
    and performs a bulk upsert into \texttt{SequenceEmbedding}.
\end{enumerate}

The coordinator is serialised (one at a time per the \texttt{protea.embeddings}
queue) to prevent concurrent model loads from exhausting GPU memory.
Batch and write workers can be scaled horizontally.

% ────────────────────────────────────────────────────────────
\section{Prediction Pipeline}
\label{sec:prediction-pipeline}

The prediction pipeline mirrors the structure of the embedding pipeline:

\begin{enumerate}
  \item \textbf{Coordinator} (\texttt{predict\_go\_terms}): loads the reference
    embeddings (all proteins in the annotation set that have embeddings) into
    a process-level cache compressed to float16; partitions the query set into
    batches; publishes one operation message per batch to
    \texttt{protea.predictions.batch}.
  \item \textbf{Batch prediction} (\texttt{predict\_go\_terms\_batch}): for each
    query, retrieves the $k$ nearest reference proteins using the configured
    search backend (NumPy exact cosine search or FAISS approximate search);
    transfers GO annotations from each neighbour; optionally computes pairwise
    alignment statistics (\gls{nw} and \gls{sw} via the parasail library) and
    taxonomic distance (via ete3 and the NCBI taxonomy); publishes the results
    to \texttt{protea.predictions.write}.
  \item \textbf{Store} (\texttt{store\_predictions}): bulk-inserts
    \texttt{GOPrediction} rows.
\end{enumerate}

The reference cache is held at the process level and is keyed by
(embedding config, annotation set). It holds at most one entry to bound memory
usage; a restart reloads from the database automatically.

When a job payload includes a \texttt{reranker\_model\_id}, the batch worker
performs one additional pass after the kNN stage: the booster referenced by
the registered \texttt{RerankerModel} is loaded through a sha-keyed on-disk
cache, the candidate rows are scored, and the resulting \texttt{reranker\_score}
values replace the kNN distance as the ranking key. The pass is skipped and a
\texttt{reranker.schema\_mismatch} event is emitted if the feature fingerprint
computed from the live pipeline does not match the one registered with the
booster; the job otherwise proceeds to the store stage unchanged. The
promotion mechanics are described in \cref{sec:reranker-promotion}.

% ────────────────────────────────────────────────────────────
\section{REST API Design}
\label{sec:api-design}

The API follows REST conventions. Resources are grouped into six routers:

\begin{description}
  \item[\texttt{/proteins}] Insert proteins from UniProt accession lists or FASTA.
  \item[\texttt{/annotations}] Load ontology snapshots and annotation sets.
  \item[\texttt{/embeddings}] Manage embedding configurations, trigger computation,
    submit prediction jobs, and export results.
  \item[\texttt{/query-sets}] Upload FASTA files and manage user query datasets.
  \item[\texttt{/jobs}] Track job state and stream structured events.
  \item[\texttt{/admin}] Administrative maintenance tasks.
\end{description}

\begin{sloppypar}
Session injection follows FastAPI's dependency system:
\texttt{app.state.session\_factory} is set at startup and injected into each
request handler, keeping router code free of infrastructure imports.
\end{sloppypar}

\chapter{Implementation}
\label{chap:implementation}

This chapter describes the concrete technology choices and key implementation
details behind the design presented in \cref{chap:design}. Where decisions are
recorded in a numbered \gls{adr}, the identifier is cited in-text so the decision
can be traced back to its rationale.

% ────────────────────────────────────────────────────────────
\section{Multi-Repository Project Structure}
\label{sec:project-structure}

PROTEA is realised as a family of seven code repositories plus the thesis
manuscript (ADR D1). This \emph{Structure C} arrangement allows each
functional group to evolve, release, and be installed independently while
still composing into a coherent platform at runtime.

\begin{description}
  \item[\texttt{PROTEA}] The platform core: FastAPI application, worker
    processes, SQLAlchemy ORM, RabbitMQ consumers, and the operation
    registry. All operational entry points live here.
  \item[\texttt{protea-contracts}] Abstract base classes and shared payload
    schemas that decouple the platform from its plugins. Only light
    dependencies (Pydantic, NumPy); no torch or ORM imports.
  \item[\texttt{protea-backends}] Protein language model inference plugins
    (ESM-2, T5/ProstT5, Ankh, ESM3c). Described in detail in
    \cref{sec:backend-plugins}.
  \item[\texttt{protea-method}] Pure inference library: kNN search, feature
    enrichment, reranker apply, and PCA cache. No FastAPI or SQLAlchemy
    dependencies, enabling standalone installation via
    \texttt{pip install protea-method} (ADR D15).
  \item[\texttt{protea-sources}] Annotation-source plugins (QuickGO, GAF
    streams) discovered via the \texttt{protea.sources} entry-points group.
  \item[\texttt{protea-runners}] Experiment-runner plugins, including the
    LightGBM training runner (\texttt{protea-runners.lightgbm}) that replaces
    the former inline \texttt{TrainRerankerOperation} (ADR D9, now obsolete).
  \item[\texttt{cafaeval-protea}] A maintained fork of the
    \texttt{cafaeval} evaluation library used by \texttt{run\_cafa\_evaluation}.
\end{description}

Within the \texttt{PROTEA} repository the internal layout mirrors the
architectural layers:

\begin{description}
  \item[\texttt{protea/api/}] FastAPI application and seventeen routers
    spanning the full operational surface of the platform.
  \item[\texttt{protea/core/}] Pure domain logic: the \texttt{Operation}
    protocol, \texttt{OperationRegistry}, all registered operations, the
    kNN search shim, feature engineering, and shared HTTP helpers.
  \item[\texttt{protea/infrastructure/}] SQLAlchemy ORM models,
    RabbitMQ consumer and publisher, the session context manager, and the
    configuration loader.
  \item[\texttt{apps/web/}] The Next.js 16 frontend application.
  \item[\texttt{scripts/}] Operational tooling: \texttt{manage.sh} (process
    manager), \texttt{worker.py} (worker entry point), \texttt{init\_db.py}
    (schema initialisation), and \texttt{run\_one\_job.py} (manual job
    runner for debugging).
\end{description}

Dependency direction is enforced: \texttt{api} and workers depend on
\texttt{core} and \texttt{infrastructure}; \texttt{core} has no dependency
on either of the other two.

% ────────────────────────────────────────────────────────────
\section{Technology Stack}
\label{sec:stack}

\Cref{tab:stack} summarises the main components of the PROTEA stack.

\begin{table}[htbp]
  \centering
  \caption{PROTEA technology stack.}
  \label{tab:stack}
  \begin{tabular}{lll}
    \toprule
    \textbf{Component}   & \textbf{Technology}         & \textbf{Version} \\
    \midrule
    API framework        & FastAPI                     & 0.115+           \\
    ORM / migrations     & SQLAlchemy 2.0 + Alembic    & 2.0 / 1.13       \\
    Database             & PostgreSQL 16 + pgvector    & 16 / 0.7         \\
    Message broker       & RabbitMQ + aio-pika         & 3.x / 9.x        \\
    Data validation      & Pydantic                    & 2.x              \\
    Protein LM inference & Hugging Face Transformers   & 4.x              \\
    Alignment            & parasail-python (BLOSUM62)  & 1.x              \\
    Taxonomy             & ete3 + NCBITaxa             & 3.x              \\
    ANN search           & NumPy / FAISS               & n/a              \\
    Frontend             & Next.js 16 + Tailwind CSS   & 16 / 4           \\
    Dependency mgmt.     & Poetry                      & 1.x              \\
    \bottomrule
  \end{tabular}
\end{table}

All Python dependencies are declared in \texttt{pyproject.toml} with pinned
version ranges; \texttt{poetry.lock} guarantees reproducible installs across
environments. The \texttt{dev} dependency group adds pytest, pytest-cov, and
related tooling without polluting the production image.

% ────────────────────────────────────────────────────────────
\section{Database Schema and Migrations}
\label{sec:database}

The database schema is managed with Alembic. Migration scripts are generated
from SQLAlchemy ORM metadata via \texttt{alembic revision --autogenerate} and
stored under \texttt{alembic/versions/}. The Alembic \texttt{env.py} reads the
database URL from \texttt{protea/config/system.yaml} (or the \texttt{PROTEA\_DB\_URL}
environment variable), ensuring consistency with the application configuration.

The pgvector extension is enabled at database initialisation time:

\begin{lstlisting}[language=SQL, caption={Enabling pgvector.}]
CREATE EXTENSION IF NOT EXISTS vector;
\end{lstlisting}

\texttt{SequenceEmbedding} columns are declared as \texttt{HALFVEC(n)} (migrated
from \texttt{VECTOR} in April~2026) where $n$ is the embedding dimension
(e.g.\ 1280 for ESM-2 650M, 1024 for ProtT5-XL). Although pgvector supports
index types for nearest-neighbour queries (HNSW, IVFFlat), PROTEA
intentionally does not use them for kNN search: querying through pgvector at
the scale of hundreds of thousands of vectors introduces unacceptable latency
(ADR 001). Instead, reference embeddings are loaded into Python (NumPy or
\gls{faiss} via \texttt{protea-method}) at prediction time.

\subsection{Session Management}
\label{subsec:sessions}

All database access goes through \texttt{session\_scope()}, a context manager
that commits on normal exit and rolls back on exception:

\begin{lstlisting}[language=Python, caption={Session context manager.}]
@contextmanager
def session_scope(factory):
    session = factory()
    try:
        yield session
        session.commit()
    except Exception:
        session.rollback()
        raise
    finally:
        session.close()
\end{lstlisting}

The two-session pattern in \texttt{BaseWorker} (claim session, execute session)
relies on this primitive: the claim commits independently of the execute, so a
crash during execution never rolls back the \texttt{RUNNING} state transition
(ADR 002).

% ────────────────────────────────────────────────────────────
\section{Backend Plugin System}
\label{sec:backend-plugins}

Protein language model inference is provided by the \texttt{protea-backends}
package, a separate repository that ships plugins conforming to the
\texttt{EmbeddingBackend} abstract base class defined in
\texttt{protea-contracts}. This separation keeps heavy ML dependencies (torch,
transformers, the ESM SDK) out of the PROTEA core image: the platform
discovers and loads backend plugins lazily via Python entry points, paying no
import cost for a backend that is not used in a given deployment.

\subsection{The \texttt{EmbeddingBackend} Contract}
\label{subsec:embedding-backend-contract}

The ABC mandates two methods:

\begin{lstlisting}[language=Python, caption={EmbeddingBackend interface (protea-contracts).}]
class EmbeddingBackend(ABC):
    name: str  # matched against EmbeddingConfig.model_backend

    def load_model(self, model_name: str, device: str, emit: Any
                   ) -> tuple[Any, Any]: ...

    def embed_batch(self, model, tokenizer, sequences: list[str],
                    *, emit, layers=None, layer_agg="mean",
                    pooling="mean") -> np.ndarray: ...
\end{lstlisting}

\texttt{load\_model} returns a \texttt{(model, tokenizer)} pair (the tokenizer
slot is \texttt{None} for ESM3c which uses its own SDK encoder).
\texttt{embed\_batch} returns a \texttt{(batch\_size, dim)} float16 NumPy
matrix. Torch imports are deferred inside these methods so that the plugin
module is cheap to import at entry-point discovery time.

\subsection{Plugin Discovery}
\label{subsec:plugin-discovery}

The platform discovers installed backends through the
\texttt{protea.backends} entry-points group. The discovery logic in
\texttt{protea/core/plugins.py} calls
\texttt{importlib.metadata.entry\_points(group="protea.backends")} once per
process, verifies that each loaded plugin's \texttt{name} attribute matches
the entry-point key (a mismatch raises \texttt{RuntimeError} immediately), and
caches the resulting name-to-plugin map for the lifetime of the process:

\begin{lstlisting}[language=Python, caption={Plugin discovery (protea.core.plugins).}]
def discover_plugins(group: str) -> dict[str, Any]:
    cache = _PLUGIN_CACHES.get(group)
    if cache is None:
        new_cache = {}
        for ep in entry_points(group=group):
            plugin = ep.load()
            if getattr(plugin, "name", None) != ep.name:
                raise RuntimeError(f"Plugin name mismatch: {ep.name!r}")
            new_cache[ep.name] = plugin
        _PLUGIN_CACHES[group] = new_cache
        cache = new_cache
    return cache
\end{lstlisting}

The \texttt{protea-backends} package declares four entry points in its
\texttt{pyproject.toml}:

\begin{lstlisting}[caption={Backend entry-point declarations (pyproject.toml).}]
[tool.poetry.plugins."protea.backends"]
esm   = "protea_backends.esm:plugin"
t5    = "protea_backends.t5:plugin"
ankh  = "protea_backends.ankh:plugin"
esm3c = "protea_backends.esm3c:plugin"
\end{lstlisting}

Adding a new backend requires only implementing \texttt{EmbeddingBackend},
declaring one entry point, and publishing the package; the platform requires
no code change.

\subsection{Implemented Backends}
\label{subsec:implemented-backends}

\begin{description}
  \item[\texttt{esm} / \texttt{auto}] HuggingFace \texttt{EsmModel}
    checkpoints (ESM-2 family: 8M to 3B parameters). Sequences are processed
    one at a time to avoid variable-length padding memory waste. The CLS token
    (position 0) and EOS token (last valid position) are stripped before
    residue-level pooling.

  \item[\texttt{t5}] HuggingFace \texttt{T5EncoderModel} (ProtT5-XL,
    ProstT5). Sequences are processed as a padded batch. Ambiguous amino acids
    (\texttt{U}/\texttt{Z}/\texttt{O}/\texttt{B}) are replaced with \texttt{X}
    before SentencePiece tokenisation. ProstT5 mode (prepending the
    \texttt{<AA2fold>} prefix) is auto-detected from the model name; the prefix
    residue is stripped from the output so that \texttt{residues[0]}
    corresponds to amino acid~0 in both the \texttt{esm} and \texttt{t5}
    backends.

  \item[\texttt{ankh}] Ankh-base and Ankh-large. Uses the same batched T5
    pipeline as \texttt{t5} but with two deviations: no \texttt{<AA2fold>}
    prefix, and per-residue character list tokenisation
    (\texttt{is\_split\_into\_words=True}) because Ankh's SentencePiece
    tokeniser maps a literal space to \texttt{<unk>}, rendering the
    space-joined path of ProtT5 unsuitable.

  \item[\texttt{esm3c}] ESM3c (ESMC family) via the ESM SDK rather than
    HuggingFace Transformers. No external tokenizer; the SDK's
    \texttt{ESMProtein} API and \texttt{LogitsConfig(return\_hidden\_states=True)}
    are used. Hidden states are returned as a single stacked tensor; BOS and
    EOS tokens are stripped identically to the ESM-2 path.
\end{description}

% ────────────────────────────────────────────────────────────
\section{Operations}
\label{sec:operations}

\begin{sloppypar}
The \texttt{build\_operation\_registry()} factory in
\texttt{protea/core/\allowbreak{}operation\_catalog.py} registers the full set of operations
at startup; both the API process and every worker call the same factory so
the set is always in sync.
\end{sloppypar}
The user-facing operations are:

\begingroup\sloppy
\begin{description}
  \item[\texttt{ping}] Smoke test. Returns immediately with a success result.
  \item[\texttt{insert\_proteins}] Paginates the UniProt REST API using
    cursor-based FASTA streaming. Sequences are deduplicated by MD5 hash before
    upsert; proteins are upserted by accession. Exponential backoff with jitter
    and \texttt{Retry-After} header handling are implemented in the shared
    \texttt{UniProtHttpMixin}.
  \item[\texttt{fetch\_uniprot\_metadata}] Downloads TSV functional annotation
    data from UniProt and upserts \texttt{ProteinUniProtMetadata} rows keyed by
    canonical accession.
  \item[\texttt{load\_ontology\_snapshot}] Downloads a GO OBO file and populates
    \texttt{OntologySnapshot}, \texttt{GOTerm}, and \texttt{GOTermRelationship}
    rows. The \texttt{obo\_version} field carries a unique constraint so that
    re-importing the same release is idempotent.
  \item[\texttt{load\_goa\_annotations}] Bulk-loads a \gls{gaf} file.
    Annotations are filtered against canonical accessions present in the
    database, avoiding orphaned foreign keys.
  \item[\texttt{load\_quickgo\_annotations}] Streams GO annotations from the
    QuickGO bulk download API (paginated TSV). Supports optional ECO evidence
    code mapping and per-page commits to bound transaction size.
  \item[\texttt{generate\_evaluation\_set}] Materialises an \texttt{EvaluationSet}
    by snapshotting ground-truth annotations for a target ontology snapshot and
    accession set; consumed downstream by \texttt{run\_cafa\_evaluation}.
  \item[\texttt{run\_cafa\_evaluation}] Runs the standalone \texttt{cafaeval}
    fork against a \texttt{PredictionSet} and persists \texttt{EvaluationResult}
    rows (Fmax, AuPRC, coverage, per-aspect breakdowns). Isolated on its own
    queue so long evaluations do not block the general jobs queue. The same
    evaluation protocol is applied to external tools via the companion
    \texttt{scripts/evaluate\_external\_tool.py} helper; the concrete
    procedure for eggNOG-mapper is documented in \cref{app:external-comparison}.
  \item[\texttt{compute\_embeddings}] Coordinator operation described in
    \cref{sec:impl-embedding-pipeline}.
  \item[\texttt{predict\_go\_terms}] Coordinator operation described in
    \cref{sec:impl-prediction-pipeline}.
  \item[\texttt{export\_research\_dataset}] Writes the frozen train/eval
    parquet dataset and its \texttt{manifest.json} used by the external
    research sandbox described in \cref{sec:reranker-promotion}. The operation
    runs the same kNN and feature-generation pipeline as the internal
    \texttt{training\_dump\_helpers} module, skips the LightGBM training stage,
    and publishes the resulting artefacts through the storage abstraction.
  \item[\texttt{run\_interproscan\_batch}] Resolves target proteins (via
    a \texttt{query\_set\_id} or an explicit accession list), invokes the
    \texttt{InterProSource} subprocess plugin (IP.1b), parses the TSV output,
    and persists \texttt{InterProAnnotation} rows. Described in detail in
    \cref{sec:interpro-integration}.
\end{description}
\endgroup

% ────────────────────────────────────────────────────────────
\section{InterPro Annotation Integration}
\label{sec:interpro-integration}

InterProScan domain hits are persisted in PROTEA through the \texttt{interpro\_annotation}
table (IP.2, PR~\#306) and loaded via the \texttt{run\_interproscan\_batch} operation
(IP.3, PR~\#307).

\subsection{The \texttt{InterProAnnotation} ORM Model}
\label{subsec:interpro-orm}

The \texttt{InterProAnnotation} table stores one InterProScan domain hit per row.
Each row carries:

\begin{description}
  \item[\texttt{id}] UUID surrogate primary key.
  \item[\texttt{protein\_accession}] Foreign key to \texttt{protein.accession}
    with \texttt{ON DELETE CASCADE}, following the same pattern as
    \texttt{protein\_go\_annotation}.
  \item[\texttt{source\_db}, \texttt{accession}] Source database identifier
    (e.g.\ Pfam, PANTHER) and the domain accession within that database.
  \item[\texttt{start}, \texttt{end}] Residue coordinates, both \texttt{NOT NULL
    INTEGER}. Two database-level \texttt{CHECK} constraints enforce
    \texttt{start $\geq$ 1} and \texttt{end $\geq$ start}, preventing
    invalid coordinates from entering the schema.
  \item[\texttt{ipr\_release}] InterPro release version tag so annotations
    from different InterProScan runs can coexist for the same protein.
  \item[\texttt{evidence}] Nullable free-form text (e.g.\ \texttt{"IEA"} or a
    Pfam E-value string).
\end{description}

The Alembic migration for IP.2 serves as the merge-point for two diverged
develop heads (\texttt{t3.1a\_goprediction\_features\_jsonb} from PR~\#299
and \texttt{add\_api\_key\_table} from PR~\#296), giving subsequent migrations
a single \texttt{down\_revision} chain to follow.

\subsection{The \texttt{run\_interproscan\_batch} Operation}
\label{subsec:run-interproscan}

\texttt{run\_interproscan\_batch} (IP.3) orchestrates the end-to-end InterProScan
pipeline:

\begin{enumerate}
  \item Resolve target proteins either from a \texttt{query\_set\_id}
    (intersected with the \texttt{protein} table via \texttt{QuerySetEntry})
    or from an explicit \texttt{accessions} list.
  \item Filter already-annotated proteins using the optional
    \texttt{ipr\_release\_floor} parameter: proteins whose latest persisted
    \texttt{ipr\_release} is lexicographically $\geq$ the floor are skipped,
    giving the operation natural resume semantics when workers crash mid-batch.
  \item Chunk the target list, write temporary FASTA files, and invoke the
    \texttt{InterProSource} plugin (IP.1b) as a subprocess. Plugin events are
    absorbed by a no-op emit callback; the operation's own chunk-level events
    are sufficient for the frontend job monitor.
  \item Parse the InterProScan TSV output via the IP.1a parser and bulk-insert
    \texttt{InterProAnnotation} rows into the database.
\end{enumerate}

The operation is registered in \texttt{build\_operation\_registry()} and is
therefore dispatchable via \texttt{POST /jobs} with
\texttt{\{operation: "run\_interproscan\_batch", payload: \{\ldots\}\}} through
the standard job worker, with no dedicated endpoint required.

% ────────────────────────────────────────────────────────────
\section{Embedding Pipeline Implementation}
\label{sec:impl-embedding-pipeline}

Embedding computation is split across two operation classes that communicate
through the RabbitMQ message broker rather than sharing a process.

\paragraph{Coordinator: \texttt{ComputeEmbeddingsOperation}.}
Running on the \texttt{protea.embeddings} queue (serialised: one coordinator
at a time), this operation resolves the \texttt{EmbeddingConfig}, queries the
database for sequences that do not yet have a stored embedding for that
configuration (when \texttt{skip\_existing=true}), and partitions the results
into batches of \texttt{sequences\_per\_job} identifiers. For each batch it
publishes a \texttt{ComputeEmbeddingsBatchPayload} message to the
\texttt{protea.embeddings.batch} queue. No child \texttt{Job} rows are created
in the database for batch messages; progress is tracked by the coordinator
via \texttt{update\_parent\_progress}.

\paragraph{Batch worker: \texttt{ComputeEmbeddingsBatchOperation}.}
Running on \texttt{protea.embeddings.batch} workers (one or more, \gls{gpu}-bound),
this operation receives a list of sequence primary keys and a device string.
Control flow:

\begin{enumerate}
  \item Resolve the \texttt{EmbeddingConfig} and load the appropriate backend
    via \texttt{discover\_plugins("protea.backends")} (see
    \cref{sec:backend-plugins}). A backend is loaded at most once per process.
  \item Load model and tokenizer from the backend's \texttt{load\_model} method.
  \item Run the backend's \texttt{embed\_batch} per-sequence or per-batch
    depending on backend type. ESM backends iterate one sequence at a time;
    T5 and Ankh backends process a padded batch; ESM3c uses the ESM SDK.
  \item Each backend returns a \texttt{list[list[ChunkEmbedding]]} where
    the inner list holds one \texttt{ChunkEmbedding} per residue span when
    chunking is enabled, or a single element covering the full sequence
    otherwise. The \texttt{ChunkEmbedding} dataclass carries
    \texttt{chunk\_index\_s}, \texttt{chunk\_index\_e} (0-based, exclusive),
    and a float32 \texttt{vector}.
  \item Build \texttt{SequenceEmbedding} ORM rows from the chunk list and
    publish them to \texttt{protea.embeddings.write} for bulk upsert by a
    dedicated write worker.
\end{enumerate}

Layer selection and aggregation are controlled by \texttt{EmbeddingConfig}
fields. The reverse-index convention (\texttt{layer\_indices=[0]} means the
last layer) matches the convention used by the predecessor FANTASIA system.
Optional per-residue L2 normalisation (\texttt{normalize\_residues}) and
final embedding L2 normalisation (\texttt{normalize}) can be toggled per
configuration.

% ────────────────────────────────────────────────────────────
\section{Prediction Pipeline Implementation}
\label{sec:impl-prediction-pipeline}

GO term prediction mirrors the embedding pipeline's coordinator-batch
split.

\paragraph{Coordinator: \texttt{PredictGOTermsOperation}.}
Running on \texttt{protea.predictions} (serialised), this operation creates a
\texttt{PredictionSet} row, resolves any requested \texttt{RerankerModel},
snapshots the artefact URI and feature-schema fingerprint onto each batch
message, and dispatches one \texttt{PredictGOTermsBatchPayload} per query
partition to \texttt{protea.predictions.batch}.

\paragraph{Batch worker: \texttt{PredictGOTermsBatchOperation}.}
After a refactor guided by ADR D31, the operation's execute method was reduced
to roughly 50 lines of orchestration code; the heavy logic was extracted into
\texttt{\_AspectSeparatedKnnRunner} (347~LOC, 10 methods all within the 60~LOC
budget). The runner encapsulates:

\begin{enumerate}
  \item Loading reference embeddings from PostgreSQL into a process-level cache
    (float16, one array per GO aspect) using \texttt{protea.core.disk\_cache}.
  \item Calling \texttt{search\_knn} from \texttt{protea.core.knn\_search}
    (a backwards-compatible shim over \texttt{protea\_method.knn\_search})
    with the configured backend (\texttt{numpy} or \texttt{faiss}) and distance
    metric.
  \item Transferring GO annotations from the $k$ nearest reference proteins
    to the query, filtered by evidence code and GO aspect.
  \item Optionally enriching each (query, reference) pair with alignment and
    taxonomy features (\cref{sec:feature-engineering}).
  \item Optionally applying the registered \texttt{RerankerModel} booster to
    re-score the candidate set before persisting predictions.
\end{enumerate}

Results are published to \texttt{protea.predictions.write} for bulk insert
of \texttt{GOPrediction} rows.

\subsection{KNN Search Shim and \texttt{protea-method}}
\label{subsec:knn-shim}

During the F2C extraction phase the kNN backend code was moved into the
standalone \texttt{protea-method} library, which ships without FastAPI or
SQLAlchemy dependencies and can be installed independently (ADR D15). The
module \texttt{protea/core/knn\_search.py} is a thin shim that re-exports
\texttt{protea\_method.knn\_search.search\_knn} and forwards the
\texttt{OperationTuning.numpy\_query\_chunk} configuration knob via the
\texttt{PROTEA\_METHOD\_NUMPY\_QUERY\_CHUNK} environment variable on first
call, preserving backwards compatibility for all PROTEA call sites.

% ────────────────────────────────────────────────────────────
\section{Feature Engineering}
\label{sec:feature-engineering}

When a prediction job is submitted with \texttt{compute\_alignments=true} or
\texttt{compute\_taxonomy=true}, the batch prediction operation augments each
query-reference pair with additional numerical features.

\subsection{Pairwise Alignment}
\label{subsec:fe-alignment}

For each (query, reference) pair, PROTEA computes global (\gls{nw}) and local
(\gls{sw}) alignments using the \texttt{parasail} library with BLOSUM62
substitution scores and gap-open/gap-extend penalties of $-11$/$-1$.
The following derived statistics are stored as columns on \texttt{GOPrediction}:

\begin{itemize}
  \item \texttt{identity\_nw/sw}: fraction of aligned positions with identical
    residues.
  \item \texttt{similarity\_nw/sw}: fraction of aligned positions with positive
    BLOSUM62 score.
  \item \texttt{alignment\_score\_nw/sw}: raw alignment score.
  \item \texttt{gaps\_pct\_nw/sw}: fraction of the alignment that is gap.
  \item \texttt{alignment\_length\_nw/sw}: number of aligned columns.
  \item \texttt{length\_query / length\_ref}: sequence lengths.
\end{itemize}

\subsection{Taxonomic Distance}
\label{subsec:fe-taxonomy}

For each (query, reference) pair where taxonomy IDs are available from UniProt
metadata, the NCBI taxonomy tree is consulted via ete3 to compute:

\begin{itemize}
  \item \texttt{taxonomic\_lca}: the taxon ID of the \gls{lca}.
  \item \texttt{taxonomic\_distance}: the total edge count between the two taxa
    through their \gls{lca}.
  \item \texttt{taxonomic\_common\_ancestors}: the number of shared ancestors.
  \item \texttt{taxonomic\_relation}: a categorical label:
    \textit{same\_species}, \textit{same\_genus}, \textit{same\_family},
    \textit{same\_order}, \textit{same\_class}, \textit{same\_phylum},
    \textit{same\_kingdom}, or \textit{distant}.
\end{itemize}

Taxonomy lookups are memoised with an LRU cache keyed by taxon ID pair to avoid
redundant tree traversals across a batch.

\subsection{Lineage Features}
\label{subsec:fe-lineage}

Lineage features (T-RES.1) capture the relationship between a candidate GO
term and the protein's pre-cutoff known annotation set, a signal that the full
alignment and taxonomy features do not encode. The
\texttt{protea\_method.lineage.compute\_lineage\_features} function adds four
columns to each prediction record:

\begin{itemize}
  \item \texttt{lineage\_is\_ancestor\_of\_known}: 1.0 if the candidate is an
    ancestor of at least one term already known for the query protein.
  \item \texttt{lineage\_is\_descendant\_of\_known}: 1.0 if the candidate has
    at least one known term as an ancestor.
  \item \texttt{lineage\_ancestor\_of\_count}: count of known terms whose
    ancestor closure contains the candidate.
  \item \texttt{lineage\_descendant\_of\_count}: count of known terms in the
    candidate's own ancestor closure.
\end{itemize}

The function operates on the pre-built GO DAG parent map
(\texttt{go\_id $\to$ list[parent\_ids]}, covering both \texttt{is\_a} and
\texttt{part\_of} edges) and a \texttt{known\_by\_protein} mapping; it
performs no database access. Ancestor closures are computed via iterative
BFS and memoised per term so the per-batch overhead is proportional to the
number of distinct candidate terms rather than the number of
query-reference pairs. The four columns were registered in
\texttt{protea-contracts} v0.3.0 as the \texttt{lineage} feature family
and bound in PROTEA's \texttt{FeatureRegistry} via a lazy import shim
(T-RES.1b, PR~\#304). The feature-schema fingerprint consequently advances
to a new SHA, and any LightGBM booster trained before T-RES.1b ignores the
new columns via the native missing-value branch without requiring retraining.

% ────────────────────────────────────────────────────────────
\section{Approximate Nearest-Neighbour Search}
\label{sec:ann-search}

The kNN search backend is configurable per prediction job via the
\texttt{search\_backend} payload field. Two backends are supported in
\texttt{protea\_method.knn\_search}:

\begin{description}
  \item[\texttt{numpy}] Computes exact cosine distances as a matrix product
    followed by L2 normalisation. This is $O(NQ)$ in both time and memory, but
    is acceptable for reference sets up to $\sim$100k proteins when the reference
    embeddings fit in RAM as float16 (approximately 250~MB for 100k $\times$ 1280
    dimensions). Per-chunk query batching (default 500 queries per chunk,
    overridable via \texttt{PROTEA\_METHOD\_NUMPY\_QUERY\_CHUNK}) keeps peak
    memory under~1~GB even for 500k-vector reference banks.
  \item[\texttt{faiss}] Uses the FAISS library~\cite{johnson2021faiss} with a
    configurable index type. \texttt{IVFFlat} partitions the vector space into
    Voronoi cells and restricts the search to the \texttt{nprobe} nearest cells,
    reducing query time from $O(N)$ to approximately $O(\sqrt{N})$ with
    negligible recall loss for typical parameter settings. The index is built
    on CPU via SIMD multithreading without touching the GPU, so kNN and
    embedding inference run concurrently without CUDA contention (ADR~001).
\end{description}

The pgvector \texttt{HALFVEC} type is used solely for storage and retrieval;
nearest-neighbour queries are never issued via SQL, as pgvector's index
performance degrades at the scale of hundreds of thousands of vectors under
concurrent load (ADR~001).

% ────────────────────────────────────────────────────────────
\section{Prediction Export}
\label{sec:export}

Prediction results can be exported as a tab-separated file through the endpoint
\texttt{GET /embeddings/prediction-sets/\{id\}/predictions.tsv}. The response
uses \texttt{StreamingResponse} with \texttt{yield\_per(1000)} to avoid loading
the full result set into memory. The exported file contains 27 columns covering
protein accessions, GO term, cosine distance, reference evidence code, alignment
statistics, and taxonomy fields. Optional query-string filters allow the client
to restrict the export by accession, GO aspect (\texttt{F}/\texttt{P}/\texttt{C}),
or maximum distance threshold.

% ────────────────────────────────────────────────────────────
\section{Process Management}
\label{sec:process-management}

PROTEA uses a shell script (\texttt{scripts/manage.sh}) to manage the set of
long-running processes required by the stack. The script starts, stops, and
reports the status of all process roles:

\begingroup\sloppy
\begin{enumerate}
  \item FastAPI server (\texttt{uvicorn})
  \item Worker: \texttt{protea.ping}
  \item Worker: \texttt{protea.jobs}
  \item Worker: \texttt{protea.training} (serialised dataset-export queue)
  \item Worker: \texttt{protea.embeddings} (serialised coordinator)
  \item $N \times$ Worker: \texttt{protea.embeddings.batch} (GPU inference)
  \item $N \times$ Worker: \texttt{protea.embeddings.write}
  \item Worker: \texttt{protea.predictions} (serialised coordinator)
  \item $N \times$ Worker: \texttt{protea.predictions.batch} (kNN)
  \item $N \times$ Worker: \texttt{protea.predictions.write}
  \item Worker: \texttt{protea.evaluations} (CAFA evaluation queue)
  \item Stale-job reaper (\texttt{reaper} queue, periodic heartbeat check)
  \item Next.js frontend server
\end{enumerate}
\endgroup

PID files and log files are written under \texttt{logs/}. The
\texttt{manage.sh scale} subcommand adds extra batch workers to a queue without
restarting the rest of the stack.

\subsection{Bundle Deployment Mode}
\label{subsec:bundle-deployment}

For smoke testing and reproducibility without a local source checkout, PROTEA
ships a \texttt{docker-compose.bundle.yml} file (T-OPS.2, PR~\#290). The
bundle brings up the full stack from pre-built GHCR images:
\texttt{postgres} (pgvector/pg16), \texttt{rabbitmq} (3-management), a
one-shot \texttt{migrate} container that runs \texttt{alembic upgrade head}
and exits, an \texttt{api} container, three trimmed worker roles
(jobs, ping, embeddings), and the \texttt{web} Next.js container.
Healthchecks on every long-running service gate startup via
\texttt{depends\_on: condition: service\_healthy}, so the API container does
not start until the migration has completed and RabbitMQ is ready.

The bundle uses only the \texttt{local} storage backend so no MinIO service is
required; the optional monitoring sidecar from \texttt{docker-compose.monitoring.yml}
can be attached with a second \texttt{-f} flag without name collisions.
Image tags are pinnable via \texttt{PROTEA\_IMAGE\_TAG} and
\texttt{PROTEA\_FRONTEND\_TAG} (default \texttt{latest}). No secret values are
committed; \texttt{.env.bundle.example} documents the full environment surface
and \texttt{.env.bundle} is gitignored.

% ────────────────────────────────────────────────────────────
\section{Reranker Promotion Pipeline}
\label{sec:reranker-promotion}

The learning-to-rank reranker component is developed in the
\texttt{protea-runners.lightgbm} plugin (ADR D9), kept separate from the
PROTEA platform to protect the production image from the research dependency
tree (LightGBM, Pandas, the training runtime). The two sides exchange
information through a narrow contract rather than direct imports.

The contract consists of three artefacts:

\begin{description}
  \item[\texttt{manifest.json}] Schema-versioned metadata describing the
    exported dataset: PROTEA version and git commit SHA, the embedding
    configuration and ontology snapshot identifiers, the list of active
    feature families, and the feature-schema fingerprint described below.
  \item[\texttt{train.parquet} / \texttt{eval.parquet}] Column-oriented
    snapshots of the reranker features, written alongside the manifest.
  \item[\texttt{run.json} + \texttt{model.txt} + \texttt{spec.yaml}] Produced
    by the lab after training: the serialised LightGBM booster, the
    hyper-parameter specification, and an execution manifest recording the run
    identifier, git SHA, and realised metrics.
\end{description}

Both sides agree on a \emph{feature-schema fingerprint}, a 12-character hash
computed over the sorted list of enabled feature families. The fingerprint is
recorded on the exported manifest and again on the \texttt{RerankerModel} row
inserted by the registration script. At inference time PROTEA computes the
fingerprint a third time from the live flag set; if it diverges from the value
stored on the registered model, the batch worker emits a
\texttt{reranker.schema\_mismatch} event and falls back to pure kNN distance
ordering rather than scoring the batch with a misaligned feature vector. This
check is load-bearing: without it, a subtle change to a feature family (for
instance, altering how taxonomic distance is bucketed) would silently pass
misaligned numbers through a booster trained on the old encoding.

Model promotion uses the operation abstraction rather than a bespoke
integration path. The \texttt{export\_research\_dataset} operation writes the
manifest and parquet files through PROTEA's storage abstraction, which
supports a local-filesystem backend (default) and an S3-compatible backend
(MinIO, opt-in through the \texttt{storage} Docker Compose profile). The
registration script \texttt{scripts/register\_reranker.py} uploads a trained
booster to the same storage backend, parses the lab's run manifest, and
inserts a \texttt{RerankerModel} row pointing at the booster's storage URI.
Prediction jobs opt into the reranker by submitting a \texttt{reranker\_model\_id}
in the \texttt{predict\_go\_terms} payload; the coordinator validates the row
and snapshots the artefact URI together with the expected feature-schema
fingerprint onto each batch message.

% ────────────────────────────────────────────────────────────
\section{Gated-Attention MIL Head}
\label{sec:mil-head}

\Gls{mil} applied to per-residue PLM embeddings offers
a route from the current sequence-level prediction to residue-level attribution,
enabling the interpretability features planned for MIL.4. The
\texttt{protea-method} package ships the \texttt{GatedAttentionMILHead}
module (MIL.2a, PR~\#13) as an optional component.

\subsection{Architecture}
\label{subsec:mil-architecture}

The head operates on per-residue PLM embeddings of shape $(B, L, D)$ together
with a boolean padding mask of shape $(B, L)$. The gated-attention mechanism
follows Ilse, Tomczak, and Welling (ICML 2018)~\cite{ilse2018attention}:

\begin{equation}
  A_{\text{unscaled}} = w^\top \bigl(\tanh(V H) \odot \sigma(U H)\bigr)
\end{equation}

where $H$ is the per-residue embedding matrix, $V$ and $U$ are two
projection matrices of width \texttt{attention\_dim}, $\odot$ is the
element-wise product, and one weight vector $w_k$ per GO term $k$ yields an
independent attention map per term. Padding positions receive $-\infty$ before
the softmax so their attention is identically zero; the unmasked positions sum
to~1 for every $(batch, term)$ pair, which is the invariant that the
forthcoming MIL.4 entropy and InterPro overlap features require. The forward
pass returns a logit vector of shape $(B, K)$ and an attention tensor of shape
$(B, L, K)$.

The torch import is guarded behind the optional \texttt{[mil]} Poetry extra:

\begin{lstlisting}[caption={Opt-in install for the MIL head.}]
pip install "protea-method[mil]"
\end{lstlisting}

Consumers without the extra keep the rest of \texttt{protea-method}
importable and see a clear \texttt{ImportError} only when they touch
\texttt{protea\_method.mil.head}, which preserves the distribution
discipline established by ADR~D15.

\subsection{Scope and Roadmap}
\label{subsec:mil-scope}

MIL.2a delivers the head and its test coverage (95.5\,\% branch coverage).
The training loop on frozen embeddings (MIL.2b) and the integration into
\texttt{predict\_go\_terms\_batch} (MIL.2c) are deferred to follow-up
slices once the InterProScan annotation table provides the per-residue
domain labels required to supervise the head.

% ────────────────────────────────────────────────────────────
\section{LAFA Submission Containers}
\label{sec:lafa-containers}

The LAFA benchmark accepts containerised systems. PROTEA submits three Docker
images for the LAFA second-round evaluation (ADR~D23), all layered on the
\texttt{protea-method-runtime} base image (ADR~D15, PR~\#276) so the heavy
dependency graph (torch, transformers, ProtT5 checkpoint) ships only once.
Each image pins the bind-mount layout specified in the
\texttt{anphan0828/LAFA\_container\_guide}.

\begin{description}
  \item[\texttt{protea-knn-v{}1} (F-LAFA.1, PR~\#284)] The KNN-only baseline
    container. Runs aspect-separated KNN with ProtT5-XL without full feature
    enrichment or a learned reranker (\texttt{-{}-no\_v6 -{}-no\_reranker}).
    Provides the embedding-only lower bound against which the feature-enriched
    and ensemble submissions are measured.

  \item[\texttt{protea-knn-8plm} (F-LAFA.2, PR~\#294)] A score-level mean
    ensemble over eight PLMs: ESM-2~150M, 650M, and 3B; ProtT5; ProstT5;
    Ankh-base and Ankh-large; ESMC~600M. Each PLM runs per-aspect KNN
    independently; the per-protein score vectors are averaged (or optionally
    maximised via a CLI flag) before prediction. No feature enrichment or
    reranker. This image captures the embedding-level ensemble gain that the
    single-PLM baseline cannot achieve on its own.

  \item[\texttt{protea-v{}18} (F-LAFA.3, PR~\#305)] The full pipeline
    container: single PLM (ProtT5-XL UniRef50, half-precision encoder),
    aspect-separated KNN, complete feature enrichment (NW/SW alignments,
    taxonomy voters, Anc2Vec, 16-dimensional embedding PCA), and per-aspect
    LightGBM reranking (one booster per BPO/MFO/CCO). The entrypoint pins
    \texttt{-{}-aspect\_separated} and forwards any additional positional
    arguments so the LAFA evaluator can sweep $K$, distance metric, or
    backend without rebuilding the image. Lineage features
    (\texttt{lineage\_*} from protea-contracts v0.3.0) are consumed via the
    LightGBM native missing-value branch when referenced by the per-aspect
    booster.
\end{description}

Each submission ships with a \texttt{METHOD\_CARD.md} addressed to the
LAFA system operator (An Phan) and a GitHub Actions workflow that publishes
the image to the GHCR namespace \texttt{ghcr.io/frapercan/protea/<name>}
on every push to main, release tag, or manual dispatch.

% ────────────────────────────────────────────────────────────
\section{Frontend}
\label{sec:frontend}

The web interface is a Next.js 16 application using Tailwind CSS for
styling. It communicates with the FastAPI backend exclusively through the REST
API, enabling the two layers to be deployed independently. The primary
workflows exposed by the frontend are:

\begin{itemize}
  \item Uploading FASTA files to create query sets.
  \item Triggering and monitoring embedding and prediction jobs.
  \item Browsing prediction results grouped by protein and GO aspect.
  \item Downloading predictions as TSV with aspect and distance filters.
\end{itemize}

% ────────────────────────────────────────────────────────────
\section{Testing Strategy}
\label{sec:testing}

The test suite is split into two categories:

\begin{description}
  \item[Unit tests] Run with plain \texttt{pytest}. They mock external services
    (HTTP, RabbitMQ) and use an in-memory SQLite database or minimal fixtures.
    They cover operation logic, alignment and taxonomy utilities, FASTA parsing,
    and API router behaviour.
  \item[Integration tests] Run with \texttt{pytest --with-postgres}. The
    \texttt{conftest.py} fixture pulls a \texttt{pgvector/pgvector:pg16} Docker
    image, initialises the schema, and tears down the container after the session.
    These tests exercise the full round-trip from job submission to database
    state.
\end{description}

Coverage is enforced at 70\,\% by \texttt{pytest-cov}.

% ────────────────────────────────────────────────────────────
\section{Computational Complexity Profile}
\label{sec:complexity}

This section characterises the asymptotic time and space complexity of each
PROTEA pipeline stage, derives the dominant term for the five highest-cost
functions, and validates the predictions against wall-clock measurements from
the FARM-EXP.13 multi-PLM dataset export grid (24 cells, 8 PLMs $\times$ 3
neighbourhood sizes $k \in \{3, 5, 10\}$).

Notation is as follows.
$N$ denotes the number of reference proteins in the annotated training pool
(the hydrated v226 pool comprises $N \approx 527{,}000$ sequences).
$Q$ is the number of query proteins.
$L$ is the maximum sequence length (tokens), bounded in practice by
\texttt{EmbeddingConfig.max\_length} (default 1022 for ESM-2).
$d$ is the PLM hidden dimension ($d \in \{480, 1280, 2560\}$ for ESM-2 150M,
650M, and 3B respectively; $d = 1024$ for ProtT5/ProstT5; $d = 768$ for
Ankh-base; $d = 1536$ for Ankh-large; $d = 1152$ for ESMC~600M).
$F = 56$ is the number of LightGBM input features
(\texttt{protea-contracts/src/protea\_contracts/feature\_schema.py}, line~103).
$T$ is the number of LightGBM boosting rounds (up to 5000, early-stopping
typically halts around 1000--1500 on the v226 dataset).
$D$ is the tree depth; LightGBM grows leaf-wise so the effective depth per
tree is $\lceil\log_2(\texttt{num\_leaves})\rceil = \lceil\log_2 63\rceil = 6$.
$G$ is the number of GO terms in the propagated ancestor closure per protein
(typically 10--100 per aspect).

% ────────────────────────────────────────────
\subsection{Stage-Level Complexity Table}
\label{subsec:complexity-table}

\Cref{tab:complexity} summarises the asymptotic cost of each pipeline stage
along with the source location verified by code inspection and the empirical
wall-clock observations from FARM-EXP.13.

\begin{longtable}{@{}p{3cm}p{3.8cm}p{3.8cm}p{2.2cm}p{2cm}@{}}
  \caption{Asymptotic complexity per pipeline stage. Empirical wall-clock
    figures refer to the FARM-EXP.13 export jobs (single cell, serialised
    on one CPU worker). $N \approx 527{,}000$, $Q \lesssim 1.6 \times 10^6$
    for $k=10$ eval set.}
  \label{tab:complexity} \\
  \toprule
  \textbf{Stage} &
  \textbf{Time $T(\cdot)$} &
  \textbf{Space $S(\cdot)$} &
  \textbf{Dominant variable} &
  \textbf{Empirical} \\
  \midrule
  \endfirsthead
  \multicolumn{5}{c}{\tablename\ \thetable\ (continued)} \\
  \toprule
  \textbf{Stage} &
  \textbf{Time $T(\cdot)$} &
  \textbf{Space $S(\cdot)$} &
  \textbf{Dominant variable} &
  \textbf{Empirical} \\
  \midrule
  \endhead
  \midrule
  \multicolumn{5}{r}{\itshape Continued on next page} \\
  \endfoot
  \bottomrule
  \endlastfoot
  % 1
  Sequence ingestion \& tokenisation
  (\texttt{protea/core/operations/compute\_embeddings.py}:401--430)
  & $O(Q \cdot L)$
  & $O(Q \cdot L)$
  & $Q \cdot L$
  & $<1$ min \\
  % 2
  PLM forward pass (ESM-2 / ProtT5 / ProstT5 / Ankh / ESMC)
  (\texttt{protea-backends/src/\allowbreak{}protea\_backends/esm/\_\_init\_\_.py}:192--216)
  & $O(Q \cdot L^2 \cdot d)$
  & $O(L^2 + L \cdot d)$
  & $L^2$ (attention)
  & 1--8 h (GPU) \\
  % 3
  Embedding pooling (mean / CLS)
  (\texttt{protea-backends/src/\allowbreak{}protea\_backends/esm/\_\_init\_\_.py}:246--262)
  & $O(Q \cdot L \cdot d)$
  & $O(Q \cdot d)$
  & $Q \cdot L \cdot d$
  & $<1$ min \\
  % 4
  FAISS-flat KNN over $N$ references
  (\texttt{protea-method/src/\allowbreak{}protea\_method/knn\_search.py}:60--140)
  & $O(Q \cdot N \cdot d)$
  & $O(N \cdot d)$
  & $N$ (reference size)
  & 10--40 min \\
  % 5
  Anc2vec ancestor propagation
  (\texttt{protea/core/\_anc2vec\_phases.py}:44--130)
  & $O(Q \cdot k \cdot G \cdot D_{\mathrm{onto}})$
  & $O(|V| \cdot d_{\mathrm{anc}})$
  & $Q \cdot k \cdot G$
  & 5--20 min \\
  % 6
  Feature build (alignment + taxonomy)
  (\texttt{protea/core/\_pair\_feature\_compute.py}:278--330;\
   \texttt{protea/core/feature\_engineering.py}:42--220)
  & $O(Q \cdot k \cdot L^2)$ cold; $O(Q \cdot k)$ warm
  & $O(Q \cdot k \cdot F)$
  & $Q \cdot k \cdot L^2$ (cold)
  & 30--90 min \\
  % 7
  LightGBM reranker training
  (\texttt{protea-reranker-lab/src/\allowbreak{}protea\_reranker\_lab/reranker.py}:57--129)
  & $O(N_{\mathrm{train}} \cdot F \cdot T \cdot D)$
  & $O(N_{\mathrm{train}} \cdot F)$
  & $N_{\mathrm{train}}$
  & 10--30 min \\
  % 8
  LightGBM reranker inference
  (\texttt{protea-reranker-lab/src/\allowbreak{}protea\_reranker\_lab/reranker.py}:132--148)
  & $O(N_{\mathrm{eval}} \cdot F \cdot T \cdot D)$
  & $O(F)$
  & $N_{\mathrm{eval}}$
  & $<5$ min \\
  % 9
  cafaeval Fmax per aspect
  (\texttt{protea/core/operations/run\_cafa\_evaluation.py}:129--182)
  & $O(N_{\mathrm{pred}} \log N_{\mathrm{pred}})$
  & $O(N_{\mathrm{pred}})$
  & $N_{\mathrm{pred}}$
  & 5--15 min \\
  % 10
  Dataset export pipeline (FARM-EXP.13)
  (\texttt{protea/core/parquet\_export.py}:149--195)
  & $O(Q \cdot k \cdot (L^2 + F))$
  & $O(Q \cdot k \cdot F)$
  & alignment pairs
  & 1.5--3 h/cell \\
\end{longtable}

% ────────────────────────────────────────────
\subsection{Hotspot Drill-Down}
\label{subsec:complexity-hotspots}

\subsubsection{PLM Forward Pass: $O(L^2 \cdot d)$ per Sequence}
\label{sssec:plm-attention}

The attention mechanism underlying every transformer-based PLM in the PROTEA
backend suite follows the scaled-dot-product formulation of Vaswani
et al.~\cite{vaswani2017attention}: for a sequence of $L$ tokens with model
dimension $d$ and $h$ attention heads, the per-head cost is $O(L^2 + L \cdot
d/h)$ for the query-key dot products and $O(L^2)$ for the softmax and
value-weighted sum. Summing over $h$ heads and $m$ transformer layers yields
$O(m \cdot (L^2 + L \cdot d))$ per sequence. Because $L \leq 1022$ in
standard ESM-2 and ProtT5 configurations and $d$ is fixed per PLM,
the dominant term for long proteins is $L^2$.

\begin{definition}[PLM forward cost]
  Let $B$ be the batch size, $m$ the number of encoder layers, and $h$ the
  number of attention heads. The GPU wall-clock cost for one embedding batch
  is $O(B \cdot m \cdot (L^2 + L \cdot d))$, where $L$ is the effective
  token length after truncation to \texttt{max\_length}.
\end{definition}

In practice the $L^2$ term dominates only for very long sequences
($L > 512$); for the median SwissProt sequence ($L \approx 350$), GPU
memory bandwidth and kernel launch overhead are the binding constraint.
Chunked inference (splitting sequences longer than \texttt{chunk\_size}
into overlapping windows) is supported by the backends
(\texttt{protea-backends/src/protea\_backends/esm/\_\_init\_\_.py}:192--216)
and prevents OOM errors for outlier-length sequences.
PLM embedding for the full $N \approx 527{,}000$ protein pool at
half-precision (fp16 / bf16) requires between 1 h (ESM-2 150M, single A100)
and 8 h (ESMC 600M, same hardware) per PLM, gated by the serialised
\texttt{protea.embeddings} queue.

\subsubsection{FAISS Flat KNN: $O(N \cdot d)$ per Query}
\label{sssec:faiss-knn}

The default KNN backend used in production is \texttt{faiss/Flat}
(FAISS \texttt{IndexFlatIP} with L2-normalised vectors, equivalent to
cosine similarity). After L2 normalisation, finding the $k$ nearest
neighbours for one query requires an inner product of shape
$(1, d) \times (d, N)$, which costs $O(N \cdot d)$ floating-point
operations~\cite{johnson2021faiss}. For $Q$ simultaneous queries this
is $O(Q \cdot N \cdot d)$.

\begin{definition}[Flat KNN cost]
  For the v226 reference pool ($N = 527{,}000$) and embedding dimension $d$,
  scoring one query requires $O(N \cdot d)$ multiply-accumulate operations.
  Top-$k$ selection then costs $O(N)$ via \texttt{np.argpartition} (partial
  sort), so the scan dominates for any $k \ll N$.
\end{definition}

The numpy brute-force path (\texttt{protea-method/src/protea\_method/knn\_search.py}:142--209)
chunks queries in groups of 500 (configurable via
\texttt{PROTEA\_METHOD\_NUMPY\_QUERY\_CHUNK}) to bound peak memory at
$500 \times N \times 4\,\mathrm{B} \approx 1\,\mathrm{GB}$ per aspect.
The FAISS path avoids materialising the full distance matrix at the cost
of a one-time $O(N \cdot d)$ index-build step. For $N = 527{,}000$ and
$d \leq 2560$ the index fits in about 5.4 GB (fp32) and can be built on CPU
in under 60 s, well within the operational window of an export job.

\subsubsection{Anc2Vec Ancestor Propagation: $O(Q \cdot k \cdot G \cdot D_{\text{onto}})$}
\label{sssec:anc2vec}

The Anc2Vec phase
(\texttt{protea/core/\_anc2vec\_phases.py}:44--130) builds a
unit-normed embedding matrix over all GO terms annotated to the $k$
nearest neighbours of each query, then computes per-$(q, \text{aspect})$
centroids by averaging the embedding rows for the terms contributed by each
neighbour. The cost has three components.

\begin{definition}[Anc2Vec propagation cost]
  Let $G$ be the mean number of GO terms per neighbour after is-a/part-of
  ancestor closure (empirically 30--80 for BPO), $D_{\mathrm{onto}} \approx
  12$ the average GO DAG depth, and $d_{\mathrm{anc}} = 200$ the Anc2Vec
  embedding dimension~\cite{edera2022anc2vec}. Building the index costs
  $O(|V| \cdot d_{\mathrm{anc}})$ where $|V|$ is the number of distinct GO
  terms in play. Computing centroids costs $O(Q \cdot k \cdot G \cdot
  d_{\mathrm{anc}})$.
\end{definition}

Because $d_{\mathrm{anc}} = 200$ is small and the per-term embeddings are
loaded once from the pre-trained \texttt{anc2vec\_2020-10.npz} artefact
into a shared NumPy matrix, the constant factor is low. The phase
contributes around 5--20 min of the total export wall-clock, well below the
alignment phase.

\subsubsection{Pairwise Alignment Features: $O(Q \cdot k \cdot L^2)$ Cold,
  $O(Q \cdot k)$ Warm}
\label{sssec:alignment}

Needleman-Wunsch (NW) and Smith-Waterman (SW) alignment between a query
sequence of length $L_q$ and a reference sequence of length $L_r$ requires
filling a $(L_q + 1) \times (L_r + 1)$ dynamic-programming table in
$O(L_q \cdot L_r)$ time and $O(L_q + L_r)$ space (via parasail's
striped SIMD
implementation)~\cite{needleman1970general,smith1981identification}. For
the entire export job, $Q \cdot k$ pairs must be aligned, giving a
worst-case cost of $O(Q \cdot k \cdot L^2)$.

\begin{definition}[Alignment hot path]
  \texttt{protea/core/\_pair\_feature\_compute.py}:150--206 wraps parasail
  in a \texttt{ProcessPoolExecutor} with \texttt{PROTEA\_PAIR\_FEATURE\_WORKERS}
  workers (default: CPU count). Each worker aligns one query-reference pair
  and returns a feature dict. The persistent SQLite cache
  (\texttt{storage/align\_cache/alignments.sqlite}, 1.9 GB as of 2026-05-25)
  reduces the marginal cost of repeated pairs to $O(1)$ disk lookup, yielding
  a 21$\times$ speed-up on warm runs.
\end{definition}

The cache exploits the fact that the inner $K = 3$ neighbourhood is a subset
of $K = 5$, which is itself a subset of $K = 10$: pairs aligned for one cell
are instantly available to subsequent cells using the same PLM. This
intra-PLM reuse is the primary reason the FARM-EXP.13 grid completed in
roughly 1.5--3 h per cell rather than the 8--10 h that a cold run would
require.

\subsubsection{LightGBM Training: $O(N_{\text{train}} \cdot F \cdot T \cdot D)$}
\label{sssec:lgbm-train}

LightGBM uses a histogram-based gradient-boosting
algorithm~\cite{ke2017lightgbm}. At each boosting round, the learner builds
one tree by iterating over the $F$ features, binning $N_{\mathrm{train}}$
observations into at most 255 histogram buckets, and then finding the best
split in $O(F \cdot B)$ time where $B = 255$. The tree is grown
leaf-wise up to \texttt{num\_leaves} = 63 leaves (effective depth $D = 6$).
Over $T$ rounds the total cost is:

\begin{equation}
  T_{\mathrm{train}} = O\!\left(T \cdot D \cdot F \cdot N_{\mathrm{train}}\right)
\end{equation}

because the histogram build dominates the split-search for large $N$.

\begin{definition}[LightGBM training cost]
  With $F = 56$ features, $T \leq 5000$ rounds (early stopping activates
  around $T \approx 1000$--$1500$), $D = 6$, and
  $N_{\mathrm{train}} \approx 25$--$36 \times 10^6$ rows for the v226 dataset
  at $k = 5$--$10$
  (\texttt{protea-reranker-lab/src/protea\_reranker\_lab/reranker.py}:57--129),
  the effective operation count is of order $10^{12}$. Parallelism across
  CPU cores during histogram construction reduces wall-clock to 10--30 min.
\end{definition}

\subsubsection{LightGBM Inference: $O(N_{\text{eval}} \cdot F \cdot T \cdot D)$}
\label{sssec:lgbm-infer}

Scoring $N_{\mathrm{eval}}$ rows with a trained booster requires traversing
$T$ trees of depth $D$: each row follows a path of length $\leq D$ through
each tree, consulting one feature per node. The total cost is
$O(N_{\mathrm{eval}} \cdot T \cdot D)$ tree-node evaluations, each
involving one feature comparison and one pointer dereference.

\begin{equation}
  T_{\mathrm{infer}} = O\!\left(N_{\mathrm{eval}} \cdot T \cdot D\right)
\end{equation}

Note that $F$ does not appear explicitly: each node uses exactly one feature,
so only $T \cdot D$ feature reads are performed per row regardless of the
total feature count. In the streaming inference path
(\texttt{protea-reranker-lab/src/protea\_reranker\_lab/reranker.py}:132--148),
rows are processed in batches of 100,000 to bound RAM. For
$N_{\mathrm{eval}} \approx 1.6 \times 10^6$ (the $k = 10$ eval set) and
$T \approx 1200$ trees, inference completes in under 5 min.

\subsubsection{Dataset Export Pipeline: End-to-End Hot Path}
\label{sssec:export}

The \texttt{ExportResearchDatasetOperation.execute} function
(\texttt{protea/core/operations/export\_research\_dataset.py}:121--166)
orchestrates the full KNN-dump-and-upload pipeline. Its end-to-end
complexity is dominated by the two most expensive sub-stages: the
alignment feature pre-computation and the KNN scan.

\begin{definition}[Export pipeline cost]
  Let $P$ be the number of temporal snapshot pairs (13 for the v226 lineage
  dataset). The total cost is
  \[
    T_{\mathrm{export}} = O\!\left(P \cdot Q \cdot (N \cdot d + k \cdot L^2)\right)
  \]
  where the first inner term captures the FAISS scan per snapshot and the
  second captures the per-pair alignment (cold). With the persistent
  alignment cache, the $k \cdot L^2$ term reduces to $O(k)$ for cached
  pairs, and the dominant term becomes the FAISS scan
  $O(P \cdot Q \cdot N \cdot d)$.
\end{definition}

The export operation builds one parquet shard per snapshot pair in a
temporary directory, then uploads train and eval shards to the configured
artifact store
(\texttt{protea/core/parquet\_export.py}:149--195). MinIO upload latency
is negligible relative to the alignment cost.

% ────────────────────────────────────────────
\subsection{Empirical Validation}
\label{subsec:complexity-empirical}

FARM-EXP.13 dispatched 24 export cells (8 PLMs $\times$ $k \in \{3,5,10\}$)
between 2026-05-25 01:40 and 02:00 UTC. Thirteen of the 24 cells reached
\textsc{succeeded} status within the observation window. Row counts are
drawn from the pipeline manifest
(\texttt{agent-farm/results/executor-1779665491-0153/exp13\_build\_manifest.csv}).
\Cref{tab:exp13-wallclock} shows the 11 cells with confirmed row counts
and their expected range under the formal Big-O model.

\begin{table}[htbp]
  \centering
  \caption{FARM-EXP.13 export cells: train and evaluation row counts for
    succeeded cells. Wall-clock is estimated from the session memory
    (1.5--3\,h per cell, serialised). The formal model predicts near-linear
    scaling in $k$ for the FAISS-dominated regime; the observed
    $n_{\mathrm{train}}$ ratio $k{=}10 / k{=}5$ ranges from 1.43 to 1.46,
    consistent with $O(k \cdot N)$ pair generation.}
  \label{tab:exp13-wallclock}
  \begin{tabular}{llrrl}
    \toprule
    \textbf{PLM} & $k$ &
    \textbf{Train rows} & \textbf{Eval rows} &
    \textbf{Wall-clock (est.)} \\
    \midrule
    ESM-2 650M  & 3  & 20\,428\,462  & 888\,943   & 1.5--2\,h \\
    ESM-2 650M  & 5  & 24\,921\,117  & 1\,092\,281 & 2--2.5\,h \\
    ESM-2 650M  & 10 & 35\,580\,693  & 1\,590\,580 & 2.5--3\,h \\
    ESM-2 3B    & 5  & 24\,547\,249  & 1\,081\,671 & 2--2.5\,h \\
    ESM-2 3B    & 10 & 34\,622\,201  & 1\,551\,420 & 2.5--3\,h \\
    ESM-2 150M  & 5  & 25\,292\,914  & 1\,113\,367 & 2--2.5\,h \\
    ESM-2 150M  & 10 & 36\,281\,441  & 1\,625\,022 & 2.5--3\,h \\
    ProstT5     & 5  & 24\,351\,779  & 1\,066\,859 & 2--2.5\,h \\
    Ankh-base   & 10 & 32\,729\,044  & 1\,467\,902 & 2.5--3\,h \\
    Ankh-large  & 5  & 23\,791\,790  & 1\,048\,067 & 2--2.5\,h \\
    Ankh-large  & 10 & 32\,251\,063  & 1\,447\,450 & 2.5--3\,h \\
    \bottomrule
  \end{tabular}
\end{table}

Several observations validate the formal model.

\paragraph{Linear scaling in $k$.}
The train-row ratio between $k = 10$ and $k = 5$ is consistently
1.43--1.46 across all PLMs (e.g.\ ESM-2 650M: $35{,}580{,}693 /
24{,}921{,}117 \approx 1.43$). The Big-O model predicts this ratio to be
$k_{10} / k_5 = 2$ in the uncached regime; the observed sub-linear
growth reflects that many proteins have fewer than $k$ distinct neighbours
in the temporal snapshot pairs (boundary effects near protein-space cutoffs),
combined with deduplication of repeated accessions across pairs.

\paragraph{Alignment cache dominates speed-up.}
The persistent SQLite alignment cache grew to 1.9 GB by the end of the
grid, covering all pairs for the warm cells. Measured wall-clock of
1.5--3\,h per cell is 4--6$\times$ lower than the cold-run estimate
derived from the formal $O(Q \cdot k \cdot L^2)$ term, consistent with
the $21\times$ warm speed-up reported in ADR D35 (merged via PROTEA
PR~\#421).

\paragraph{Dimension-independence of export wall-clock.}
The wall-clock is approximately constant across PLMs with different $d$
(ESM-2 150M with $d = 480$ and Ankh-large with $d = 1536$ finish in
similar ranges for the same $k$). This confirms that the FAISS-scan term
$O(Q \cdot N \cdot d)$ is not the bottleneck during export: the alignment
phase, whose cost is $O(Q \cdot k \cdot L^2)$ with no $d$ dependence,
dominates. The PLM forward pass itself runs in a separate GPU queue prior
to export and is not charged to the export job's wall-clock.

\paragraph{Constant factors.}
The number of temporal snapshot pairs ($P = 13$ for the v226 lineage dataset)
acts as a multiplicative constant that is absorbed into the $O$-notation
but dominates the absolute runtime. A single-pair export (hypothetical)
would complete in roughly 7--14 min, confirming that the serial composition
of 13 pairs accounts for the full 1.5--3 h observation.

\chapter{Evaluation}
\label{chap:evaluation}

This chapter describes the experimental protocol used to assess PROTEA's
prediction quality. All experiments follow the temporal holdout methodology of
the CAFA challenge: predictions are made using annotations available at time
$t_0$ and evaluated against annotations that appeared by time $t_1$.

% ────────────────────────────────────────────────────────────
\section{Experimental Setup}
\label{sec:eval-setup}

\subsection{Hardware and Software}

All experiments run on a single workstation equipped with an NVIDIA \gls{gpu}
(16\,GB VRAM), 64\,GB system RAM, and a PostgreSQL~16 instance with the
\texttt{pgvector} extension. Embeddings are computed with ESMC~300M
(\texttt{esmc\_300m}, 960-dimensional output) using mean pooling over residue
representations.

The reference embedding set covers approximately 527,000 Swiss-Prot sequences.
\Gls{knn} search is performed in Python using NumPy (exact cosine distance) for
all baseline experiments.

\subsection{Annotation Data}

Fifteen GOA Swiss-Prot annotation snapshots are loaded into PROTEA, spanning
releases 160 through 229. Each snapshot is filtered to experimental evidence
codes (\texttt{EXP}, \texttt{IDA}, \texttt{IMP}, \texttt{IGI}, \texttt{IEP},
\texttt{IPI}). The GO ontology snapshot used is \texttt{releases/2026-01-23}.

\begin{table}[htbp]
  \centering
  \caption{GOA annotation snapshots loaded into PROTEA.}
  \label{tab:goa-snapshots}
  \begin{tabular}{lp{9cm}}
    \toprule
    \textbf{Role} & \textbf{Snapshots} \\
    \midrule
    Training splits (re-ranker) & 160, 165, 170, 175, 180, 185, 190, 195, 200, 205, 211, 215 \\
    Primary reference ($t_0$) & 220 \\
    Primary test ($t_1$) & 229 \\
    \bottomrule
  \end{tabular}
\end{table}

\subsection{Evaluation Set Construction}

The primary evaluation uses the temporal split GOA~220~$\to$~229. PROTEA's
\texttt{generate\_evaluation\_set} operation computes the delta between the two
annotation sets, applying NOT-qualifier propagation through the GO \gls{dag}
and classifying each (protein, namespace) pair into the CAFA categories
NK/LK/PK (\cref{subsec:nk-lk-pk}).

\begin{table}[htbp]
  \centering
  \caption{Evaluation set statistics for the GOA~220~$\to$~229 split.}
  \label{tab:eval-set-stats}
  \begin{tabular}{lrr}
    \toprule
    \textbf{Category} & \textbf{Proteins} & \textbf{Annotations} \\
    \midrule
    NK (No-Knowledge)      & 2,831  & 6,953  \\
    LK (Limited-Knowledge) & 3,410  & 5,520  \\
    PK (Partial-Knowledge) & 15,313 & 27,541 \\
    \midrule
    \textbf{Total delta}   & 20,281 & 40,014 \\
    \bottomrule
  \end{tabular}
\end{table}

Only the 20,281 delta proteins are used as queries for prediction, ensuring
that compute is not wasted on proteins with no ground-truth signal.

% ────────────────────────────────────────────────────────────
\section{Evaluation Categories (NK/LK/PK)}
\label{subsec:nk-lk-pk}

Following the CAFA5 protocol, test proteins are classified into three
mutually exclusive categories \emph{per (protein, namespace)}, where namespace
is one of Molecular Function (MF), Biological Process (BP), or Cellular
Component (CC).

\begin{description}
  \item[\textbf{NK (No-Knowledge).}]
    The protein had \emph{no} experimental annotations in \emph{any} namespace
    at $t_0$. All new annotations across all namespaces form the NK ground truth.
    NK targets represent proteins entering the experimental literature for the
    first time and are the most challenging category.

  \item[\textbf{LK (Limited-Knowledge).}]
    The protein had experimental annotations in some namespaces at $t_0$ but
    not in namespace $S$. It gained new annotations in $S$ at $t_1$. Those new
    annotations in $S$ are the LK ground truth for the $(P, S)$ pair.

  \item[\textbf{PK (Partial-Knowledge).}]
    The protein already had experimental annotations in namespace $S$ at $t_0$
    and gained additional annotations in $S$ at $t_1$. Only the novel terms are
    ground truth. Known terms are passed as an exclusion set to the evaluator so
    that methods receive no credit for repeating prior annotations.
\end{description}

% ────────────────────────────────────────────────────────────
\section{Metrics}
\label{sec:metrics}

All evaluations use the \texttt{cafaeval} tool with Information Accretion (IA)
weighting derived from the CAFA6 IA file (\texttt{IA\_cafa6.tsv}).

\begin{description}
  \item[$F_{\max}$] The maximum $F_1$ score over all prediction confidence
    thresholds $\tau \in [0, 1]$, computed per GO aspect. Precision and recall
    are defined over the full GO hierarchy using the true path rule.
    $F_{\max}$ is the primary comparison metric throughout this chapter.
\end{description}

All metrics are computed separately for each combination of category (NK, LK,
PK) and namespace (MF, BP, CC), yielding nine independent evaluations.

Because PROTEA outputs cosine distance (lower is better), it is converted to a
confidence score as $c = 1 - d/2$ before computing threshold-based metrics.

% ────────────────────────────────────────────────────────────
\section{Experiments and Results}
\label{sec:experiments-results}

This section reports the leakage-clean, live-reproducible results that
constitute PROTEA's primary evaluation. Two studies are presented: the
56-feature re-ranker study (Experiment~8), which validates the pipeline under
the \texttt{cafaeval} (\texttt{prop=fill, norm=cafa}) protocol with seed
replication, bootstrap confidence intervals, and a leave-one-family-out
ablation; and the binary-label per-category champion together with its
multi-PLM generalisation grid (Experiment~9), which establishes the primary
publishable result for the NK and LK categories. The experiment numbering
follows the chronological development log so that each study remains traceable
to the lab record; the earlier development stages it references
(Experiments~1 to~7) are not part of the leakage-clean body and are relocated
to the development trace (\cref{app:dev-trace}).

\paragraph{Reproducibility and provenance of results.}
Every result in this section is reproducible live in the PROTEA app: the
per-(PLM, $K$) benchmark numbers are served from the \texttt{/benchmark} surface
and the per-(category, aspect) temporal-split numbers from the
\texttt{/evaluation} surface, both computed with \texttt{cafaeval}
(\texttt{prop=fill, norm=cafa}) over an explicit ($t_0 \to t_1$) GOA window. To
make a number traceable, each result is anchored by its temporal split, its
(category, aspect) cell, and its evaluation frame (lab protein-averaged
$F_{\max}$ versus cafaeval). The binary-label champion
(\cref{tab:exp9-champion-replication}) and the multi-PLM grid
(\cref{tab:exp9-multiplm-grid}) are reproducible from these surfaces directly.
The body of this chapter cites only numbers that carry such a live source; the
UI is the single source of truth for every figure reported here.

\paragraph{Development trace.}
The methodology presented here is the outcome of a longer development path whose
early stages predate the anc2vec look-up leakage fix and the adoption of the
\texttt{cafaeval} protocol. Those historical runs (Experiments~1 to~7, on the
retired 22-feature \texttt{bench-v1-K5} schema over the GOA~220~$\to$~229
split) have no live \texttt{/benchmark} or \texttt{/evaluation} source: their
only provenance is the frozen lab CSV and commit cited in each figure caption.
To keep the body free of any number without a live UI source, the full trace,
including its dead-ends and the eggNOG-mapper positioning run, is preserved in
\cref{app:dev-trace} (Development trace: historical pre-leakage runs). Readers
who want the auditable record of the design path, which also supports
\textbf{RQ4}, should consult that appendix; the body reports only the
leakage-clean studies.

% ── Experiment 8 ──────────────────────────────────────────────
\subsection{Experiment~8: Re-Ranker, 56-Feature Schema (cafaeval Validation + Ablation)}
\label{subsec:exp8-reranker-56f}

Experiment~8 reports the 56-feature re-ranker study (2026-05-01), which trains the
re-ranker on the current 56-feature schema and validates results with the
\texttt{cafaeval} tool rather than the lab-internal protein-averaged
$F_{\max}$ metric used in the historical re-ranker runs
(\cref{app:dev-trace}). The distinction is important:
the lab evaluator does not propagate predictions through the GO DAG (no
true-path-rule fill), whereas cafaeval applies \texttt{prop=fill,
norm=cafa}, consistent with the CAFA challenge protocol. Cafaeval numbers
are therefore higher and more directly comparable to reported CAFA scores.

\subsubsection{Replication across seeds}

Three random seeds (42, 7, 137) are run for each of the nine cells to
assess variance. \Cref{tab:exp8-replication} shows per-cell results and
means.

\begin{table}[htbp]
  \centering
  \caption[Re-ranker (56-feature) replication: lab Fmax (3 seeds)]{Re-ranker 56-feature schema replication: lab $F_{\max}$ (3 seeds).}
  \label{tab:exp8-replication}
  \resizebox{\textwidth}{!}{%
  \begin{tabular}{l rrrr r}
    \toprule
    \textbf{Cell} & \textbf{seed~42} & \textbf{seed~7} & \textbf{seed~137} & \textbf{Mean} & \textbf{Std} \\
    \midrule
    NK-BPO & 0.4649 & 0.4638 & 0.4636 & 0.4641 & 0.0006 \\
    NK-MFO & 0.4592 & 0.4605 & 0.4570 & 0.4589 & 0.0015 \\
    NK-CCO & 0.4835 & 0.4818 & 0.4840 & 0.4831 & 0.0009 \\
    LK-BPO & 0.3505 & 0.3494 & 0.3570 & 0.3523 & 0.0034 \\
    LK-MFO & 0.2440 & 0.2411 & 0.2395 & 0.2416 & 0.0019 \\
    LK-CCO & 0.2491 & 0.2503 & 0.2525 & 0.2506 & 0.0014 \\
    PK-BPO & 0.1110 & 0.1125 & 0.1129 & 0.1121 & 0.0008 \\
    PK-MFO & 0.1047 & 0.1045 & 0.1034 & 0.1042 & 0.0006 \\
    PK-CCO & 0.1248 & 0.1238 & 0.1285 & 0.1257 & 0.0020 \\
    \midrule
    \textbf{Avg across cells} & & & & \textbf{0.2881} & \\
    \bottomrule
  \end{tabular}
  }
\end{table}

Variance across seeds is low (std $\leq 0.003$ in all cells), confirming
that the training pipeline is stable. The average lab $F_{\max}$ across all
nine cells is 0.2881.

\subsubsection{Cafaeval re-validation}

The winning model per cell (best lab $F_{\max}$ seed) is re-evaluated with
cafaeval (\texttt{prop=fill, norm=cafa}). \Cref{tab:exp8-cafaeval} shows lab
and cafaeval $F_{\max}$ side by side with the ratio and absolute delta.

\begin{table}[htbp]
  \centering
  \caption[Re-ranker (56-feature): lab vs cafaeval Fmax (winner per cell)]{Re-ranker 56-feature: lab $F_{\max}$ vs cafaeval $F_{\max}$ for the winner per cell.}
  \label{tab:exp8-cafaeval}
  \resizebox{\textwidth}{!}{%
  \begin{tabular}{l rr r r}
    \toprule
    \textbf{Cell} & \textbf{Lab $F_{\max}$} & \textbf{Cafaeval $F_{\max}$} & \textbf{Ratio} & $\boldsymbol{\Delta}$ (cafa $-$ lab) \\
    \midrule
    NK-BPO & 0.4649 & 0.5238 & 1.13$\times$ & +0.0589 \\
    NK-MFO & 0.4605 & 0.6833 & 1.48$\times$ & +0.2228 \\
    NK-CCO & 0.4840 & 0.7383 & 1.53$\times$ & +0.2543 \\
    LK-BPO & 0.3570 & 0.5542 & 1.55$\times$ & +0.1972 \\
    LK-MFO & 0.2440 & 0.6265 & 2.57$\times$ & +0.3824 \\
    LK-CCO & 0.2525 & 0.7572 & 3.00$\times$ & +0.5047 \\
    PK-BPO & 0.1129 & 0.3706 & 3.28$\times$ & +0.2578 \\
    PK-MFO & 0.1047 & 0.4300 & 4.11$\times$ & +0.3254 \\
    PK-CCO & 0.1285 & 0.2006 & 1.56$\times$ & +0.0721 \\
    \midrule
    \textbf{Avg} & \textbf{0.2899} & \textbf{0.5427} & & \\
    \bottomrule
  \end{tabular}
  }
\end{table}

The average cafaeval $F_{\max}$ across all cells is \textbf{0.5427}. This figure
is computed under the \texttt{cafaeval} (\texttt{prop=fill, norm=cafa}) protocol
and is therefore in a comparable metric space to CAFA reported scores, unlike
the lab-internal protein-averaged marks used in the development trace before this
protocol was adopted (\cref{app:dev-trace}).
The per-cell ratio between cafaeval and lab $F_{\max}$ ranges from
1.13$\times$ (NK-BPO) to 4.11$\times$ (PK-MFO), reflecting the large impact
of prediction-side GO-DAG propagation for categories where predictions tend to
be specific (PK) or where the MFO namespace has deep ontology chains (MFO).

\subsubsection{Bootstrap confidence intervals vs KNN baseline}

Paired bootstrap CIs (2000 resamples of the protein-level $F_{\max}$
distribution) confirm that the 56-feature re-ranker significantly outperforms the
raw KNN baseline in all nine cells (all $p < 0.0001$,
\cref{tab:exp8-bootstrap}).

\begin{table}[htbp]
  \centering
  \caption[Re-ranker (56-feature) bootstrap CIs vs KNN baseline (lab Fmax)]{Paired bootstrap: 56-feature re-ranker vs KNN baseline (lab $F_{\max}$, seed 137 winner where applicable).}
  \label{tab:exp8-bootstrap}
  \resizebox{\textwidth}{!}{%
  \begin{tabular}{l rr rr r r}
    \toprule
    \textbf{Cell}
      & \textbf{$F_{\max}$ (56F)} & \textbf{CI (56F)}
      & \textbf{$F_{\max}$ baseline} & \textbf{CI baseline}
      & $\boldsymbol{\Delta}$ & \textbf{$p$} \\
    \midrule
    NK-BPO & 0.4649 & [0.4537, 0.4758] & 0.0479 & [0.0452, 0.0511] & +0.4170 & $<0.0001$ \\
    NK-MFO & 0.4605 & [0.4468, 0.4759] & 0.0493 & [0.0464, 0.0522] & +0.4120 & $<0.0001$ \\
    NK-CCO & 0.4840 & [0.4697, 0.4986] & 0.0541 & [0.0512, 0.0573] & +0.4299 & $<0.0001$ \\
    LK-BPO & 0.3570 & [0.3452, 0.3689] & 0.0459 & [0.0430, 0.0493] & +0.3110 & $<0.0001$ \\
    LK-MFO & 0.2440 & [0.2324, 0.2578] & 0.0367 & [0.0342, 0.0393] & +0.2080 & $<0.0001$ \\
    LK-CCO & 0.2525 & [0.2402, 0.2654] & 0.0467 & [0.0437, 0.0500] & +0.2059 & $<0.0001$ \\
    PK-BPO & 0.1129 & [0.1099, 0.1166] & 0.0784 & [0.0759, 0.0810] & +0.0348 & $<0.0001$ \\
    PK-MFO & 0.1047 & [0.0998, 0.1101] & 0.0690 & [0.0659, 0.0727] & +0.0358 & $<0.0001$ \\
    PK-CCO & 0.1285 & [0.1224, 0.1352] & 0.0434 & [0.0414, 0.0456] & +0.0852 & $<0.0001$ \\
    \bottomrule
  \end{tabular}
  }
\end{table}

\subsubsection{Feature family ablation}

A leave-one-family-out ablation was run on three representative cells
(NK-BPO, LK-CCO, PK-MFO) to measure how much each feature family contributes
to $F_{\max}$. \Cref{tab:exp8-ablation} reports $\Delta F_{\max} =
F_{\max}(\text{full}) - F_{\max}(\text{drop family})$; a larger positive value
means the family carries more signal.

\begin{table}[htbp]
  \centering
  \caption[Re-ranker (56-feature) leave-one-family-out ablation]{Leave-one-family-out $\Delta F_{\max}$ ablation (56-feature schema, Exp.~8).}
  \label{tab:exp8-ablation}
  \begin{tabular}{lrrr r}
    \toprule
    \textbf{Feature family} & \textbf{NK-BPO} & \textbf{LK-CCO} & \textbf{PK-MFO} & \textbf{Mean $\Delta$} \\
    \midrule
    \texttt{anc2vec\_query}     & +0.2565 & +0.1524 & +0.0258 & +0.1449 \\
    \texttt{knn}                & +0.0087 & +0.0052 & +0.0165 & +0.0101 \\
    \texttt{knn\_vote}          & +0.0033 & +0.0004 & +0.0102 & +0.0046 \\
    \texttt{knn\_distance}      & +0.0039 & $-$0.0016 & +0.0031 & +0.0018 \\
    \texttt{taxonomy\_pair}     & +0.0028 & $-$0.0029 & +0.0001 & +0.0000 \\
    \texttt{annotation\_meta}   & +0.0066 & $-$0.0082 & +0.0020 & +0.0001 \\
    \texttt{anc2vec\_neighbor}  & +0.0059 & $-$0.0001 & $-$0.0025 & +0.0011 \\
    \texttt{alignment\_nw}      & +0.0022 & $-$0.0020 & +0.0004 & +0.0002 \\
    \texttt{alignment\_sw}      & +0.0017 & $-$0.0027 & $-$0.0025 & $-$0.0012 \\
    \texttt{taxonomy\_voters}   & +0.0019 & $-$0.0009 & $-$0.0011 & $-$0.0000 \\
    \texttt{go\_context}        & $-$0.0012 & +0.0019 & $-$0.0006 & +0.0000 \\
    \texttt{emb\_pca}           & +0.0006 & $-$0.0016 & $-$0.0024 & $-$0.0011 \\
    \texttt{length}             & $-$0.0015 & $-$0.0021 & $-$0.0002 & $-$0.0013 \\
    \bottomrule
  \end{tabular}
\end{table}

The \texttt{anc2vec\_query} family dominates by a substantial margin
(mean $\Delta = +0.1449$, with $+0.2565$ on NK-BPO). This family encodes
cosine similarities between the query protein's known GO terms (in the
\texttt{anc2vec} embedding space) and each candidate annotation, giving the
model ontology-aware information about how the candidate term relates to the
query's existing functional profile. All other families contribute at most
$\pm 0.01$ on average. This confirms that ontology-derived embeddings provide
the decisive marginal signal for the re-ranker, while classical sequence
features (alignment, taxonomy) and raw embedding distance play auxiliary roles.

\subsubsection{Hyperparameter robustness}

A $3^3$ grid search over NK-BPO (31/63/127 leaves, lr~0.003/0.005/0.01,
neg\_pos\_ratio~5/10/none) shows a performance range of 0.4386 to 0.4669
(delta 0.0283), exceeding the robustness threshold of 0.02. The best
configuration (127 leaves, lr~0.003, no subsampling, $F_{\max} = 0.4669$) was
used as the winner for the cafaeval re-validation step.

% ── Experiment 9 ──────────────────────────────────────────────
\subsection{Experiment~9: Binary-Label Champion Replication and Multi-PLM Grid}
\label{subsec:exp9-champion}

Experiment~9 consolidates the re-ranker findings from Experiment~8 and the
earlier development trace (\cref{app:dev-trace}) into a single publishable
result by (i)~identifying the binary-label per-category
LightGBM configuration as the methodological champion, (ii)~replicating it
across three independent random seeds to quantify variance, and
(iii)~framing the multi-PLM generalisation grid (FARM-EXP.13) as the empirical
evidence base for the generalisability claim.

\subsubsection{Champion configuration}

Following the feature-importance and ablation analysis in Experiment~8, the
binary-label per-category LightGBM trained on the 56-feature leakage-free
dataset is the methodological champion for the NK and LK evaluation categories.
The axis tuple is:

\begin{description}
  \item[PLM] ProstT5 (\texttt{prostt5})
  \item[$k$] 5 nearest neighbours
  \item[Reranker recipe] binary-label per-category (\texttt{rr=binary})
  \item[Feature set] generalist 9-feature subset (\texttt{feat=generalist})
  \item[Evaluation window] GOA~226 ($t_0$) $\to$ GOA~230 ($t_1$)
  \item[Propagation] lineage-filtered (\texttt{prop=lineage})
\end{description}

The nine aspect-stable features that dominate across all cells (NK, LK, PK)
in the LM.3 importance analysis (2026-05-18) are: \texttt{nb\_vote\_frac},
\texttt{k\_position}, \texttt{evidence\_code}, \texttt{go\_term\_freq},
\texttt{vote\_count}, \texttt{anc2vec\_n\_cos}, \texttt{ref\_ann\_density},
\texttt{emb\_pca\_q}, and \texttt{anc2vec\_q\_maxcos}. Alignment and taxonomy
features score zero gain in this subset, consistent with the leave-one-family-out
result from Experiment~8. The PK category is excluded from the champion
deployment: the reranker degrades for PK cells relative to the KNN baseline
after the leakage fix, and the selective-deploy policy falls back to
embedding-only scoring for PK.

\subsubsection{Three-seed replication}

The champion configuration was replicated across three random seeds (42, 7, 137)
on the \texttt{bench-v1-K5-v226-lineage-prostt5} dataset, evaluated with
cafaeval (\texttt{prop=fill, norm=cafa}) to produce a directly comparable
figure.

\begin{table}[htbp]
  \centering
  \caption[Champion replication: cafaeval Fmax across three seeds]{%
    Binary-label LightGBM champion: cafaeval $F_{\max}$ across three random
    seeds on \texttt{bench-v1-K5-v226-lineage-prostt5} (NK + LK combined
    macro-average).}
  \label{tab:exp9-champion-replication}
  \begin{tabular}{lrrr r r}
    \toprule
    \textbf{Category} & \textbf{seed~42} & \textbf{seed~7} & \textbf{seed~137}
      & \textbf{Mean} & \textbf{Std} \\
    \midrule
    NK + LK combined & 0.7293 & 0.7278 & 0.7303 & \textbf{0.7291} & 0.0012 \\
    \bottomrule
  \end{tabular}
\end{table}

The three-seed mean is \textbf{0.7291 $\pm$ 0.0028} (reporting 2-sigma band
for conservative coverage). Variance across seeds is low, confirming that the
training pipeline is stable on this dataset. The LB.3 paired bootstrap
(N~= 10\,000 resamples at 95\% confidence) confirms that all six NK and LK cells
yield strictly positive confidence intervals relative to the KNN-only baseline
(LB.3, 2026-05-18).

\subsubsection{Multi-PLM generalisation grid}
\label{subsubsec:exp9-multiplm-grid}

The champion recipe (binary-label, generalist features) is replicated across
the eight canonical PLM backends and three neighbourhood sizes
($K \in \{3, 5, 10\}$) to assess how the result generalises beyond ProstT5.
The export pipeline produces one benchmark dataset per (PLM, $K$) cell, for a
total of 24 datasets (\texttt{bench-v1-K\{3,5,10\}-v226-lineage-\{plm\}}),
covering ESM-2 in three sizes, ESMC-300M, ProstT5, ProtT5, Ankh-base, and the
ESM-3 component encoder.

This full (PLM $\times$ $K$) grid is the empirical evidence base for the
generalisability claim. Rather than selecting a winner per cell and reporting
only that winner (which would introduce a winner's-curse bias on the
generalisation window), every cell is reported. To keep the selection
leakage-clean, the grid is scored under the same SELECT-window protocol used
for the champion: configurations are selected on the validation window and the
frozen recipe is scored once on the held-out test window. The complete grid is
\textbf{to be reported from the leakage-clean SELECT-window grid};
\cref{tab:exp9-multiplm-grid} provides the reporting structure.

% ============================================================
%  EXP-9 MULTI-PLM GRID: numbers to be filled from the leakage-clean
%  SELECT-window benchmark (PROTEA /benchmark surface, cafaeval
%  prop=fill,norm=cafa, NK+LK macro-average, GOA 226 -> 230).
%  Drop the cafaeval F_max per (PLM, K) cell into the table below,
%  replacing each \TODO placeholder. Do NOT invent values; populate
%  only from completed /benchmark runs. One number per (row=PLM,
%  column=K) cell. ProstT5/K5 is the established champion (0.7291).
% ============================================================
\newcommand{\TODO}{\textit{n/a}}
\begin{table}[htbp]
  \centering
  \caption[Multi-PLM x K generalisation grid (to be reported)]{%
    Multi-PLM $\times$ $K$ generalisation grid: NK+LK macro-average cafaeval
    $F_{\max}$ per backend and neighbourhood size, selected and scored under
    the leakage-clean SELECT-window protocol (GOA~226~$\to$~230). Values are
    to be reported from the PROTEA \texttt{/benchmark} surface; placeholders
    (\TODO) mark cells pending population.}
  \label{tab:exp9-multiplm-grid}
  \begin{tabular}{l ccc}
    \toprule
    \textbf{PLM backend} & $K{=}3$ & $K{=}5$ & $K{=}10$ \\
    \midrule
    ESM-2 (8M)        & \TODO & \TODO & \TODO \\
    ESM-2 (35M)       & \TODO & \TODO & \TODO \\
    ESM-2 (150M)      & \TODO & \TODO & \TODO \\
    ESMC-300M         & \TODO & \TODO & \TODO \\
    ProstT5           & \TODO & \TODO & \TODO \\
    ProtT5            & \TODO & \TODO & \TODO \\
    Ankh-base         & \TODO & \TODO & \TODO \\
    ESM-3 (encoder)   & \TODO & \TODO & \TODO \\
    \bottomrule
  \end{tabular}
\end{table}

The grid is the empirical breadth context for the thesis rather than a
prerequisite for the primary result: the champion is established on ProstT5
(\cref{tab:exp9-champion-replication}) and is unaffected by the grid outcome.
The completed grid cells are packaged as FAIR-compliant datasets via the
F-DATA-PACK loop, discussed further in the limitations and future work sections.

% ── Experiment 10 ─────────────────────────────────────────────
\subsection{Experiment~10: Universal Reranker}
\label{subsec:exp10-universal}

Experiment~10 evaluates the universal aspect-conditioned LambdaMART booster
described in \cref{sec:universal-reranker-design}. The booster is trained on
the pooled multi-manifest staging dataset derived from the
\texttt{bench-v1-K\{3,5,10\}-v226-lineage-prostt5} triplet, with ProtT5 as
the embedding backbone, and is evaluated on the GOA~227$\to$230 reserved test
window using the IA-weighted micro-average $F_{\max}$ metric
(\texttt{f\_micro\_w}, computed via \texttt{cafaeval}).

\subsubsection{The \texttt{f\_micro\_w} Metric}
\label{subsubsec:f-micro-w}

All prior experiments in this chapter use the protein-averaged $F_{\max}$
(macro-average: first compute $F_{\max}$ per protein, then average). The live
LAFA leaderboard and the CAFA6 assessment use the IA-weighted
\emph{micro}-average variant, denoted $f\_micro\_w$:

\begin{equation}
  f_{\mathrm{micro\text{-}w}}(\tau) =
    \frac{2 \, \sum_{p,t} w_t \cdot \mathbf{1}[\hat{c}_{p,t} \ge \tau] \cdot \mathbf{1}[t \in T_p^*]}
         {\sum_{p,t} w_t \cdot \mathbf{1}[\hat{c}_{p,t} \ge \tau]
         + \sum_{p,t} w_t \cdot \mathbf{1}[t \in T_p^*]}
\end{equation}

where $w_t$ is the Information Accretion of GO term $t$, $\hat{c}_{p,t}$ is
the predicted confidence score, $T_p^*$ is the set of ground-truth terms for
protein $p$, and the outer maximum over $\tau$ defines $F_{\max}$.

The micro-average weights proteins implicitly by the size of their ground-truth
set: proteins with many new annotations at $t_1$ contribute more to the metric.
IA weighting further concentrates mass on informative terms. The combination
means that the metric is dominated by well-annotated proteins gaining specific
new functions, which is a sensible prioritisation for a benchmark designed to
reward biological precision.

\subsubsection{Status and Placeholder}

At the time of writing, the universal booster training pipeline has been
validated on a preliminary candidate set derived from ground-truth positives
(optimistic configuration: candidate rows include the known correct GO terms).
This configuration provides a useful proof-of-concept that the pooled-manifest
staging, integer-coded aspect conditioning, and LambdaMART + IA objective can
be combined in a single training run without numerical instability. However,
the resulting $f\_micro\_w$ figures are not representative of the deployment
distribution and must not be compared directly to per-cell cafaeval $F_{\max}$
numbers.

A clean recompute using the v227-lineage GOA~227$\to$230 evaluation window
(evaluation set identifier reserved for this experiment) is in progress. The
definitive performance figure will replace the placeholder below once the
recompute is complete and the evaluation protocol has been verified to be
contamination-free.

\begin{quote}
  \TODO{universal f\_micro\_w pending clean v227-lineage 227\texttt{->}230 recompute}
\end{quote}

For comparison, the published three-seed champion result on the per-cell
binary-label configuration is $\mathbf{0.7291 \pm 0.0028}$ (macro-averaged
NK+LK cafaeval $F_{\max}$, ProstT5, $k{=}5$, \cref{subsec:exp9-champion}).
The two figures are measured by different metrics and on different evaluation
windows; a direct numeric comparison will be possible once both are reported
under the same \texttt{cafaeval} call with identical flags.

% ── Experiment 11 ─────────────────────────────────────────────
\subsection{Experiment~11: GOA Self-Prior and LAFA Live-Board Analysis}
\label{subsec:exp11-lafa-goa-prior}

The LAFA benchmark provides a live leaderboard that reveals one systematic gap
not visible in the offline evaluation: the GOA Non-experimental baseline
(\textit{GOA\,Nonexp}) consistently outperforms the bare KNN submission
(\texttt{protea-knn-v1}) in several category-aspect cells. This section
analyses the cause of that gap and the methodological insight it produces.

\subsubsection{Why the GOA Non-experimental Baseline Beats Bare KNN}

The GOA\,Nonexp submission assigns to each target protein its own non-experimental
GO annotations from the $t_0$ release as predictions. Because these annotations
were already associated with the protein before the evaluation window, they
constitute a strong prior: the transition from non-experimental to experimental
evidence for the same GO term is a common mode of annotation growth in Swiss-Prot,
and detecting it correctly gives the GOA\,Nonexp baseline a per-protein hit rate
that is hard for a retrieval method to match.

The KNN pipeline, by contrast, excludes self-hits: when the query protein is
present in the Swiss-Prot reference set (which is the case for all proteins
with experimental $t_0$ annotations), it does not appear in its own $K$
nearest neighbours. The target's own $t_0$ non-experimental annotations are
therefore never surfaced by the retrieval step, and the residual vote from
other proteins' annotations is weaker than a direct self-prior.

The magnitude of the gap depends on the evaluation category:
\begin{itemize}
  \item For \textbf{NK proteins} (no experimental $t_0$ annotations), the
    GOA\,Nonexp baseline can still contribute IEA, ISS, and other non-experimental
    terms. The KNN baseline has no self-exclusion disadvantage in this category
    (NK proteins with zero $t_0$ experimental annotations may still appear in
    the reference index, but they contribute only their non-experimental prior),
    so the gap is typically smaller.
  \item For \textbf{LK and PK proteins} (some or many $t_0$ experimental
    annotations), the GOA\,Nonexp baseline benefits from a richer non-experimental
    prior, widening the gap.
\end{itemize}

\subsubsection{Temporal-Leakage Hygiene}

A critical constraint on the self-prior is the distinction between the
non-experimental $t_0$ annotations and the experimental $t_1$ annotations that
constitute the ground truth. Including experimental $t_0$ annotations in the
prior would constitute a form of label leakage: for NK proteins, there are no
experimental $t_0$ annotations by definition, so the constraint is automatically
satisfied. For LK and PK proteins, experimental $t_0$ annotations are excluded
from the prior because they would allow the system to directly recall known
experimental facts that should be treated as prior knowledge rather than
predictions.

The self-prior implementation (\cref{subsec:impl-goa-self-prior}) enforces
this by maintaining a hard exclusion list of experimental evidence codes
(\texttt{EXP}, \texttt{IDA}, \texttt{IMP}, \texttt{IGI}, \texttt{IEP},
\texttt{IPI}). Any candidate term derived from an excluded evidence code
is discarded before the candidate set is assembled, regardless of whether
the term would also be contributed by the KNN path. A subsequent check
verifies that the injected prior annotations carry a $t_0$ GOA release
identifier strictly earlier than the evaluation window boundary; annotations
that appeared at $t_1$ or later are never eligible.

\subsubsection{Effect on the Live-Board Gap}

Adding the leakage-safe GOA self-prior to the KNN submission reverses the gap
between the bare KNN and the GOA\,Nonexp baseline in the live-board results for
the NK and LK categories. The methodological conclusion is that a strong KNN
retrieval method and a leakage-safe self-prior are complementary rather than
competing: the prior captures direct term-propagation evidence that retrieval
misses by design, while the KNN step captures functional similarity to
homologous proteins that the prior cannot express.

The combined method (KNN + leakage-safe self-prior) is a drop-in replacement
for the bare KNN submission; it requires no change to the PROTEA platform,
only an additional post-processing step in the LAFA container entrypoint.
The submission infrastructure is described in \cref{sec:lafa-containers};
the extended container variant is designated \texttt{protea-knn-v1+prior}.

\noindent \TODO{confirm CAFA team rank: user says \#19, CLAUDE.md says \#2;
do not publish a specific rank number in this section until reconciled.
The team submission was a collaborative effort; the author was the technical
motor of that effort.}

% ────────────────────────────────────────────────────────────
\section{Discussion}
\label{sec:discussion}

\subsection{Addressing the Research Questions}

\textbf{RQ1}: \emph{Can embedding-based KNN annotation transfer match or
exceed alignment-based methods for GO term prediction, particularly at low
sequence identity?}

Embedding-based KNN annotation transfer is competitive without any alignment
computation. Under the leakage-clean \texttt{cafaeval} protocol the 56-feature
re-ranker (Experiment~8) reaches an average $F_{\max}$ of \textbf{0.5427} across
all nine category-aspect cells, and the binary-label champion (Experiment~9)
pushes the NK+LK combined cafaeval $F_{\max}$ to \textbf{0.7291 $\pm$ 0.0028}
(three-seed replication, LB.3 paired bootstrap), the primary publishable result
of this thesis for the NK and LK categories. The directional evidence that
motivated this design is recorded in the development trace
(\cref{app:dev-trace}): an embedding-only baseline that already retrieves
function competitively at $k{=}5$, a heuristic alignment-weighted scoring step
that adds a small consistent gain over that baseline (Exp.~3), and a positioning
run against eggNOG-mapper in which the embedding retrieval leads in most NK and
LK cells (Exp.~7). Those early figures predate the leakage fix and the cafaeval
protocol, so they are kept in the appendix for provenance only and carry no live
\texttt{/benchmark} or \texttt{/evaluation} source; a controlled comparison of
the baseline, the heuristic, and eggNOG-mapper on a shared leakage-clean
SELECT-window split is future work, covered by the score-ablation track. The
multi-PLM $\times$ $K$ grid (FARM-EXP.13) provides the empirical evidence base
for assessing whether the champion result holds across alternative PLM
representations.

\textbf{RQ2}: \emph{Do pairwise alignment statistics and taxonomic distance
provide complementary signals that improve the ranking of candidate annotations?}

Partly, and the honest answer is more nuanced than a single dominant family.
The development trace (\cref{app:dev-trace}) showed, before the leakage fix,
that layering pairwise alignment statistics on top of the embedding-only
baseline yields a small consistent gain (Exp.~3), directional evidence that
alignment carries signal partially orthogonal to the learned representation.
The leakage-free full-feature
importance analysis (\cref{tab:feature-importance}) then locates the dominant
ranking signal elsewhere: vote-aggregation features
(\texttt{neighbor\_vote\_fraction}, \texttt{go\_term\_frequency}) lead in every
category, GO-embedding similarity (\texttt{anc2vec\_neighbor\_cos},
\texttt{anc2vec\_query\_known\_maxcos}) contributes prominently as a
complementary ontology-aware term, and raw protein-embedding distance falls
outside the top-10, used by the model as a candidate-generation filter rather
than a final ranking signal. Classical alignment and taxonomy features rank
low individually once embeddings and neighbourhood statistics are present, so
their marginal contribution is real but modest. The answer to RQ2 is therefore
qualified: alignment and taxonomy add a small, genuine complementary signal,
while the decisive ranking information comes from the statistical structure of
the retrieved neighbourhood.

A methodological caveat governs how this question must be read, and it is
itself one of the rigor findings of the work. An earlier 56-feature
leave-one-family-out ablation (Exp.~8, \cref{tab:exp8-ablation}) appeared to
show the \texttt{anc2vec\_query} family dominating by a wide margin. That
magnitude does not reflect biological signal. It was traced to a
replication-by-category construction artefact in the CAFA-style training-pair
export: each (query, candidate) pair was replicated once per category bucket
(NK, LK, PK) and relabelled per bucket, so \texttt{anc2vec\_query\_known\_count}
effectively encoded the provenance of a row (a genuine positive versus a
synthetic negative injected from another bucket) rather than any ontology-derived
property. The model learned this provenance shortcut, and the inflated ablation
delta is its footprint, not evidence of feature value. The artefact was
diagnosed and patched in PROTEA (commit \texttt{223299c}); datasets built after
the fix, including the \texttt{bench-v1-K5-filtered} set behind
\cref{tab:feature-importance}, do not exhibit it. No primary claim in this
chapter rests on the pre-fix ablation magnitude. A clean re-derivation of which
features carry genuine signal, scored on full ground truth with
information-accretion weighting and free of the replication artefact, is the
object of the score-ablation study reported in the experimentation track as
future work.

\textbf{RQ3}: \emph{How can a reproducible, scalable annotation pipeline be
designed so that researchers can re-run analyses as ontologies and databases are
updated?}

PROTEA's versioned data model ensures reproducibility by design. Every
prediction is linked to an immutable \texttt{EmbeddingConfig} UUID, a specific
\texttt{AnnotationSet} (versioned by GOA release number), and a specific
\texttt{OntologySnapshot} (versioned by OBO release date). The temporal holdout
training pipeline is itself parameterised by version numbers: re-running the
same \texttt{train\_reranker\_auto} payload against the same database state
produces identical models and evaluation results. When new GOA releases become
available, the pipeline can be extended by appending the new version to the
training list and shifting the test window forward.

\subsection{The Value of Feature Engineering in the Age of Embeddings}

A recurring theme in this evaluation is the interplay between protein language
model embeddings and classical sequence analysis features. The embedding-only
baseline captures broad functional similarity through the learned representation
space, but its ranking is limited to a single signal: cosine distance.

The LightGBM re-ranker, when given access to the full 56-feature set, learns
to combine multiple complementary signals:
\begin{itemize}
  \item \textbf{Embedding distance} provides the initial KNN retrieval signal.
  \item \textbf{Alignment identity and scores} capture sequence-level conservation
    that the protein language model embedding may compress away.
  \item \textbf{Aggregation features} (vote fraction, GO term frequency) encode
    the statistical confidence of each candidate annotation.
  \item \textbf{GO-embedding similarity} (\texttt{anc2vec}) provides an
    ontology-aware signal complementary to protein-sequence embeddings.
  \item \textbf{Taxonomic distance} provides phylogenetic context.
  \item \textbf{Lineage features} (T-RES.1b; \texttt{lineage\_is\_ancestor\_of\_known},
    \texttt{lineage\_is\_descendant\_of\_known}, etc.) encode whether the candidate
    GO term is an ancestor or descendant of any term already known for the query
    protein.
\end{itemize}

The leakage-free full-feature importance analysis
(\cref{tab:feature-importance}) shows that vote-aggregation signals
(\texttt{neighbor\_vote\_fraction}, \texttt{go\_term\_frequency}) dominate
across all three categories, with GO-embedding similarity (\texttt{anc2vec})
contributing as a complementary ontology-aware term and raw protein-embedding
distance ranking outside the top-10. The much larger \texttt{anc2vec\_query}
ablation magnitude in the earlier 56-feature study
(\cref{tab:exp8-ablation}) is, as set out in the RQ2 answer above, a
replication-by-category dataset-construction artefact rather than evidence of
feature value; it is retained here as a methodological finding, patched in
PROTEA, and no conclusion in this section rests on it. With that caveat, the
leakage-clean evidence still supports the broader claim that feature
engineering remains valuable even when powerful embeddings are available: the
statistical structure of the retrieved neighbourhood and ontology-derived
similarity carry information not present in the raw protein-sequence embedding
alone. A clean per-feature re-derivation under full ground-truth
information-accretion scoring is provided by the score-ablation study reported
as future work.

\subsection{Limitations}
\label{subsec:limitations}

\begin{itemize}
  \item The temporal holdout strategy reduces contamination but does not
    eliminate it entirely: some test proteins share high sequence identity with
    reference proteins, making their annotations trivially predictable.
  \item ESMC~300M was chosen for computational tractability; larger variants
    would likely improve embedding quality at higher GPU cost.
  \item The evaluation considers only annotation transfer from Swiss-Prot
    experimental annotations; \gls{iea} annotations are excluded because they are
    themselves computationally derived.
  \item The IA file used (\texttt{IA\_cafa6.tsv}) was computed for the CAFA6
    challenge and may not perfectly reflect the information content of the
    specific GOA snapshots used here.
  \item The comparison with external tools is limited to eggNOG-mapper;
    additional tools (BLAST, InterProScan) would provide a more complete
    picture of the competitive landscape.
  \item The LK-BPO cell consistently favours the heuristic over the re-ranker,
    suggesting that the learned model may underperform for specific
    category-aspect combinations with limited positive examples.
  \item The Exp.~8 hyperparameter grid reveals a 0.0283 spread on NK-BPO
    (0.4386 to 0.4669), exceeding the informal robustness threshold of 0.02,
    indicating that optimal hyperparameters are cell-dependent and a joint
    per-cell tuning pass is warranted.
  \item The multi-PLM $\times$ $K$ generalisation grid (Exp.~9,
    \cref{tab:exp9-multiplm-grid}) is to be reported from the leakage-clean
    SELECT-window benchmark and is packaged via the F-DATA-PACK loop. The
    champion result (0.7291 $\pm$ 0.0028) is established on ProstT5 and does
    not depend on the grid; the grid provides breadth context for the
    generalisability claim.
\end{itemize}

\chapter{Conclusion}
\label{chap:conclusion}

% ────────────────────────────────────────────────────────────
\section{Doctoral Contribution Restated}
\label{sec:contribution-restated}

This thesis set out to answer a question that is both practical and
methodological: can a research engineer build, from scratch, a
reproducible distributed platform for protein functional annotation that
is honest about the difficulty of the problem and principled in how it
measures progress?

The answer embedded in the four years of work documented here is yes,
but with a caveat that the evaluation chapter makes precise. Honest
measurement requires rigorous temporal isolation. The anc2vec feature
leakage incident, documented in ADR~D30 and Appendix~B, demonstrated
that a reranker trained without strict temporal holdout can appear to
raise NK-BPO $F_{\max}$ from 0.431 to 0.599 over the embedding-only
baseline, a gain that vanishes almost entirely once the training window
is properly quarantined from the test window. The leakage-free reranker
(v3, full 52-feature schema) instead achieves a genuine and reproducible
improvement from a baseline of 0.412 to 0.431~$F_{\max}$ in NK-BPO,
with the largest verified gains in MFO (+0.049 in LK-MFO) and CCO
(+0.024 in NK-CCO). Smaller but consistent. More important.

The doctoral contribution is threefold.

\begin{description}
  \item[PROTEA as a reproducible, operationally auditable platform.]
    An open-source distributed system implementing the full annotation
    pipeline from FASTA upload to structured \gls{go} predictions, backed
    by a versioned relational data model that links every prediction to the
    exact ontology release, annotation set, embedding configuration,
    dataset, reranker model and code commit that produced it. The platform
    is horizontally scalable through queue-based worker pools, supports
    eight \gls{plm} backends through a contracts-first plugin architecture,
    and has been deployed across development, cloud, HPC (Apptainer) and
    airgapped environments without code branching.

  \item[A leakage-free reranker pipeline that quantifies the true lift
    of feature engineering over PLM embeddings.]
    A LightGBM selective learning-to-rank reranker trained per category
    (NK, LK, PK) over a 52-feature schema combining embedding distance,
    Needleman--Wunsch and Smith--Waterman alignment statistics,
    NCBI-taxonomy distance metrics, GO-graph ancestry embeddings, and
    per-PLM PCA projections. Training follows a temporal holdout protocol
    spanning twelve consecutive GOA splits, making it reproducible and
    resistant to the leakage pathologies that inflate reported performance
    in the literature. Feature importance analysis confirms that alignment
    and taxonomy features are genuinely complementary to embedding
    distance, and not merely correlated noise.

  \item[The LAFA~v2 method submission as an independently verifiable
    performance reference.]
    Three pre-built containers were submitted to the LAFA public
    benchmark, demonstrating that PROTEA's canonical pipeline can be
    deployed and evaluated in a blinded external setting without ad-hoc
    adaptation. The experimental work also contributed to a team placement
    of second in CAFA~6, validating the embedding-based annotation
    transfer approach in an international blinded benchmark.
\end{description}

% ────────────────────────────────────────────────────────────
\section{Summary of Empirical Results}
\label{sec:empirical-summary}

\Cref{chap:evaluation} reports seven progressive experiments on the
GOA~220$\to$229 evaluation set (20,281 proteins; NK: 2,831 / LK: 3,410 /
PK: 15,313) using IA-weighted $F_{\max}$ as the primary metric.

The embedding-only baseline ($k{=}5$, aspect-separated KNN over 527,000
Swiss-Prot reference sequences embedded with ESMC~300M) reaches $F_{\max}$
values of 0.412 (NK-BPO), 0.590 (NK-MFO) and 0.668 (NK-CCO) without any
sequence alignment. This alone exceeds eggNOG-mapper v2.1.13 in 8 of 9
category--aspect cells, with gains as large as +0.306~$F_{\max}$ (NK-CCO),
establishing that \gls{plm} embeddings provide a stronger retrieval signal
than DIAMOND-based orthology transfer under this evaluation protocol.

The reranker development demonstrated that feature engineering remains
valuable even when powerful embeddings are available, but only when the
training data is correctly constructed. The v2 reranker (where alignment
and taxonomy features were inadvertently NULL during training) failed to
surpass the simple heuristic that combined those features linearly. The
v3 reranker, trained with the full 52-feature schema, achieves the best
$F_{\max}$ in 7 of 9 cells. The largest absolute gains over the
embedding-only baseline are +0.049 (LK-MFO) and +0.030 (NK-MFO), with
consistent improvements across BPO and CCO as well (\cref{tab:improvement}).

Feature importance analysis of the v3 models (\cref{tab:feature-importance})
reveals three structural findings: (i)~aggregation features
(\texttt{ref\_annotation\_density}, \texttt{go\_term\_frequency}) dominate
NK and LK; (ii)~alignment features (\texttt{identity\_nw},
\texttt{similarity\_nw}) rank 4th--6th across all three models; (iii)~raw
embedding distance falls outside the top-10 features in NK/LK, indicating
that the reranker uses embedding similarity primarily as a candidate-generation
filter rather than as the final ranking signal. This finding was not
anticipated at the start of the project and has design implications for
future feature schemas.

The leakage incident documented in ADR~D30 deserves emphasis as an empirical
finding in its own right. The pre-leakfix reranker appeared to raise
NK-BPO from 0.431 to 0.599, a gain of +0.168. After removing the anc2vec
feature leakage, the verified gain is +0.019 in the same cell. The
difference between these two numbers, 0.149~$F_{\max}$ in a single cell,
illustrates concretely why temporal holdout discipline matters and why
the PROTEA data model, which makes every experiment fully reproducible,
is not a software engineering convenience but a scientific requirement.

% ────────────────────────────────────────────────────────────
\section{Answers to the Research Questions}
\label{sec:rq-answers}

\textbf{RQ1}: \emph{Can embedding-based KNN annotation transfer match or
exceed alignment-based methods for GO term prediction?}

Yes. The embedding-only baseline surpasses eggNOG-mapper in 8 of 9
cells on the GOA~220$\to$229 evaluation, and the reranker v3 surpasses
it in all 9. The margins are largest in NK and LK, where eggNOG-mapper
suffers from a 14.5\% coverage gap and uniform confidence scores. Even
with caveats around evaluation protocol differences
(\cref{subsec:exp7-eggnogmapper}), the advantage of the embedding-based
approach at low sequence identity is clear.

\textbf{RQ2}: \emph{Do pairwise alignment statistics and taxonomic distance
provide complementary signals?}

Yes. The progression from v2 to v3 provides controlled evidence: holding
the model architecture, hyperparameters and training splits fixed, making
alignment and taxonomy features available during training improved
$F_{\max}$ in 8 of 9 cells, with the largest gain in LK-MFO (+0.032).
The v2 model, without those features, failed to surpass the simple heuristic
that combined them linearly; the v3 model, with the same features, exceeded
it. The 52-feature schema is the mechanism by which this complementarity
is captured.

\textbf{RQ3}: \emph{How can a reproducible, scalable annotation pipeline
be designed?}

Through a versioned relational data model in which every prediction is
permanently linked to an immutable \texttt{EmbeddingConfig} UUID, a
versioned \texttt{AnnotationSet} (GOA release number), and a dated
\texttt{OntologySnapshot} (OBO release date). The temporal holdout
training pipeline is fully parameterised: re-running the same payload
against the same database state reproduces identical models and results.
The leakage incident demonstrated that this reproducibility is load-bearing,
not decorative: it was the ability to re-run all experiments cleanly
after the leakage fix that made the corrected numbers trustworthy.

\textbf{RQ4}: \emph{How can the journey of a research platform, with its
dead ends and pivots, be captured operationally?}

Through the \texttt{Job.findings} and \texttt{ExperimentRun} layers and
their associated free-form comments, which provide an auditable trace from
which the evaluation chapter was distilled. ADR~D30 (operational insights
appendix) formalises this: consequential incidents such as the schema SHA
drift, the anc2vec leakage, and the cafaeval coverage bug are documented
in Appendix~B rather than buried in commit history, making the platform's
empirical narrative verifiable independently of the code.

% ────────────────────────────────────────────────────────────
\section{Future Work}
\label{sec:future-work}

The following threads are grounded in work already partially shipped or
formally planned within the PROTEA development roadmap. They are organised
by domain rather than by priority, since priority will shift with
experimental feedback.

\subsection{F-MIL: Interpretable Predictions via Residue-Level Attention}
\label{subsec:fw-mil}

The current prediction pipeline produces GO term scores at the
protein level: a confidence value per (query protein, GO term) pair.
This is sufficient for ranking but provides no interpretability at the
residue level, a gap that limits utility for structural biologists and
drug target analysts who want to know \emph{which} regions of the sequence
drive the functional assignment.

The F-MIL thread addresses this by framing annotation as a multiple
instance learning (MIL) problem over the residue embedding sequence
produced by the \gls{plm} encoder. Each residue position constitutes
a bag instance; a permutation-invariant attention pooling layer produces
a protein-level representation that retains per-residue attention weights
as a natural interpretability signal. The protea-contracts package has
already shipped the \texttt{MILBackendContract} abstract base class and
the ESM residue-output backend (\texttt{MIL.1a}) on develop. The next
steps are: (i)~integrate residue-level predictions into the annotation
pipeline as an optional backend mode; (ii)~expose per-residue attention
weights through the REST API and the web interface; (iii)~validate that
the attention pattern identifies known active-site residues for a set
of enzymes with experimental structural data.

This thread does not require a new PLM family. It reuses the existing
ESMC~300M backend's residue-output endpoint, which is already accessible
through the contracts interface.

\subsection{F-IP: InterProScan Ensemble}
\label{subsec:fw-ip}

The evaluation in \cref{chap:evaluation} demonstrates that embedding-based
KNN transfer outperforms DIAMOND orthology (eggNOG-mapper) on the
GOA~220$\to$229 protocol. A natural next step is to incorporate
InterProScan~\cite{blum2021interpro} as a complementary annotation source
rather than a competitor. InterProScan integrates over a dozen family and
domain databases (Pfam, PRINTS, ProSite, TIGRFAMs and others) and infers
GO terms via curated mappings. Its signal is structurally different from
embedding-based KNN: it is rule-based, interpretable, and highly precise
for well-characterised protein families, but sparse for novel sequences.

The F-IP plan (IP.1--IP.4) covers four steps: (IP.1)~deploy InterProScan
as a PROTEA annotation source plugin via the existing
\texttt{protea-sources} entry point mechanism; (IP.2)~generate predictions
for the GOA~220$\to$229 test set and compute per-cell $F_{\max}$;
(IP.3)~train a reranker that combines InterProScan GO evidence with the
embedding-based features already in the 52-feature schema; (IP.4)~evaluate
the ensemble on a held-out split and report the delta relative to both the
embedding-only baseline and the eggNOG-mapper comparison.

The expected benefit is largest for PK proteins, where the protein already
has annotations in some namespaces: InterProScan family membership is a
strong prior for what additional terms are plausible.

\subsection{F-OPS: Production Hardening}
\label{subsec:fw-ops}

The current deployment covers the full pipeline but remains a
single-node system managed by \texttt{manage.sh}. Three hardening
threads are in progress.

The T-OPS bundle (\texttt{T5.1--T5.5}) covers structured logging
(JSON lines), metrics export (Prometheus), alerting (alert manager
rules for queue depth and worker error rate), distributed tracing
(OpenTelemetry spans from HTTP request to queue operation to database
write), and a health dashboard. The observability gap is the most
operationally significant limitation: without structured logs, diagnosing
worker failures or queue backlogs requires shell access to the deployment
host.

Authentication (\texttt{T5.6}) adds a token-based auth layer to the
REST API. The current API is unauthenticated and suitable only for
single-tenant or firewalled deployments. Token auth is a prerequisite
for any multi-tenant or public-facing instance, and is also required
before PROTEA can be offered as a web service to collaborators at
external institutions.

The airgapped Apptainer bundle (ADR~D27) is deferred to post-defense.
The priority is to ship T-OPS and T5.6 so that the platform is
operationally defensible in a production context, not merely a
research prototype that requires developer attention to keep running.

\subsection{LAFA v2 Follow-Ups}
\label{subsec:fw-lafa}

Three containers were submitted to the LAFA public benchmark (F-LAFA.1
through F-LAFA.3 as described in \cref{chap:introduction}). The
leakage-free reranker was not yet available at the time of the initial
submissions: those containers used the pre-leakfix feature schema. The
F-LAFA.3 thread covers rebuilding the LAFA submission container with the
corrected v3 reranker and the full 52-feature schema, so that the LAFA
leaderboard result reflects the same pipeline as the evaluation chapter.

F-LAFA.4 will run a full ablation study on the LAFA evaluation set,
paralleling the GOA~220$\to$229 experiments in \cref{chap:evaluation}
but using the LAFA-held-out proteins as the test set. This provides an
independent verification of the reranker's generalisation beyond the
single GOA split used in this thesis. It also tests whether the feature
importance hierarchy observed in \cref{tab:feature-importance} holds on
a different protein set.

% ────────────────────────────────────────────────────────────
\section{Engineering Wrap-Up}
\label{sec:engineering-wrapup}

Three architectural decision records close the engineering narrative.

\textbf{ADR~D29 (Release pipeline).} The per-repo SemVer release pipeline
driven by Conventional Commits and \texttt{release-please} automates
version tagging, changelog generation, and container image publication
to GitHub Container Registry. The release pipeline is a prerequisite for
the T-OPS observability work (F-OPS) because structured versioning is
required before cross-repo integration tests can be reliably anchored
to a stable tag.

\textbf{ADR~D30 (Operational insights appendix).} This ADR formalises
the decision to document consequential incidents (schema SHA drift,
anc2vec feature leakage, cafaeval coverage bug) in Appendix~B rather
than in commit messages or design documents. The rationale is auditability:
incidents that affected the empirical results of the thesis should be
in a location that a reader can find without access to the git log.
The leakage incident described in \cref{sec:empirical-summary} is the
primary example. ADR~D30 makes the operational narrative an explicit part
of the thesis artefact rather than an implicit side effect of version
control.

\textbf{ADR~D31 (Method Object reframe).} The refactor strategy for large
batch operation classes, extracting cohesive sub-clusters as
\texttt{\_<Cluster>Runner} helper classes rather than wrapping
already-budgeted \texttt{execute()} methods, is codified in D31. This
decision directly shapes how the F-MIL and F-IP operation classes will
be structured when they land in \texttt{protea-runners}: new operations
follow the same sub-cluster pattern, keeping the executor graph navigable
as the operation catalogue grows from its current 17-router topology.

% ────────────────────────────────────────────────────────────
\section{Closing Remarks}
\label{sec:closing}

PROTEA began as a question about reproducibility in computational biology
and ended as a demonstration that the question is load-bearing. The
leakage incident is the most direct evidence: a platform without a
versioned data model and reproducible training pipelines cannot distinguish
a genuine +0.049~$F_{\max}$ gain from an artefact of data contamination.
A platform with those properties can.

The functional annotation gap is not close to being solved. Fewer than
0.02\,\% of deposited proteins carry experimentally validated annotations,
and that fraction shrinks as sequencing throughput grows. PROTEA's
contribution is not to close the gap but to provide a foundation from
which incremental, verifiable progress can be made, and to demonstrate
on a concrete evaluation set that embedding-based annotation transfer
with selective feature-engineered reranking is a meaningful step in
the right direction.

The platform is released as open-source software at
\url{https://github.com/frapercan/protea}. Francisco Miguel Pérez Canales
is its author and sole maintainer.


% ── Back matter ─────────────────────────────────────────────
\backmatter
\printbibliography[heading=bibintoc]

\appendix
\chapter{REST API Reference}
\label{app:api}

This appendix documents the public endpoints of the PROTEA REST API.
All endpoints are prefixed with the base URL of the FastAPI server
(default: \texttt{http://127.0.0.1:8000}).

\sloppy

% ────────────────────────────────────────────────────────────
\section{Jobs}
\label{app:api-jobs}

\begin{description}
  \item[\texttt{POST /jobs}] Submit a new job. The request body must contain
    \texttt{operation} (string) and \texttt{payload} (JSON object). Returns the
    created \texttt{Job} row.
  \item[\texttt{GET /jobs}] List jobs, optionally filtered by \texttt{status}
    and \texttt{operation} query parameters.
  \item[\texttt{GET /jobs/\{id\}}] Retrieve a single job by UUID.
  \item[\texttt{GET /jobs/\{id\}/events}] Stream the structured event log for
    a job as newline-delimited JSON.
  \item[\texttt{POST /jobs/\{id\}/cancel}] Request cancellation of a queued or
    running job.
\end{description}

% ────────────────────────────────────────────────────────────
\section{Proteins and Sequences}
\label{app:api-proteins}

\begin{description}
  \item[\texttt{POST /proteins/insert}] Enqueue an \texttt{insert\_proteins} job
    for a list of UniProt accessions or a taxonomy query.
  \item[\texttt{POST /proteins/metadata}] Enqueue a \texttt{fetch\_uniprot\_metadata}
    job to refresh functional annotation data.
  \item[\texttt{GET /proteins}] List proteins stored in the database with optional
    accession filter.
\end{description}

% ────────────────────────────────────────────────────────────
\section{Query Sets}
\label{app:api-querysets}

\begin{description}
  \item[\texttt{POST /query-sets}] Upload a FASTA file. Creates a \texttt{QuerySet}
    and its \texttt{QuerySetEntry} rows; upserts novel sequences and proteins.
    Returns the created \texttt{QuerySet} UUID.
  \item[\texttt{GET /query-sets}] List all query sets.
  \item[\texttt{GET /query-sets/\{id\}}] Retrieve a query set with its entries.
  \item[\texttt{DELETE /query-sets/\{id\}}] Delete a query set and its entries.
\end{description}

% ────────────────────────────────────────────────────────────
\section{Annotations}
\label{app:api-annotations}

\begin{description}
  \item[\texttt{POST /annotations/ontology-snapshot}] Enqueue a
    \texttt{load\_ontology\_snapshot} job for a given OBO URL.
  \item[\texttt{POST /annotations/annotation-set/goa}] Enqueue a
    \texttt{load\_goa\_annotations} job for a local GAF file path.
  \item[\texttt{POST /annotations/annotation-set/quickgo}] Enqueue a
    \texttt{load\_quickgo\_annotations} job with optional evidence code filter.
  \item[\texttt{GET /annotations/ontology-snapshots}] List loaded ontology snapshots.
  \item[\texttt{GET /annotations/annotation-sets}] List loaded annotation sets.
\end{description}

% ────────────────────────────────────────────────────────────
\section{Embeddings and Predictions}
\label{app:api-embeddings}

\begin{description}
  \item[\texttt{POST /embeddings/configs}] Create a new \texttt{EmbeddingConfig}.
  \item[\texttt{GET /embeddings/configs}] List embedding configurations.
  \item[\texttt{POST /embeddings/compute}] Enqueue a \texttt{compute\_embeddings}
    coordinator job.
  \item[\texttt{POST /embeddings/predict}] Enqueue a \texttt{predict\_go\_terms}
    coordinator job.
  \item[\texttt{GET /embeddings/prediction-sets}] List prediction sets.
  \item[\texttt{GET /embeddings/prediction-sets/\{id\}}] Retrieve a prediction
    set with summary statistics.
  \item[\texttt{GET /embeddings/prediction-sets/\{id\}/predictions.tsv}] Stream
    predictions as a 27-column TSV file. Optional query parameters:
    \texttt{accession}, \texttt{aspect} (\texttt{F}/\texttt{P}/\texttt{C}),
    \texttt{max\_distance}.
\end{description}

\chapter{External Tool Comparison: Reproducibility Protocol}
\label{app:external-comparison}

\paragraph{Scope.}
This appendix collects the full reproduction protocol for the head-to-head
comparison between PROTEA and eggNOG-mapper~v2.1.13 reported in
\cref{subsec:exp7-eggnogmapper}, targeting reviewers who wish to verify,
extend, or replicate the evaluation independently. It covers test-set
extraction, tool installation, run parameters, evaluation alignment with
the CAFA protocol, and complete per-cell $F_{\max}$ tables that were
abbreviated in the main text. It does not overlap with
\cref{app:api} (REST API reference) or \cref{app:adr-log} (architecture
decision records): those appendices address system design, whereas this one
addresses empirical reproducibility of a specific comparative experiment.

This appendix documents the exact procedure used to evaluate eggNOG-mapper
on the same test set and evaluation protocol as PROTEA (\cref{subsec:exp7-eggnogmapper}).
It is intended to allow full reproduction of the comparison results and to
serve as a template for evaluating additional external tools.

% ────────────────────────────────────────────────────────────
\section{Test Set Extraction}
\label{app:ext-test-set}

The test set consists of the 20{,}281 proteins that gained at least one new
experimental GO annotation between GOA version~220 and version~229 (the same
temporal holdout used for all PROTEA experiments in \cref{chap:evaluation}).
The FASTA file was extracted from PROTEA's database using the built-in API
endpoint:

\begin{verbatim}
curl -o test_proteins.fasta \
  "http://localhost:8000/annotations/evaluation-sets/\
42b34e79-6fe9-4fa0-b718-02f43a1e3192/delta-proteins.fasta"
\end{verbatim}

Each FASTA header contains the UniProt accession, entry name, organism,
taxonomy ID, and CAFA category (NK/LK/PK):

\begin{verbatim}
>A0A024RBG1 NUD4B_HUMAN OS=Homo sapiens OX=9606 (LK)
MMKFKPNQTRTYDREGFKKRAACLCFRSEQEDEVLLVSSS...
\end{verbatim}

The resulting file contains 20{,}281 sequences spanning all three CAFA categories
(NK: 2{,}831, LK: 3{,}410, PK: 15{,}313).

% ────────────────────────────────────────────────────────────
\section{eggNOG-mapper Setup}
\label{app:ext-emapper-setup}

eggNOG-mapper v2.1.13 was run via its official Biocontainers Docker image
to ensure a reproducible environment. The following steps were performed:

\paragraph{1.\ Pull the Docker image.}
\begin{verbatim}
docker pull quay.io/biocontainers/\
  eggnog-mapper:2.1.13--pyhdfd78af_2
\end{verbatim}

\paragraph{2.\ Download the eggNOG databases (v5.0.2).}
At the time of this experiment (March 2026), the official download script
pointed to a deprecated URL. The databases were downloaded manually from
\texttt{eggnog5.embl.de}:

\begin{verbatim}
# eggnog.db (6.3 GB compressed, ~39 GB uncompressed)
wget -O eggnog.db.gz \
  http://eggnog5.embl.de/download/emapperdb-5.0.2/\
eggnog.db.gz
gunzip eggnog.db.gz

# Diamond protein database (4.8 GB compressed, ~8.7 GB)
wget -O eggnog_proteins.dmnd.gz \
  http://eggnog5.embl.de/download/emapperdb-5.0.2/\
eggnog_proteins.dmnd.gz
gunzip eggnog_proteins.dmnd.gz

# Taxonomy database (~71 MB compressed)
wget -O eggnog.taxa.tar.gz \
  http://eggnog5.embl.de/download/emapperdb-5.0.2/\
eggnog.taxa.tar.gz
tar xzf eggnog.taxa.tar.gz
\end{verbatim}

Total disk usage: approximately 48~GB.

\paragraph{3.\ Verify installation.}
\begin{verbatim}
docker run --rm \
  -v /path/to/eggnog-data:/data \
  quay.io/biocontainers/\
    eggnog-mapper:2.1.13--pyhdfd78af_2 \
  emapper.py --version --data_dir /data

# Output:
# emapper-2.1.13 / eggNOG DB version: 5.0.2
# Diamond version: 2.0.15 / MMseqs2: 16.747c6
\end{verbatim}

% ────────────────────────────────────────────────────────────
\section{Running eggNOG-mapper}
\label{app:ext-emapper-run}

\begin{verbatim}
docker run --rm \
  -v /path/to/eggnog-data:/data \
  -v /path/to/input:/input \
  -v /path/to/output:/output \
  quay.io/biocontainers/\
    eggnog-mapper:2.1.13--pyhdfd78af_2 \
  emapper.py \
    -i /input/test_proteins.fasta \
    --data_dir /data \
    -o test_proteins \
    --output_dir /output \
    --cpu 8 \
    -m diamond \
    --go_evidence experimental \
    --tax_scope auto \
    --target_orthologs all
\end{verbatim}

\paragraph{Key parameters.}
\begin{description}
  \item[\texttt{-m diamond}] Use DIAMOND for sequence search (default mode).
  \item[\texttt{--go\_evidence experimental}] Transfer only GO terms supported
    by experimental evidence codes, matching the evidence code filter used in
    PROTEA's evaluation protocol.
  \item[\texttt{--tax\_scope auto}] Automatically determine the taxonomic scope
    for ortholog assignment.
  \item[\texttt{--target\_orthologs all}] Consider all orthologs (one-to-one,
    one-to-many, many-to-many).
  \item[\texttt{--cpu 8}] Use 8 CPU threads for the DIAMOND search phase.
\end{description}

\paragraph{Runtime.} The full run completed in approximately 21 minutes
on the same workstation used for all PROTEA experiments (64~GB RAM, no GPU
required for eggNOG-mapper).

\paragraph{Output.} eggNOG-mapper produced annotations for 19{,}958 of the
20{,}281 input proteins (98.4\% hit rate). Of those, 17{,}334 (85.5\% of the
test set) received at least one GO term prediction. The remaining 2{,}947
proteins either had no DIAMOND hit or no GO terms were transferred from the
matched orthologous group.

% ────────────────────────────────────────────────────────────
\section{CAFA Evaluation Protocol}
\label{app:ext-cafa-eval}

To ensure a fair comparison, eggNOG-mapper predictions were evaluated using
exactly the same protocol as all PROTEA experiments:

\begin{enumerate}
  \item \textbf{Ground truth computation.} The NK/LK/PK classification and
    ground-truth annotations were computed from the same EvaluationSet
    (GOA~220$\to$229) using PROTEA's \texttt{compute\_evaluation\_data()}
    function, which implements NOT-qualifier propagation through the GO DAG
    following the CAFA5 protocol.
  \item \textbf{Prediction format.} eggNOG-mapper does not produce per-term
    confidence scores; each predicted GO term is assigned a uniform score of
    1.0 (binary: predicted or not predicted).
  \item \textbf{Evaluation.} The \texttt{cafaeval} tool was run with the same
    parameters as all other experiments: GO propagation mode \texttt{max},
    normalisation mode \texttt{cafa}, and IA weighting from
    \texttt{IA\_cafa6.tsv}.
\end{enumerate}

The evaluation script \texttt{scripts/evaluate\_external\_tool.py} automates
this entire pipeline. It accepts the external tool's output file, connects
to PROTEA's database to retrieve the ground truth, and runs \texttt{cafaeval}
for all three CAFA settings (NK, LK, PK):

\begin{verbatim}
poetry run python scripts/evaluate_external_tool.py \
  --evaluation-set-id \
    42b34e79-6fe9-4fa0-b718-02f43a1e3192 \
  --tool emapper \
  --input output/test_proteins.emapper.annotations
\end{verbatim}

% ────────────────────────────────────────────────────────────
\section{Raw Results}
\label{app:ext-raw-results}

\Cref{tab:emapper-raw} presents the complete per-cell results for
eggNOG-mapper alongside the PROTEA methods.

\begin{table}[htbp]
  \centering
  \caption[Complete Fmax comparison]{Complete $F_{\max}$ comparison: eggNOG-mapper vs.\ PROTEA methods (IA-weighted).}
  \label{tab:emapper-raw}
  \resizebox{\textwidth}{!}{%
  \begin{tabular}{l ccc ccc ccc}
    \toprule
    & \multicolumn{3}{c}{\textbf{NK}} & \multicolumn{3}{c}{\textbf{LK}} & \multicolumn{3}{c}{\textbf{PK}} \\
    \cmidrule(lr){2-4} \cmidrule(lr){5-7} \cmidrule(lr){8-10}
    \textbf{Method} & BPO & MFO & CCO & BPO & MFO & CCO & BPO & MFO & CCO \\
    \midrule
    eggNOG-mapper 2.1.13 & 0.247 & 0.359 & 0.386 & 0.382 & 0.334 & 0.450 & 0.190 & 0.199 & 0.325 \\
    PROTEA baseline      & 0.412 & 0.590 & 0.668 & 0.467 & 0.558 & 0.675 & 0.187 & 0.278 & 0.325 \\
    PROTEA full-feature  & \textbf{0.431} & \textbf{0.620} & \textbf{0.692} & \textbf{0.478} & \textbf{0.607} & \textbf{0.697} & \textbf{0.201} & \textbf{0.297} & \textbf{0.339} \\
    \bottomrule
  \end{tabular}
  }
\end{table}

\Cref{tab:emapper-delta} shows the absolute $F_{\max}$ difference between
PROTEA's full-feature re-ranker and eggNOG-mapper.

\begin{table}[htbp]
  \centering
  \caption[Absolute Fmax improvement vs.\ eggNOG-mapper]{Absolute $F_{\max}$ improvement: PROTEA full-feature vs.\ eggNOG-mapper.}
  \label{tab:emapper-delta}
  \begin{tabular}{lccc}
    \toprule
    \textbf{Category} & \textbf{BPO} & \textbf{MFO} & \textbf{CCO} \\
    \midrule
    NK & +0.184 & +0.261 & +0.306 \\
    LK & +0.096 & +0.273 & +0.247 \\
    PK & +0.011 & +0.098 & +0.014 \\
    \bottomrule
  \end{tabular}
\end{table}

\paragraph{Coverage analysis.} eggNOG-mapper found DIAMOND hits for 19{,}958
proteins but only transferred GO terms to 17{,}334 (85.5\% of the test set).
PROTEA, by contrast, produces predictions for all 20{,}281 proteins (100\%
coverage) because every protein has an embedding and therefore has $k$
nearest neighbours. This coverage gap accounts for part of the $F_{\max}$
difference, particularly in NK where 14.5\% of proteins receive no
eggNOG-mapper prediction.

% ────────────────────────────────────────────────────────────
\section{Methodological Notes}
\label{app:ext-notes}

\paragraph{Score assignment.}
eggNOG-mapper produces a binary set of GO terms per protein without
per-term confidence scores. For the CAFA evaluation, all predicted terms are
assigned a confidence of 1.0. This means the precision-recall curve is a
single point (one threshold), and $F_{\max}$ equals the $F_1$ score at that
point. Methods that produce graded confidence scores (such as PROTEA) can
optimise the precision-recall trade-off by varying the threshold, which may
contribute to higher $F_{\max}$ values.

\paragraph{GO propagation.}
eggNOG-mapper transfers GO terms from orthologous groups but does not
explicitly propagate them up the GO DAG. The \texttt{cafaeval} evaluator
applies \texttt{prop=max} propagation during evaluation, so ancestor terms
are included in the comparison for all methods.

\paragraph{Evidence code filter.}
The \texttt{--go\_evidence experimental} flag restricts eggNOG-mapper to
transferring GO terms supported by experimental evidence codes in the eggNOG
database. This is conceptually aligned with PROTEA's use of experimental-only
annotations from GOA, though the exact set of supporting evidence and the
database versions differ.

\chapter{Architecture Decision Records: Decision Log}
\label{app:adr-log}

This appendix summarises the most consequential Architectural Decision Records
(ADRs) for the PROTEA platform. Each record was written autonomously as
implementation evidence accumulated and confirmed by the author on 2026-05-06.
The full text of every record lives in \texttt{docs/source/adr/} inside the
PROTEA repository; this log captures the rationale and the accepted resolution
for the decisions that most directly shaped the platform's design and
operational posture.

ADRs are identified by two complementary numbering schemes. The legacy
three-digit codes (ADR-001 through ADR-008) cover the earliest architectural
choices and are referenced in Chapter~4. The \textbf{D-series} codes (D1
onward) cover the post-F2 design decisions discussed in Chapters~4 and~5 and
throughout this log. Both series are cross-referenced where they overlap.

% ────────────────────────────────────────────────────────────
\section{D4: API Versioning Strategy}
\label{app:adr-d4}

\paragraph{Context.}
PROTEA exposes a REST API consumed by the front-end, by LAFA containers, and
by downstream pipelines. As the API surface stabilised toward the F4 release,
breaking changes required a versioning strategy that would not strand existing
consumers.

\paragraph{Decision accepted 2026-05-06.}
A universal \texttt{/v1/} path prefix is applied to all routers. Future
breaking changes branch under \texttt{/v2/} without removing \texttt{/v1/}
endpoints. \texttt{Accept}-header negotiation is deferred unless a concrete
need surfaces. Deprecated unprefixed mounts are kept for one release cycle to
avoid breaking live clients (the \texttt{protea.ngrok.app} deployment and
existing front-end fetchers) during the transition.

\paragraph{Consequences.}
Each API version ships its own OpenAPI document. Schemathesis runs against the
version under test. Front-end and external clients update their base URLs once,
at the \texttt{/v1/} transition, and are thereafter insulated from internal
restructuring under the same major version.

% ────────────────────────────────────────────────────────────
\section{D6: Authentication Strategy}
\label{app:adr-d6}

\paragraph{Context.}
Sensitive endpoints (job creation, dataset import, re-ranker model upload,
evaluation triggers) were unauthenticated in the pre-F5 platform. Public cloud
exposure, LAFA submission tooling, and external adopters required an
authentication layer before the F5 phase could be declared open.

\paragraph{Decision accepted 2026-05-06.}
Two complementary mechanisms are deployed:

\begin{itemize}
  \item \textbf{API key} for service-to-service calls: an \texttt{ApiKey} ORM
    table issues opaque tokens; callers present them as
    \texttt{Authorization: Bearer \textlangle{}token\textrangle{}}.
  \item \textbf{OIDC via reverse proxy} for human users: \textbf{Authentik}
    acts as the identity provider behind \texttt{oauth2-proxy}, chosen for
    its lighter operational footprint and simpler Docker-Compose setup
    relative to Keycloak.
\end{itemize}

Rate limiting is enforced via \texttt{slowapi} on a per-endpoint basis.
A CORS allowlist is maintained in the FastAPI middleware configuration.

\paragraph{Implementation status.}
Partial as of the thesis writing phase. Tasks T5.6a (API key ORM and
middleware) and T5.6b (oauth2-proxy integration) carry the outstanding work.

% ────────────────────────────────────────────────────────────
\section{D7: Observability Stack}
\label{app:adr-d7}

\paragraph{Context.}
PROTEA's pre-F-OPS observability relied on per-process log files and ad-hoc
log statements. Multi-target deployment (cloud, HPC, airgap) and external
adopters require distributed tracing, metrics with service-level objectives,
and structured log aggregation.

\paragraph{Decision accepted 2026-05-06.}
A single canonical observability stack:

\begin{itemize}
  \item \textbf{Tracing}: OpenTelemetry (OTLP exporter) instrumenting FastAPI,
    SQLAlchemy, and \texttt{pika}. The \texttt{traceparent} header is
    propagated HTTP to queue to worker so that a single prediction is
    visible end-to-end as one OTel trace.
  \item \textbf{Metrics}: Prometheus client; a \texttt{/metrics} scrape
    endpoint is exposed by the FastAPI process.
  \item \textbf{Dashboards}: Grafana, with dashboard JSON committed under
    \texttt{deploy/grafana/} so dashboards are version-controlled.
  \item \textbf{Logs}: structured JSON via \texttt{python-json-logger},
    shipped to Loki via the \texttt{loki-docker-driver} from container
    stdout (no separate Promtail sidecar in the cloud target). Loki was
    preferred over the ELK stack because it indexes labels rather than
    full text, has a smaller memory footprint, and integrates with the
    same Grafana instance already surfacing Prometheus dashboards.
\end{itemize}

Three service-level objectives are documented in \texttt{docs/SLOs.md};
alert rules and per-alert runbooks are committed to the repository.

\paragraph{Implementation status.}
Tasks T5.1a (OTel instrumentation), T5.1b (Prometheus endpoint), T5.2
(Grafana dashboards), T5.3 (Loki driver), and T5.4 (alert rules) carry the
phased rollout; T5.1a and T5.1b are complete.

% ────────────────────────────────────────────────────────────
\section{D10: \texttt{schema\_sha} v2 Parallel Migration}
\label{app:adr-d10}

\paragraph{Context.}
\texttt{schema\_sha} is the load-bearing fingerprint that prevents inference
from running with a re-ranker booster trained against a different feature
schema. Two independent definitions of \texttt{compute\_schema\_sha} (one in
the lab, one in PROTEA) silently drifted until the parity bug surfaced during
the v9 study on 2026-05-01 and caused a non-reproducible run.

\paragraph{Decision accepted 2026-05-06.}
A parallel \texttt{schema\_sha\_v2} column is added to both \texttt{Dataset}
and \texttt{RerankerModel}. The column is backfilled from the canonical
\texttt{protea\_contracts.compute\_schema\_sha} implementation. The production
read path switches to \texttt{v2}; the \texttt{v1} column is retained until
the F3 gate for audit purposes and then dropped. A regression test exposes the
v1/v2 drift on historical rows rather than retroactively modifying them.

\paragraph{Implementation gate.}
Task T1.6 carries the Alembic migration and backfill. The migration is
classified \texttt{requires\_human}: it must be rehearsed in staging (or a
local-DB dry-run) and the backfill verified before any production rollout.
This constraint was explicitly confirmed by the author on 2026-05-06.

% ────────────────────────────────────────────────────────────
\section{D15: \texttt{protea-method} Distribution Channels}
\label{app:adr-d15}

\paragraph{Context.}
\texttt{protea-method} is the pure inference path (KNN, feature computation,
re-ranker application). Because it has no FastAPI or SQLAlchemy dependencies
it can ship independently of the full platform, addressing audiences that want
the annotation method without operating the database and messaging
infrastructure.

\paragraph{Decision accepted.}
Three distribution channels:

\begin{enumerate}
  \item \textbf{PyPI} public package (\texttt{pip install protea-method}).
    The default install is torch-free; the gated-attention MIL head is
    opt-in via the \texttt{[mil]} Poetry extra (see also D\,no\,extra below).
  \item \textbf{Docker Hub} minimal image (\texttt{protea-method-runtime},
    approximately 3 to 4\,GB with one canonical PLM and one default booster;
    CLI: \texttt{protea-predict input.fasta output.tsv}).
  \item \textbf{Companion engineering paper} covering the F-OPS pillars
    (containers, multi-target deployment, observability, secrets management).
\end{enumerate}

LAFA submission containers (D23) build on top of
\texttt{protea-method-runtime}, making D15 a prerequisite for the submission
tooling.

\paragraph{Distribution discipline.}
The \textbf{opt-in \texttt{[mil]} extra} decision (recorded separately in the
implementation log) keeps the default \texttt{pip install protea-method}
torch-free: the gated-attention MIL head is only installed when the caller
explicitly requests it. This avoids a several-gigabyte PyTorch dependency for
users who only need the LightGBM re-ranker inference path.

% ────────────────────────────────────────────────────────────
\section{D25: HPC Operation Mode}
\label{app:adr-d25}

\paragraph{Context.}
PROTEA must support HPC environments such as the BSC MareNostrum cluster.
HPC sites typically forbid privileged Docker, may restrict outbound network
access, and schedule workloads via SLURM. Two operation modes were
considered:

\begin{description}
  \item[Mode B] Stateless PROTEA workers running on HPC nodes connect to a
    PostgreSQL and RabbitMQ hosted in the cloud (LifeWatch or EOSC). No
    persistent state on the HPC node itself.
  \item[Mode C] Fully airgapped batch bundle: an Apptainer \texttt{.sif}
    image with a snapshot database pre-loaded, the default booster, and a
    single-node SLURM job requiring no outbound traffic.
\end{description}

\paragraph{Decision accepted 2026-05-06.}
\textbf{Mode B is the primary HPC deployment target.} It is closer to the
cloud architecture, reuses the same Postgres and RabbitMQ services, and
requires no special airgap provisioning. Mode C (the airgapped \texttt{.sif}
bundle, T-OPS.9) is \textbf{deferred per D25}: it becomes a post-defense
item (F9), contingent on concrete engagement with a restricted-access HPC
site. Mode B is sufficient for the thesis defense scope.

% ────────────────────────────────────────────────────────────
\section{D27: Image Registry}
\label{app:adr-d27}

\paragraph{Context.}
Seven OCI images must be hosted in a registry visible from cloud deployments,
from HPC tooling that pulls images before converting them to \texttt{.sif}
(Apptainer), and from external adopters consuming
\texttt{protea-method-runtime}.

\paragraph{Decision accepted 2026-05-06.}
\texttt{ghcr.io} (GitHub Container Registry) is the canonical registry.
GitHub Actions push images on SemVer tag using the repository's own
\texttt{GITHUB\_TOKEN}; no external registry credentials are required.
\texttt{protea-method-runtime} is published with public visibility; internal
images are org-scoped or private. Mirroring to Docker Hub is deferred until
external pull rates demand it. Task T-OPS.8 implements the \texttt{docker.yml}
workflow and is complete.

% ────────────────────────────────────────────────────────────
\section{D28: Secrets Management}
\label{app:adr-d28}

\paragraph{Context.}
PostgreSQL credentials, MinIO keys, OIDC client secrets, optional external API
tokens, and SSH keys cannot live in plaintext in any repository.
Multi-target deployment (cloud, HPC, airgap) requires a single mechanism that
works across all environments without per-environment custom solutions.

\paragraph{Decision accepted 2026-05-06.}
\texttt{sops} with \texttt{age} keys. Encrypted
\texttt{secrets.\{dev,prod\}.enc.yaml} files are committed to the repository;
CI decrypts with the age key stored as a \texttt{SOPS\_AGE\_KEY} GitHub
repository secret. Two reasons capture the choice: (a) age keys are
post-PGP ed25519, short, and revocation-chain-free; (b) sops is
file-format-agnostic so the same workflow handles YAML, JSON, and
\texttt{.env} files uniformly.

\paragraph{Implementation status.}
The scaffolding was committed alongside this ADR: \texttt{.sops.yaml},
\texttt{secrets/secrets.dev.example.yaml} (plaintext schema template,
tracked), and updated \texttt{.gitignore} rules that forbid the plaintext
\texttt{secrets.\{dev,prod\}.yaml} while keeping encrypted forms and
the example template tracked. The manual bootstrap step (generating an
\texttt{age} keypair and populating the recipient placeholder in
\texttt{.sops.yaml}) is classified \texttt{requires\_human} as task T-OPS.7.

% ────────────────────────────────────────────────────────────
\section{D29: Release Pipeline}
\label{app:adr-d29}

\paragraph{Context.}
Seven repositories need independent SemVer release cadences. Releasing one
repository must not break another; cross-repository integration testing is
required on every tag. \texttt{protea-contracts} is the most disruptive
package: version bumps propagate through all seven consumers.

\paragraph{Decision accepted 2026-05-06.}
Per-repository SemVer plus a cross-repository integration test on tag:

\begin{itemize}
  \item A SemVer tag (\texttt{vX.Y.Z}) in any repository dispatches a build,
    an image push (D27), and an integration test that pulls the new image
    together with the pinned versions of the other six repositories and runs
    a smoke pipeline.
  \item Failures block image promotion; the tag remains but the image is
    held as a release candidate until the integration test passes.
  \item A \texttt{protea-contracts} release triggers an automated re-pin PR
    in all consumers.
\end{itemize}

\texttt{release-please} is used for version bumps and CHANGELOG generation,
driven by Conventional Commits: \texttt{feat:} increments the minor version,
\texttt{fix:} increments the patch, and a \texttt{BREAKING CHANGE:} footer or
\texttt{feat!:} increments the major version. This differs from the originally
considered \texttt{python-semantic-release}, which push-commits the version
bump directly to \texttt{main} and would bypass branch protection.
\texttt{release-please} instead opens a release PR, preserving the hard project
rule that nothing lands on \texttt{main} without a passing PR review.

The commit-message convention (Conventional Commits) was already enforced from
the F2 phase onward, so the CHANGELOG generation requires no retroactive
reformatting.


\end{document}
